% Chap2: EMPHASIS ON STORYLINE, THE PLOT, THE SITUATION, ETC — Literature in English Upper Sixth — Cours signé OnBuch+
% Official MINESEC syllabus. Compiler avec : tectonic LIT02-chap2-emphasis-on-storyline-the-plot-the-situation-etc.tex
\newcommand{\DOCMATIERE}{Literature in English}
\newcommand{\DOCNIVEAU}{Upper Sixth}
\newcommand{\DOCMODULE}{Upper Sixth — Literature in English}
\newcommand{\DOCLECON}{2}
\newcommand{\DOCTITRE}{Chap2: EMPHASIS ON STORYLINE, THE PLOT, THE SITUATION, ETC}
% OnBuch+ — préambule « cours signés » ENRICHI (compilé avec tectonic / XeLaTeX)
% Système visuel « Pop » d'OnBuch+ : fond crème (jamais blanc), bordures encre
% épaisses, ombres « dures » décalées, titres Archivo Black en capitales.
% Version MATHÉMATIQUES : graphes (pgfplots/tikz), boîtes pédagogiques
% (définition, propriété, méthode, attention, le savais-tu, exercice résolu,
% exercice, TP, bilan), page de garde et signature OnBuch+ ancrée en bas.
%
% Macros attendues dans main.tex AVANT \input{preamble} :
%   \DOCMATIERE  \DOCNIVEAU  \DOCMODULE  \DOCLECON (numéro)  \DOCTITRE
\documentclass[11pt]{article}

\usepackage{fontspec}
\usepackage[english]{babel}
\usepackage[a4paper,margin=2cm,top=3.4cm,bottom=2.8cm,headheight=1.4cm,footskip=1.3cm]{geometry}
\usepackage{amsmath,amssymb,amsfonts}
\usepackage{mathtools} % \xmapsto, \coloneqq... souvent utilisés par les modèles
\usepackage{cancel} % simplifications barrées (bilans de forces)
% \abs{z} (module, valeur absolue) et \norm{u} (norme), utilisés par les modèles
\providecommand{\abs}{}\let\abs\relax\DeclarePairedDelimiter{\abs}{\lvert}{\rvert}
\providecommand{\norm}{}\let\norm\relax\DeclarePairedDelimiter{\norm}{\lVert}{\rVert}
\usepackage[table]{xcolor} % [table] : fournit aussi \rowcolors (alternance de lignes)
\usepackage{pagecolor}
\usepackage{tikz}
\usetikzlibrary{arrows.meta,positioning,calc,shapes.geometric,shapes.misc,shapes.arrows,decorations.pathmorphing,decorations.pathreplacing,decorations.markings,patterns,shadows,mindmap,trees,fit,backgrounds,matrix,intersections,angles,quotes}
\usepackage{pgfplots}
\usepackage[version=4]{mhchem} % équations nucléaires \ce{^{235}_{92}U}
\usepackage[european,siunitx]{circuitikz} % schémas électriques
\pgfplotsset{compat=1.18}
\usepgfplotslibrary{fillbetween,groupplots} % aires entre courbes (calcul intégral)
\usepackage{enumitem}
\usepackage{fancyhdr}
% [most] charge toutes les bibliothèques tcolorbox (listings, minted, raster,
% documentation...) et gonfle inutilement la mémoire TeX sur un document long
% (~300 boîtes) : on ne charge que ce qui est réellement utilisé.
\usepackage[skins,breakable]{tcolorbox}
\usepackage{graphicx}
\usepackage{colortbl}
\usepackage{booktabs}
\usepackage{tabularx}
\usepackage{diagbox} % case d'en-tête diagonale des tableaux à double entrée
% Colonnes extensibles centrée / gauche / droite (C, L, R), souvent utilisées par les modèles
\newcolumntype{C}{>{\centering\arraybackslash}X}
\newcolumntype{L}{>{\raggedright\arraybackslash}X}
\newcolumntype{R}{>{\raggedleft\arraybackslash}X}
\usepackage{array}
\usepackage{multicol}
\usepackage{siunitx}
% parse-numbers=false : accepte aussi \SI{2,5 \times 10^{-3}}{...}
\sisetup{locale=UK,output-decimal-marker={.},group-minimum-digits=4,parse-numbers=false}
\DeclareSIUnit\ohm{\ensuremath{\symup{\Omega}}} % U+2126 absent de la police
\DeclareSIUnit\division{div} % réglages d'oscilloscope (ms/div, V/div)
\DeclareSIUnit\tr{tr} % tours (vitesse de rotation, ex. \SI{1500}{\tr\per\minute})
\usepackage{qrcode}
% \degree : commande fréquemment utilisée par les modèles pour un angle ou une
% température en dehors de \SI{}{} (non fournie par les paquets chargés).
\providecommand{\degree}{^{\circ}}
% \qed (fin de démonstration), utilisé par les modèles sans amsthm
\providecommand{\qed}{\hfill$\square$}
% \barre : double barre des valeurs interdites dans un tableau de variations (tkz-tab)
\providecommand{\barre}{\ensuremath{\|}}
% environnement proof (amsthm non chargé) utilisé par les modèles
\ifdefined\proof\else\newenvironment{proof}[1][Proof]{\par\noindent\textit{#1.}\ }{\qed\par}\fi
\usepackage{lastpage}
\usepackage{hyperref}
\hypersetup{hidelinks}

% ============================================================
% Palette « Pop » OnBuch+ — mêmes hex que la classe P (plus_ui.dart)
% ============================================================
\definecolor{poporange}{HTML}{F59321}
\definecolor{popdark}{HTML}{E07A0C}
\definecolor{popcream}{HTML}{FFF6E8}
\definecolor{popink}{HTML}{1C1714}
\definecolor{poppurple}{HTML}{7A5AE0}
\definecolor{popgreen}{HTML}{2E9E6A}
\definecolor{poppink}{HTML}{E0457B}
\definecolor{popblue}{HTML}{2F7DE1}
\definecolor{popgold}{HTML}{C98A12}
\definecolor{popmuted}{HTML}{8A7F76}
\definecolor{popline}{HTML}{EADFD2}
\definecolor{popgray}{HTML}{8A7F76}
\colorlet{popgrey}{popgray}
% Alias tolérants (le modèle nomme parfois les couleurs différemment)
\colorlet{popwhite}{white}
\colorlet{popblack}{popink}
\colorlet{popred}{poppink}
\colorlet{popyellow}{popgold}
% Teintes claires pour fonds de tableaux / zones
\colorlet{popblueL}{popblue!10!white}
\colorlet{poporangeL}{poporange!14!white}
\colorlet{popgreenL}{popgreen!12!white}
\colorlet{poppurpleL}{poppurple!12!white}
\colorlet{poppinkL}{poppink!12!white}
\colorlet{popgoldL}{popgold!14!white}

\pagecolor{popcream}

% ============================================================
% Typographie
% ============================================================
\defaultfontfeatures{Ligatures=TeX}
\setmainfont{PlusJakartaSans-Medium.ttf}[Path=../../../fonts/,BoldFont=PlusJakartaSans-Bold.ttf]
% Maths en sans-serif (Fira Math) avec chiffres et lettres droites de Plus
% Jakarta, pour que les formules se fondent dans le texte.
\usepackage[math-style=ISO,bold-style=ISO]{unicode-math}
\setmathfont{FiraMath-Regular.otf}[Path=../../../fonts/]
\setmathfont{PlusJakartaSans-Medium.ttf}[Path=../../../fonts/,range={up/num,up/{Latin,latin},\mathup}]
% (icomma retiré : décimales avec point en anglais)
% Caractères mathématiques Unicode absents des polices de texte (Plus Jakarta,
% Archivo Black) : rendus par leur équivalent mathématique, en texte comme en
% formule — sinon ils s'affichent en carré vide (ex. « Arithmétique dans ℤ »).
\usepackage{newunicodechar}
\newunicodechar{ℝ}{\ensuremath{\mathbb{R}}}
\newunicodechar{ℤ}{\ensuremath{\mathbb{Z}}}
\newunicodechar{ℂ}{\ensuremath{\mathbb{C}}}
\newunicodechar{ℕ}{\ensuremath{\mathbb{N}}}
\newunicodechar{ℚ}{\ensuremath{\mathbb{Q}}}
\newunicodechar{𝔻}{\ensuremath{\mathbb{D}}}
\newunicodechar{α}{\ensuremath{\alpha}}
\newunicodechar{β}{\ensuremath{\beta}}
\newunicodechar{γ}{\ensuremath{\gamma}}
\newunicodechar{δ}{\ensuremath{\delta}}
\newunicodechar{ε}{\ensuremath{\varepsilon}}
\newunicodechar{θ}{\ensuremath{\theta}}
\newunicodechar{λ}{\ensuremath{\lambda}}
\newunicodechar{μ}{\ensuremath{\mu}}
\newunicodechar{σ}{\ensuremath{\sigma}}
\newunicodechar{τ}{\ensuremath{\tau}}
\newunicodechar{φ}{\ensuremath{\varphi}}
\newunicodechar{ψ}{\ensuremath{\psi}}
\newunicodechar{ω}{\ensuremath{\omega}}
\newunicodechar{Σ}{\ensuremath{\Sigma}}
% majuscules produites par \MakeUppercase dans les titres (α-aminés -> Α)
\newunicodechar{Α}{\ensuremath{\alpha}}
\newunicodechar{Β}{\ensuremath{\beta}}
\newunicodechar{Γ}{\ensuremath{\Gamma}}
\newunicodechar{Δ}{\ensuremath{\Delta}}
\newunicodechar{Ε}{\ensuremath{\varepsilon}}
\newunicodechar{Θ}{\ensuremath{\Theta}}
\newunicodechar{Λ}{\ensuremath{\Lambda}}
\newunicodechar{Μ}{\ensuremath{\mu}}
\newunicodechar{Τ}{\ensuremath{\tau}}
\newunicodechar{Φ}{\ensuremath{\Phi}}
\newunicodechar{Ψ}{\ensuremath{\Psi}}
\newunicodechar{Ω}{\ensuremath{\Omega}}
\newunicodechar{↦}{\ensuremath{\mapsto}}
\newunicodechar{∈}{\ensuremath{\in}}
\newunicodechar{∉}{\ensuremath{\notin}}
\newunicodechar{∧}{\ensuremath{\wedge}}
\newunicodechar{⇔}{\ensuremath{\Leftrightarrow}}
\newunicodechar{⟺}{\ensuremath{\Longleftrightarrow}}
\newunicodechar{⇒}{\ensuremath{\Rightarrow}}
\newunicodechar{⟹}{\ensuremath{\Longrightarrow}}
\newunicodechar{∘}{\ensuremath{\circ}}
\newunicodechar{∪}{\ensuremath{\cup}}
\newunicodechar{∩}{\ensuremath{\cap}}
\newunicodechar{≡}{\ensuremath{\equiv}}
\newunicodechar{⊂}{\ensuremath{\subset}}
\newunicodechar{⊥}{\ensuremath{\perp}}
\newunicodechar{∃}{\ensuremath{\exists}}
\newunicodechar{∀}{\ensuremath{\forall}}
\newunicodechar{∛}{\ensuremath{\sqrt[3]{\,}}}
\newunicodechar{∜}{\ensuremath{\sqrt[4]{\,}}}
\newunicodechar{⃗}{\ensuremath{{}^{\to}}}
\newunicodechar{ⁿ}{\ifmmode^{n}\else\textsuperscript{\ensuremath{n}}\fi}
\newunicodechar{ˣ}{\ifmmode^{x}\else\textsuperscript{\ensuremath{x}}\fi}
\newunicodechar{ᵇ}{\ifmmode^{b}\else\textsuperscript{\ensuremath{b}}\fi}
\newunicodechar{ʳ}{\ifmmode^{r}\else\textsuperscript{\ensuremath{r}}\fi}
\newunicodechar{ᵘ}{\ifmmode^{u}\else\textsuperscript{\ensuremath{u}}\fi}
\newunicodechar{ʸ}{\ifmmode^{y}\else\textsuperscript{\ensuremath{y}}\fi}
\newunicodechar{ᵃ}{\ifmmode^{a}\else\textsuperscript{\ensuremath{a}}\fi}
\newunicodechar{ᵉ}{\ifmmode^{e}\else\textsuperscript{\ensuremath{e}}\fi}
\newunicodechar{ᵏ}{\ifmmode^{k}\else\textsuperscript{\ensuremath{k}}\fi}
\newunicodechar{ᵖ}{\ifmmode^{p}\else\textsuperscript{\ensuremath{p}}\fi}
\newunicodechar{ᴺ}{\ifmmode^{N}\else\textsuperscript{\ensuremath{N}}\fi}
\newunicodechar{ᶜ}{\ifmmode^{c}\else\textsuperscript{\ensuremath{c}}\fi}
\newunicodechar{ʰ}{\ifmmode^{h}\else\textsuperscript{\ensuremath{h}}\fi}
\newunicodechar{ᵠ}{\ifmmode^{q}\else\textsuperscript{\ensuremath{q}}\fi}
\newunicodechar{ₙ}{\ifmmode_{n}\else\textsubscript{\ensuremath{n}}\fi}
\newunicodechar{ₐ}{\ifmmode_{a}\else\textsubscript{\ensuremath{a}}\fi}
\newunicodechar{ᵢ}{\ifmmode_{i}\else\textsubscript{\ensuremath{i}}\fi}
\newunicodechar{ᵣ}{\ifmmode_{r}\else\textsubscript{\ensuremath{r}}\fi}
\newunicodechar{ₖ}{\ifmmode_{k}\else\textsubscript{\ensuremath{k}}\fi}
\newunicodechar{ᵦ}{\ifmmode_{\beta}\else\textsubscript{\ensuremath{\beta}}\fi}
\newunicodechar{ₓ}{\ifmmode_{x}\else\textsubscript{\ensuremath{x}}\fi}
\newfontfamily\popdisplay{ArchivoBlack-Regular.ttf}[Path=../../../fonts/]
\newfontfamily\poplogo{SpaceGrotesk-Bold.ttf}[Path=../../../fonts/]
\newfontfamily\popsemi{PlusJakartaSans-SemiBold.ttf}[Path=../../../fonts/]
\newfontfamily\popxbold{PlusJakartaSans-ExtraBold.ttf}[Path=../../../fonts/]
\newfontfamily\popxboldls{PlusJakartaSans-ExtraBold.ttf}[Path=../../../fonts/,LetterSpace=3.0]

\setlist[enumerate]{leftmargin=1.7em,itemsep=5pt,topsep=4pt}
\setlist[itemize]{leftmargin=1.7em,itemsep=4pt,topsep=4pt,label={\color{poporange}\textbullet}}
\setlength{\parindent}{0pt}
\setlength{\parskip}{5pt}
\renewcommand{\arraystretch}{1.25}

% pgfplots : style Pop par défaut
\pgfplotsset{
  every axis/.append style={axis line style={popink,line width=1pt},tick style={popink},
    grid style={popline,line width=0.5pt},label style={font=\small},tick label style={font=\footnotesize}},
  popcurve/.style={poporange,line width=1.8pt,no markers,smooth},
}
% Même style disponible dans un \draw TikZ simple (hors pgfplots)
\tikzset{popcurve/.style={poporange,line width=1.8pt}}
\tikzset{popgrid/.style={popline,very thin}}
\colorlet{popcurve}{poporange} % le nom du style est parfois utilisé comme couleur

% ============================================================
% Pill / squiggle
% ============================================================
\newcommand{\poppill}[3]{%
  \tikz[baseline=-0.55ex]\node[draw=popink,line width=1.2pt,rounded corners=2.6pt,%
    fill=#2,inner xsep=6.5pt,inner ysep=3pt]{%
    \popxboldls\fontsize{7}{7}\selectfont\color{#3}\MakeUppercase{#1}};%
}
\newcommand{\popsquiggle}[1]{%
  \tikz[baseline=-0.4ex]{
    \draw[#1,line width=2.6pt,line cap=round]
      plot [smooth] coordinates {(0,0.09) (0.34,0.01) (0.68,0.21) (1.02,0.01) (1.36,0.10)};
  }%
}
% Mot-clé surligné (marqueur orange sous le texte)
\newcommand{\cle}[1]{{\popxbold\color{popdark}#1}}

% ============================================================
% En-tête / pied
% ============================================================
\pagestyle{fancy}
\fancyhf{}
\renewcommand{\headrulewidth}{0pt}
\renewcommand{\footrulewidth}{0pt}
\lhead{\raisebox{-0.15ex}{\poplogo\fontsize{13}{13}\selectfont\color{popink}OnBuch\color{poporange}+}}
\rhead{\poppill{\DOCMATIERE\ \textbullet\ \DOCNIVEAU\ \textbullet\ Chapter \DOCLECON}{white}{popink}}
\lfoot{\popxbold\fontsize{7.6}{7.6}\selectfont\color{popmuted}\MakeUppercase{Signed course OnBuch+}\ \textbullet\ {\popsemi\DOCTITRE}}
\rfoot{\tikz[baseline=-0.6ex]\node[rounded corners=8.5pt,fill=popink,minimum height=17pt,minimum width=17pt,inner xsep=5pt,inner sep=0]{\popxbold\fontsize{8}{8}\selectfont\color{popcream}\thepage\,/\,\pageref*{LastPage}};}

% ============================================================
% Titres
% ============================================================
\newcommand{\chaptitre}[2]{%
  \begin{tcolorbox}[enhanced,colback=poporange,colframe=popink,boxrule=2.6pt,arc=11pt,
      drop shadow=popink,left=15pt,right=15pt,top=12pt,bottom=11pt,before skip=3pt,after skip=18pt]
    \poppill{#2}{popink}{popcream}\par\vspace{9pt}
    {\raggedright\hyphenpenalty=10000\exhyphenpenalty=10000\popdisplay\fontsize{22}{24}\selectfont\color{popcream}\MakeUppercase{#1}\par}
  \end{tcolorbox}
}
% Section numérotée : pastille orange avec le numéro + titre Archivo
\newcounter{popsec}
\newcommand{\coursec}[1]{%
  \stepcounter{popsec}\setcounter{popsub}{0}%
  \par\vspace{16pt}\noindent
  \begin{minipage}{\linewidth}
  \tikz[baseline=-0.7ex]\node[draw=popink,line width=1.6pt,fill=poporange,rounded corners=4pt,minimum size=22pt,inner sep=0]{\popdisplay\fontsize{12}{12}\selectfont\color{popcream}\thepopsec};%
  \hspace{8pt}{\popdisplay\fontsize{15}{17}\selectfont\color{popink}\MakeUppercase{#1}}\par
  \vspace{4pt}\noindent{\color{poporange}\rule{36pt}{3.6pt}}
  \end{minipage}\par\nopagebreak\vspace{8pt}
}
\newcounter{popsub}
\newcommand{\courssub}[1]{%
  \stepcounter{popsub}%
  \par\vspace{9pt}\noindent{\popxbold\fontsize{11.5}{13}\selectfont\color{popink}{\color{poporange}\thepopsec.\thepopsub}\hspace{6pt}#1}\par\nopagebreak\vspace{4pt}
}

% ============================================================
% Boîtes pédagogiques — même recette : fond blanc, bordure épaisse colorée,
% ombre dure encre, badge plein. Titre optionnel [#1].
% ============================================================
\tcbset{popbase/.style={breakable,enhanced,colback=white,boxrule=2.3pt,arc=9pt,
  shadow={3pt}{-3pt}{0pt}{popink},left=14pt,right=14pt,top=11pt,bottom=11pt,before skip=11pt,after skip=15pt,
  fontupper=\popsemi\fontsize{10.6}{14.5}\selectfont\color{popink},
  fontlower=\popsemi\fontsize{10.6}{14.5}\selectfont\color{popink}}}
% Coche dessinée (le glyphe ✓ n'existe pas dans les polices du document)
\newcommand{\popcheck}[1]{\tikz[baseline=-0.5ex]\draw[#1,line width=1.6pt,line cap=round,line join=round](-0.13,0.0)--(-0.04,-0.09)--(0.13,0.1);}
\providecommand{\coche}{\popcheck{popgreen}}
\providecommand{\resultat}[1]{\par\smallskip\noindent\fcolorbox{popgreen}{popgreenL}{\parbox{\dimexpr\linewidth-2\fboxsep-2\fboxrule\relax}{\textbf{Result.} #1}}\par\smallskip}
\newcommand{\popboxhead}[3]{\poppill{#1}{#2}{white}\ifx\relax#3\relax\else\hspace{6pt}{\popxbold\fontsize{10.6}{12}\selectfont\color{popink}#3}\fi\par\vspace{7pt}}

\newtcolorbox{aretenir}[1][]{popbase,colframe=popgreen,before upper={\popboxhead{Key points}{popgreen}{#1}}}
\newtcolorbox{exemplebox}[1][]{popbase,colframe=poporange,before upper={\popboxhead{Exemple}{poporange}{#1}}}
\newtcolorbox{definition}[1][]{popbase,colframe=popblue,before upper={\popboxhead{Definition}{popblue}{#1}}}
\newtcolorbox{propriete}[1][]{popbase,colframe=poppurple,before upper={\popboxhead{Property}{poppurple}{#1}}}
\newtcolorbox{methode}[1][]{popbase,colframe=popgold,before upper={\popboxhead{Method}{popgold}{#1}}}
\newtcolorbox{attention}[1][]{popbase,colframe=poppink,before upper={\popboxhead{Warning}{poppink}{#1}}}
\newtcolorbox{experience}[1][]{popbase,colframe=popblue,colback=popblueL,before upper={\popboxhead{Experiment}{popblue}{#1}}}
% « Le savais-tu ? » : carte violette pleine (contexte, culture, Cameroun)
\newtcolorbox{savaistu}[1][]{popbase,colback=poppurple,colframe=popink,
  fontupper=\popsemi\fontsize{10.4}{14.2}\selectfont\color{white},
  before upper={\poppill{Did you know?}{white}{poppurple}\ifx\relax#1\relax\else\hspace{6pt}{\popxbold\color{white}#1}\fi\par\vspace{7pt}}}
% Exercice résolu : énoncé en haut, \tcblower puis la solution sur fond crème
\newcounter{popexo}\newcounter{popexr}
\newtcolorbox{exoresolu}[1][]{popbase,colframe=poporange,colbacklower=popcream,
  segmentation style={popink,line width=1.2pt,dashed},
  before upper={\stepcounter{popexr}\popboxhead{Worked example \thepopexr}{poporange}{#1}},
  before lower={\poppill{Solution}{popink}{popcream}\par\vspace{6pt}}}
% Exercice (non résolu en place) — la correction est dans la partie Corrigés
\newtcolorbox{exercice}[1][]{popbase,colframe=popink,
  before upper={\stepcounter{popexo}\popboxhead{Exercise \thepopexo}{popink}{#1}}}
% Corrigé
\newtcolorbox{corrige}[1][]{popbase,colframe=popgreen,colback=popgreenL,
  before upper={\popboxhead{Solution}{popgreen}{#1}}}
% Objectifs / prérequis / bilan
\newtcolorbox{objectifs}{popbase,colframe=popink,colback=poporangeL,
  before upper={\popboxhead{Chapter objectives}{popink}{}}}
\newtcolorbox{prerequis}{popbase,colframe=popink,colback=popblueL,
  before upper={\popboxhead{Prerequisites}{popblue}{}}}
\newtcolorbox{bilan}{popbase,colframe=popink,colback=white,boxrule=2.8pt,
  before upper={\popboxhead{Fiche bilan}{poporange}{L'essentiel à connaître pour l'examen}}}
% Figure encadrée avec légende
\newtcolorbox{popfigure}[1][]{enhanced,breakable=false,colback=white,colframe=popline,boxrule=1.4pt,arc=8pt,
  left=10pt,right=10pt,top=10pt,bottom=8pt,before skip=10pt,after skip=12pt,halign=center,
  lower separated=false,halign lower=center,
  fontlower=\popsemi\fontsize{9}{11.5}\selectfont\color{popmuted},
  #1}
\newcommand{\legende}[1]{\par\vspace{6pt}{\popsemi\fontsize{9}{11.5}\selectfont\color{popmuted}#1}}

% ============================================================
% Signature OnBuch+
% ============================================================
\newcommand{\poplogoinline}[1]{{\poplogo\fontsize{#1}{#1}\selectfont\color{popink}OnBuch\color{poporange}+}}
\newcommand{\signatureonbuch}{%
  \begin{tcolorbox}[enhanced,colback=popink,colframe=popink,boxrule=2.2pt,arc=10pt,
      drop shadow=poporange,left=14pt,right=14pt,top=10pt,bottom=10pt,before skip=0pt,after skip=0pt]
    \tikz[baseline=-0.7ex]\node[circle,fill=popgreen,draw=popcream,line width=1.3pt,minimum size=22pt,inner sep=0]{\popcheck{white}};%
    \hspace{9pt}\begin{minipage}[c]{0.62\linewidth}
      {\popxbold\fontsize{12}{13}\selectfont\color{popcream}Signed\ \textbullet\ {\poplogo OnBuch\color{poporange}+}}\par\vspace{2pt}
      {\raggedright\popsemi\fontsize{8}{10.5}\selectfont\color{popline}\MakeUppercase{OnBuch+ teaching team}\\\MakeUppercase{Follows the official MINESEC syllabus}\par}
    \end{minipage}\hfill
    {\poplogo\fontsize{22}{22}\selectfont\color{popcream}OnBuch\color{poporange}+}
  \end{tcolorbox}%
}

% ============================================================
% Page de garde
% #1 titre · #2 matière · #3 niveau · #4 module · #5 n° leçon · #6 durée ·
% #7 description / objectifs (texte ou itemize)
% ============================================================
\newcommand{\pagegarde}[7]{%
  \thispagestyle{empty}
  \newgeometry{a4paper,margin=1.7cm,top=1.6cm,bottom=1.4cm}%
  \begin{tikzpicture}[remember picture,overlay]
    \fill[poporange!18] ([xshift=-3.2cm,yshift=-3.6cm]current page.north east) circle (4.6cm);
    \fill[poppurple!14] ([xshift=2.4cm,yshift=7.6cm]current page.south west) circle (3.8cm);
  \end{tikzpicture}%
  {\poplogo\fontsize{30}{30}\selectfont\color{popink}OnBuch\color{poporange}+}\hfill
  \raisebox{0.6ex}{\poppill{Signed course \textbullet\ Premium}{popink}{popcream}}\par
  \vspace{1pt}\popsquiggle{poppurple}\par
  \vspace{8pt}
  \begin{tcolorbox}[enhanced,colback=poporange,colframe=popink,boxrule=2.8pt,arc=13pt,
      drop shadow=popink,left=18pt,right=18pt,top=12pt,bottom=13pt,after skip=10pt]
    \poppill{#2 \textbullet\ #3}{popink}{popcream}\hspace{5pt}\poppill{#4}{white}{popink}\par\vspace{8pt}
    {\popdisplay\fontsize{14}{14}\selectfont\color{popink}CHAPTER #5}\par\vspace{4pt}
    {\raggedright\hyphenpenalty=10000\exhyphenpenalty=10000\popdisplay\fontsize{25}{27}\selectfont\color{popcream}\MakeUppercase{#1}\par}
    \vspace{8pt}
    {\popxbold\fontsize{9.5}{9.5}\selectfont\color{popink}\MakeUppercase{Suggested time: #6}}
  \end{tcolorbox}
  \begin{tcolorbox}[enhanced,colback=white,colframe=popink,boxrule=2.2pt,arc=10pt,
      drop shadow=popink,left=16pt,right=16pt,top=10pt,bottom=10pt,after skip=8pt]
    \poppill{In this chapter}{poppurple}{white}\par\vspace{5pt}
    \popsemi\fontsize{9.3}{12.6}\selectfont\color{popink}#7
  \end{tcolorbox}
  \noindent\begin{minipage}[t]{0.487\linewidth}
    \begin{tcolorbox}[enhanced,colback=popgreen,colframe=popink,boxrule=2.2pt,arc=10pt,
        drop shadow=popink,left=11pt,right=11pt,top=10pt,bottom=10pt,height=3.1cm,valign=center]
      \tcbox[colback=white,colframe=popink,boxrule=1.3pt,arc=5pt,boxsep=0pt,nobeforeafter,
        left=3pt,right=3pt,top=3pt,bottom=3pt,valign=center]{\qrcode[height=1.9cm]{https://play.google.com/store/apps/details?id=com.luvvix.onbuch&hl=en}}%
      \hspace{8pt}\begin{minipage}[c]{0.5\linewidth}
        {\popdisplay\fontsize{11}{12.5}\selectfont\color{white}\MakeUppercase{Download OnBuch}}\par\vspace{4pt}
        {\raggedright\popsemi\fontsize{8.4}{11}\selectfont\color{white}Free \textbullet\ Google Play\\AI tutor, past papers, results\par}
      \end{minipage}
    \end{tcolorbox}
  \end{minipage}\hfill
  \begin{minipage}[t]{0.487\linewidth}
    \begin{tcolorbox}[enhanced,colback=poppurple,colframe=popink,boxrule=2.2pt,arc=10pt,
        drop shadow=popink,left=11pt,right=11pt,top=10pt,bottom=10pt,height=3.1cm,valign=center]
      \tikz[baseline=-0.7ex]\node[circle,fill=white,draw=popink,line width=1.3pt,minimum size=1.6cm,inner sep=0]{\popdisplay\fontsize{24}{24}\selectfont\color{poppurple}+};%
      \hspace{8pt}\begin{minipage}[c]{0.6\linewidth}
        {\popdisplay\fontsize{11}{12.5}\selectfont\color{white}\MakeUppercase{Go OnBuch+}}\par\vspace{4pt}
        {\raggedright\popsemi\fontsize{8.4}{11}\selectfont\color{white}All signed courses, unlimited AI tutor, PDF sheets — OnBuch+ tab in the app\par}
      \end{minipage}
    \end{tcolorbox}
  \end{minipage}
  \vspace{8pt}
  \popcontenu
  \vfill
  \signatureonbuch
  \restoregeometry
  \newpage
}

% Rangée « Ce cours contient » (page de garde)
\newcommand{\poptile}[3]{% #1 couleur #2 gros texte #3 légende
  \begin{tcolorbox}[enhanced,colback=#1,colframe=popink,boxrule=2pt,arc=9pt,shadow={2.5pt}{-2.5pt}{0pt}{popink},
      left=6pt,right=6pt,top=9pt,bottom=9pt,halign=center,valign=center,height=1.75cm,width=\linewidth,nobeforeafter]
    {\popdisplay\fontsize{12.5}{13}\selectfont\color{white}\MakeUppercase{#2}}\par\vspace{3pt}
    {\centering\popsemi\fontsize{7.8}{9.5}\selectfont\color{white}#3\par}
  \end{tcolorbox}}
\newcommand{\popcontenu}{%
  \par\noindent{\popxboldls\fontsize{7.5}{7.5}\selectfont\color{popmuted}THIS SIGNED COURSE CONTAINS}\par\vspace{7pt}
  \noindent\begin{minipage}[t]{0.235\linewidth}\poptile{popblue}{Course}{complete, illustrated, step by step}\end{minipage}\hfill
  \begin{minipage}[t]{0.235\linewidth}\poptile{poporange}{Methods}{and worked examples}\end{minipage}\hfill
  \begin{minipage}[t]{0.235\linewidth}\poptile{poppink}{Exercises}{exam style + solutions}\end{minipage}\hfill
  \begin{minipage}[t]{0.235\linewidth}\poptile{popgreen}{Summary sheet}{the essentials to remember}\end{minipage}\par}

% Sommaire
\newtcolorbox{sommaire}{enhanced,colback=white,colframe=popink,boxrule=2.2pt,arc=10pt,
  drop shadow=popink,left=18pt,right=18pt,top=13pt,bottom=13pt,before skip=6pt,after skip=20pt,
  before upper={\poppill{Contents}{poppurple}{white}\par\vspace{9pt}\popsemi\fontsize{10.6}{14.5}\selectfont\color{popink}}}

% Signature de fin de cours — poussée tout en bas de la dernière page
\newcommand{\findecours}{%
  \par\vfill\nopagebreak
  \begin{center}\popsquiggle{poporange}\end{center}\vspace{4pt}
  \signatureonbuch
}
\AtBeginDocument{\def\llbracket{\mathopen{[\mkern-3.5mu[}}\def\rrbracket{\mathclose{]\mkern-3.5mu]}}} % intervalles d'entiers (glyphe ⟦ absent de la police)
% Glyphes absents de Fira Math : redéfinis à partir de glyphes disponibles
\AtBeginDocument{%
  \def\setminus{\mathbin{\backslash}}\let\smallsetminus\setminus
  \def\vdots{\mathord{\vbox{\baselineskip3.2pt\lineskiplimit0pt\kern4pt\hbox{.}\hbox{.}\hbox{.}}}}%
  \def\ddots{\mathinner{\mkern1mu\raise6.5pt\hbox{.}\mkern2mu\raise3.6pt\hbox{.}\mkern2mu\raise0.7pt\hbox{.}\mkern1mu}}%
  \def\triangle{\mathord{\scalebox{0.9}{$\symup{\Delta}$}}}%
  \def\nabla{\mathord{\raisebox{\depth}{\scalebox{1}[-1]{$\symup{\Delta}$}}}}%
  \def\checkmark{\popcheck{popdark}}%
  \def\star{\mathbin{\ast}}\def\bigstar{\mathord{\ast}}%
  \def\diamond{\mathbin{\scalebox{0.6}{$\Diamond$}}}%
  \def\bigcup{\mathop{\mathchoice{\raisebox{-0.3em}{\scalebox{1.7}{$\cup$}}}{\scalebox{1.25}{$\cup$}}{\cup}{\cup}}\displaylimits}%
  \def\bigcap{\mathop{\mathchoice{\raisebox{-0.3em}{\scalebox{1.7}{$\cap$}}}{\scalebox{1.25}{$\cap$}}{\cap}{\cap}}\displaylimits}%
  \def\lesseqqgtr{\mathrel{\lessgtr}}\def\bigoplus{\mathop{\oplus}\displaylimits}%
}
\providecommand{\ssi}{if and only if} % abréviation fréquente
\AtBeginDocument{\providecommand{\wideparen}{\overparen}} % arc de cercle

% Fonctions usuelles que les modèles utilisent sans que amsmath les fournisse
\providecommand{\arsinh}{\operatorname{arsinh}}\providecommand{\arcosh}{\operatorname{arcosh}}
\providecommand{\artanh}{\operatorname{artanh}}\providecommand{\arsech}{\operatorname{arsech}}
\providecommand{\arcsch}{\operatorname{arcsch}}\providecommand{\arcoth}{\operatorname{arcoth}}
\providecommand{\arcsinh}{\operatorname{arsinh}}\providecommand{\arccosh}{\operatorname{arcosh}}
\providecommand{\arctanh}{\operatorname{artanh}}\providecommand{\sech}{\operatorname{sech}}
\providecommand{\csch}{\operatorname{csch}}\providecommand{\arcsec}{\operatorname{arcsec}}
\providecommand{\arccsc}{\operatorname{arccsc}}\providecommand{\arccot}{\operatorname{arccot}}
\providecommand{\sgn}{\operatorname{sgn}}\providecommand{\lcm}{\operatorname{lcm}}
\providecommand{\Var}{\operatorname{Var}}\providecommand{\Cov}{\operatorname{Cov}}

%% FIGFIX-BEGIN v5 — correctifs globaux des figures (docs/onbuch-generation/tools/figfix)
\usepackage{adjustbox}\usepackage{etoolbox}\tcbuselibrary{raster}
% toute figure TikZ/pgfplots plus large que la ligne est réduite à la largeur disponible
\BeforeBeginEnvironment{tikzpicture}{\begin{adjustbox}{max width=\linewidth}}
\AfterEndEnvironment{tikzpicture}{\end{adjustbox}}
% légendes pgfplots lisibles par-dessus les courbes
\pgfplotsset{every axis/.append style={legend style={fill=white,fill opacity=0.92,text opacity=1,draw=black!15,inner sep=3pt}}}
% étiquettes d'arêtes (syntaxe edge["..."]) avec fond blanc
\tikzset{every edge quotes/.append style={fill=white,inner sep=1.2pt}}
%% FIGFIX-END
\begin{document}

\pagegarde{Chap2: EMPHASIS ON STORYLINE, THE PLOT, THE SITUATION, ETC}{Literature in English}{Upper Sixth}{Upper Sixth — Literature in English}{2}{14 periods}{This chapter deepens the learner's mastery of narrative analysis by following the plot development of three set texts: Orwell's Nineteen Eighty-Four (Parts 2 and 3), Chaucer's General Prologue and Merchant's Tale, and Hansberry's A Raisin in the Sun (Act 2). Through close reading of storyline, plot structure and situation, learners connect literary events to themes of totalitarianism, social satire, race and the American Dream.
\vspace{4pt}
\begin{multicols}{2}\small
\begin{itemize}
\item By the end of this chapter I can summarise and analyse the plot of Nineteen Eighty-Four Part 2 (the love affair, the Book, the arrest) and Part 3 (the Ministry of Love, Room 101).
\item By the end of this chapter I can explain how Orwell uses storyline and situation to denounce totalitarian manipulation of truth, love and thought.
\item By the end of this chapter I can characterise Chaucer's pilgrims (Cook, Manciple, Shipman, Doctor, Wife of Bath, Parson, Plowman, Summoner, Pardoner, Host) and link each portrait to social satire.
\item By the end of this chapter I can analyse the structure, prologue and marriage debate of The Merchant's Tale, including January's delusion and the advice of Placebo and Justinus.
\item By the end of this chapter I can trace the plot of A Raisin in the Sun Act 2 and relate its climaxes (the lost money, the Clybourne Park decision) to the American Dream, racism and cultural identity.
\item By the end of this chapter I can distinguish plot, storyline, situation, climax and sub-plot, and use these terms accurately in an essay.
\end{itemize}\end{multicols}}

\begin{sommaire}
\begin{multicols}{2}
\begin{enumerate}
\item Section 1: Nineteen Eighty-Four Part 2 — Love as Rebellion (Chapters 7–9)
\item Section 2: Nineteen Eighty-Four Part 3 — The Destruction of the Self (Chapters 1–4)
\item Section 3: The General Prologue I — Skilled Workers, Traders and Professionals
\item Section 4: The General Prologue II — Women, Ideal Christians and Corrupt Churchmen
\item Section 5: The Merchant's Tale — Structure, Prologue and the Marriage Debate
\item Section 6: A Raisin in the Sun Act 2 Scene 1 — American Dream, Money and Cultural Identity
\item Section 7: A Raisin in the Sun Act 2 Scenes 2–3 — Loss, Racism and the Clybourne Park Decision
\item Integration activity
\item Methods and worked examples
\item Exercises
\item Solutions to the exercises
\item Summary sheet
\end{enumerate}
\end{multicols}
\end{sommaire}

\begin{objectifs}
\begin{itemize}
\item By the end of this chapter I can summarise and analyse the plot of Nineteen Eighty-Four Part 2 (the love affair, the Book, the arrest) and Part 3 (the Ministry of Love, Room 101).
\item By the end of this chapter I can explain how Orwell uses storyline and situation to denounce totalitarian manipulation of truth, love and thought.
\item By the end of this chapter I can characterise Chaucer's pilgrims (Cook, Manciple, Shipman, Doctor, Wife of Bath, Parson, Plowman, Summoner, Pardoner, Host) and link each portrait to social satire.
\item By the end of this chapter I can analyse the structure, prologue and marriage debate of The Merchant's Tale, including January's delusion and the advice of Placebo and Justinus.
\item By the end of this chapter I can trace the plot of A Raisin in the Sun Act 2 and relate its climaxes (the lost money, the Clybourne Park decision) to the American Dream, racism and cultural identity.
\item By the end of this chapter I can distinguish plot, storyline, situation, climax and sub-plot, and use these terms accurately in an essay.
\item By the end of this chapter I can write a coherent GCE-style essay analysing how a writer uses plot development to convey theme.
\item By the end of this chapter I can compare how the three set texts use narrative situations to criticise their societies.
\end{itemize}
\end{objectifs}

\begin{prerequis}
\begin{itemize}
\item Plot elements (exposition, rising action, climax, falling action, resolution) studied in Lower Sixth
\item Nineteen Eighty-Four Part 1: Winston's world, the Party, Big Brother, the telescreen
\item The General Prologue: Chaucer's frame narrative and the portraits already studied (Knight to Franklin)
\item A Raisin in the Sun Act 1: the Younger family, the insurance cheque, Mama's plant
\item Basic essay-writing technique: introduction, developed paragraphs, textual evidence, conclusion
\end{itemize}
\end{prerequis}

\newpage

\coursec{Section 1: Nineteen Eighty-Four Part 2 — Love as Rebellion (Chapters 7–9)}

Imagine a young man in Yaoundé who discovers that every message he sends is read, every friendship he forms is reported, and every dream he nurtures is reshaped by those in power. Now imagine a family in a cramped Douala neighbourhood waiting for a life‑changing sum of money, each member dreaming of a different future. Finally, picture a group of pilgrims on the road from Bamenda to Buea, each telling a story that secretly reveals who they truly are. These are not fantasies: they are the worlds of Orwell's \emph{Nineteen Eighty‑Four}, Hansberry's \emph{A Raisin in the Sun}, and Chaucer's \emph{Canterbury Tales}. In this chapter, we learn to follow a storyline closely, to trace how plot, situation and character interlock, and to see how great writers turn narrative into a mirror of society — skills the GCE examiner rewards above all.

\courssub{Recap of Part 1 Situation}

Before Part 2 opens, Winston Smith is a low‑ranking member of the Outer Party in Oceania. He works at the Ministry of Truth, where his job is to falsify historical records so that the Party's predictions always appear correct. In secret he keeps a diary — an act punishable by death — and writes \cle{Down with Big Brother}. He feels a fatalistic certainty: \emph{We are the dead}. His hatred of the Party is private, unspoken, and he believes the proles (the proletarian masses) are the only possible source of rebellion, yet they remain unconscious of their power. This situation establishes the \cle{exposition}: a solitary consciousness trapped in a total surveillance state, yearning for truth but convinced that resistance is futile.

\begin{definition}[Rising action]
The \textbf{rising action} is the series of events that build tension and complicate the central conflict after the exposition and before the climax. In a novel, it typically introduces new characters, deepens relationships, and raises the stakes so that the eventual climax carries maximum emotional weight.
\end{definition}

\courssub{Chapter 7 — Memory, the Proles and the Abolition of Private Loyalty}

Chapter 7 begins with Winston's dream of his mother and sister sinking into a dark abyss, a memory he has repressed because the Party demands that loyalty be directed solely to Big Brother. The dream forces him to confront the \cle{mutability of the past}: the Party erases personal history just as it rewrites public records. Winston realises that the proles, though politically ignorant, retain a humanity the Party cannot fully extinguish — they still love, sing, and remember. This insight becomes the intellectual foundation for his later rebellion.

\begin{aretenir}[Key points from Chapter 7]
\begin{itemize}
\item The dream of the mother and sister symbolises the Party's destruction of private affection.
\item The proles represent a reservoir of uncorrupted human feeling.
\item The Party's slogan \emph{Who controls the past controls the future} is dramatised through Winston's personal memory loss.
\end{itemize}
\end{aretenir}

\courssub{Chapter 8 — The Room above Charrington's Shop}

Winston rents the room above Mr Charrington's junk shop, believing it to be a sanctuary free from telescreens. Here he

\coursec{Section 2: Nineteen Eighty-Four Part 3 — The Destruction of the Self (Chapters 1–4)}

\courssub{Transition from Part 2: The Shattering of the Private World}

In Section 1 we followed Winston and Julia as they carved out a fragile sanctuary in the room above Mr Charrington’s shop. Their love affair, the paperweight, the nursery rhyme, and the illusion of a past that could be reclaimed — all these formed a \textbf{counter‑narrative} to the Party’s monopoly on truth. Part 2 ends with the brutal intrusion of the Thought Police: the paperweight is smashed, the telescreen behind the picture reveals itself, and the lovers are dragged apart. Part 3 opens not with a rescue or a rebellion, but with the cold, clinical machinery of the Ministry of Love. The transition is total: every symbol of hope from Part 2 is systematically destroyed, and the novel’s focus shifts from \emph{external resistance} to \emph{internal annihilation}.

\begin{aretenir}[Key Transition]
Part 2 builds the \emph{possibility} of freedom; Part 3 demonstrates the \emph{impossibility} of that freedom under totalitarianism. The climax of the novel is not a battle but a breakdown — Winston’s mind is the battlefield.
\end{aretenir}

\courssub{Chapter 1: The Ministry of Love — The Architecture of Terror}

Winston finds himself in a high‑windowed cell, monitored by four telescreens that leave no blind spot. The physical environment is designed to erase privacy: constant light, the hum of machinery, the impossibility of sleep. He meets other prisoners — the \emph{skull‑faced man} (a starving, broken intellectual), and \emph{Parsons}, who has been denounced by his own seven‑year‑old daughter for thoughtcrime uttered in his sleep. Parsons’ pride in his daughter’s vigilance illustrates the Party’s success in turning the family into an extension of the surveillance state.

The chapter’s climax is the arrival of \textbf{O’Brien}. Winston had believed O’Brien to be a fellow conspirator; now O’Brien appears as his interrogator and torturer. The revelation is not a twist of plot alone — it is the collapse of Winston’s last intellectual anchor. O’Brien says, \emph{“They got me a long time ago,”} implying that the Party’s power is so absolute that even the rebellion was a Party‑orchestrated trap.

\begin{definition}[Climax]
The \cle{climax} of a narrative is the turning point where the central conflict reaches its highest intensity. In \emph{Nineteen Eighty‑Four}, the climax is not a single dramatic event but the prolonged process of Winston’s capitulation under torture in the Ministry of Love. It is the moment when the protagonist’s will is broken and the Party’s victory becomes total.
\end{definition}

\begin{exemplebox}[The Skull‑Faced Man as a Mirror]
The skull‑faced man is a \emph{foil} to Winston: he has already undergone the process Winston is about to endure. His physical emaciation and mental surrender foreshadow Winston’s fate. When he offers Winston a piece of bread and is instantly punished, we see that even compassion is a crime in the Ministry of Love.
\end{exemplebox}

\courssub{Chapter 2: The Interrogation — Power as an End in Itself}

O’Brien begins the interrogation with a calm, almost pedagogical tone. He explains the Party’s philosophy: \emph{“The Party seeks power entirely for its own sake. We are not interested in the good of others; we are interested solely in power. Not wealth or luxury or long life or happiness: only power, pure power.”} This is the \cle{ideological core} of the novel. Power is not a means to an end; it is the end.

The famous \emph{fingers} scene dramatises the Party’s claim to control reality itself. O’Brien holds up four fingers and demands Winston see five. When Winston refuses, the pain increases. The point is not to force a lie but to force a \emph{genuine belief} that 2+2=5 if the Party says so. Reality is not external; it is whatever the Party declares.

\begin{propriete}[O’Brien’s Philosophy of Power]
\begin{itemize}
\item Power is collective and eternal — the individual dies, the Party is immortal.
\item Power is exercised by \emph{tearing human minds to pieces and putting them together again in new shapes of your own choosing}.
\item Objective truth does not exist; truth is what the Party says it is.
\item The past is mutable; records are continuously rewritten.
\end{itemize}
\end{propriete}

\begin{exoresolu}[Analysing the Fingers Scene]
\textbf{Question:} How does the “How many fingers?” episode illustrate the Party’s theory of reality control?
\textbf{Solution:}
\begin{enumerate}
\item \textbf{Physical coercion} — pain is used to break resistance.
\item \textbf{Psychological manipulation} — O’Brien does not want a forced confession; he wants Winston to \emph{see} five fingers.
\item \textbf{Epistemological claim} — the Party asserts that reality is inside the human mind, and the Party controls the mind.
\item \textbf{Result} — Winston eventually “sees” five fingers, demonstrating that objective reality has been replaced by Party‑defined reality.
\end{enumerate}
\end{exoresolu}

\begin{popfigure}
\begin{tikzpicture}[scale=0.9, every node/.style={font=\small}]
% Part 2 column
\node[draw, fill=popblueL, rounded corners, minimum width=5cm, minimum height=0.8cm] (p2title) at (0,4) {Part 2: Hope, Love, Privacy};
\node[draw, fill=popgreenL, rounded corners, minimum width=5cm, minimum height=0.6cm] (p2a) at (-2.5,3) {Paperweight (past)};
\node[draw, fill=popgreenL, rounded corners, minimum width=5cm, minimum height=0.6cm] (p2b) at (2.5,3) {Room above shop (sanctuary)};
\node[draw, fill=popgreenL, rounded corners, minimum width=5cm, minimum height=0.6cm] (p2c) at (-2.5,2) {Julia (love/rebellion)};
\node[draw, fill=popgreenL, rounded corners, minimum width=5cm, minimum height=0.6cm] (p2d) at (2.5,2) {Diary (private thought)};
\node[draw, fill=popgreenL, rounded corners, minimum width=5cm, minimum height=0.6cm] (p2e) at (0,1) {Nursery rhyme (culture)};

% Part 3 column
\node[draw, fill=poppink, rounded corners, minimum width=5cm, minimum height=0.8cm] (p3title) at (0,-1) {Part 3: Despair, Betrayal, Control};
\node[draw, fill=poporange, rounded corners, minimum width=5cm, minimum height=0.6cm] (p3a) at (-2.5,-2) {Paperweight smashed};
\node[draw, fill=poporange, rounded corners, minimum width=5cm, minimum height=0.6cm] (p3b) at (2.5,-2) {Room invaded by Thought Police};
\node[draw, fill=poporange, rounded corners, minimum width=5cm, minimum height=0.6cm] (p3c) at (-2.5,-3) {Julia betrayed / broken};
\node[draw, fill=poporange, rounded corners, minimum width=5cm, minimum height=0.6cm] (p3d) at (2.5,-3) {Diary used as evidence};
\node[draw, fill=poporange, rounded corners, minimum width=5cm, minimum height=0.6cm] (p3e) at (0,-4) {Rhyme completed by O’Brien (mockery)};

% Arrows showing destruction
\draw[->, thick, popdark] (p2a.south) -- (p3a.north);
\draw[->, thick, popdark] (p2b.south) -- (p3b.north);
\draw[->, thick, popdark] (p2c.south) -- (p3c.north);
\draw[->, thick, popdark] (p2d.south) -- (p3d.north);
\draw[->, thick, popdark] (p2e.south) -- (p3e.north);
\end{tikzpicture}
\legende{Contrast diagram: every symbol of hope in Part 2 is systematically destroyed in Part 3.}
\end{popfigure}

\courssub{Chapter 3: The Three Stages of Reintegration}

O’Brien outlines the process Winston must undergo: \textbf{learning}, \textbf{understanding}, and \textbf{acceptance}. This is not mere torture; it is a \emph{re‑education} designed to make Winston love Big Brother.

\begin{methode}[O’Brien’s Three Stages of Reintegration]
\begin{enumerate}
\item \textbf{Learning} — Winston is forced to confess to crimes he did not commit, to accept Party doctrine intellectually. He learns the “correct” answers.
\item \textbf{Understanding} — He begins to grasp the \emph{logic} of the Party: power for power’s sake, the mutability of the past, the non‑existence of objective truth. He sees \emph{why} the Party must control reality.
\item \textbf{Acceptance} — The final stage: Winston \emph{loves} Big Brother. His self is destroyed and rebuilt in the Party’s image. He no longer merely obeys; he believes.
\end{enumerate}
\end{methode}

The chapter culminates in O’Brien’s terrifying vision of the future: \emph{“If you want a picture of the future, imagine a boot stamping on a human face — forever.”} This image encapsulates the novel’s warning: totalitarianism does not merely oppress; it \emph{perfects} oppression by eliminating even the possibility of resistance.

\begin{aretenir}[The Boot Stamping on a Human Face]
This metaphor is the \cle{central image} of the novel’s political philosophy. It signifies:
\begin{itemize}
\item The permanence of Party rule (no revolution, no decay).
\item The intimacy of oppression (the face is the seat of identity).
\item The absence of hope (forever).
\end{itemize}
\end{aretenir}

\courssub{Chapter 4: Apparent Progress and the Final Betrayal of the Self}

Winston appears to make progress. He writes on a slate, practices doublethink, and even dreams of the Golden Country. But in a moment of semi‑consciousness, he cries out \emph{“Julia! Julia! Julia, my love!”} — revealing that his \emph{heart} still belongs to her. O’Brien declares him “cured” intellectually but not yet “whole”. The final test is Room 101.

Before Room 101, Winston clings to the idea that he is the \emph{“last man in Europe”} — the solitary guardian of truth. This is an illusion. O’Brien shows him that the Party has already won: the proles will never rebel, the Brotherhood is a myth, and Julia has betrayed him (or so he believes). The chapter ends with Winston’s realisation that he still \emph{hates} Big Brother. He has not yet reached acceptance.

\begin{attention}[The “Last Man” Illusion]
Do not mistake Winston’s momentary defiance for heroism. Orwell presents it as a \emph{delusion} — the Party allows it only to make the final breaking more complete. The tragedy is that Winston’s resistance was always within the Party’s design.
\end{attention}

\courssub{Plot Analysis: Part 3 as Climax and Falling Action}

Part 3 functions as both \cle{climax} and \cle{falling action}. The climax is the extended interrogation (Chapters 1–3); the falling action is Winston’s final capitulation in Room 101 (Chapter 4–6, though only Chapters 1–4 are in this section). Every hope raised in Part 2 is reversed:

\begin{center}
\begin{tabularx}{\linewidth}{|l|X|X|}
\hline
\rowcolor{popblueL}
\textbf{Part 2 Hope} & \textbf{Part 3 Reality} \\
\hline
Love as rebellion & Love used as a lever for betrayal \\
Privacy in the room & Total surveillance in the cell \\
The paperweight (beauty, past) & Smashed — past destroyed \\
Julia as partner & Julia as victim / betrayer \\
O’Brien as ally & O’Brien as torturer \\
\hline
\end{tabularx}
\end{center}

The love affair is not merely ended; it is \emph{weaponised}. Winston’s betrayal of Julia in Room 101 (anticipated here) is the ultimate proof that the Party can destroy the deepest human bonds.

\courssub{Situation Analysis: The Inversion of the Prison Narrative}

Traditional prison narratives (e.g., \emph{The Count of Monte Cristo}, \emph{The Shawshank Redemption}) offer \emph{escape}, \emph{martyrdom}, or \emph{spiritual triumph}. \emph{Nineteen Eighty‑Four} inverts all three:

\begin{itemize}
\item \textbf{No escape:} The Ministry of Love is inescapable; even death is not a release because the Party controls the mind before the body dies.
\item \textbf{No martyrdom:} Winston does not die for truth; he is forced to \emph{deny} truth and \emph{love} the lie.
\item \textbf{No spiritual triumph:} The self is not preserved; it is \emph{destroyed and rebuilt}.
\end{itemize}

Orwell’s warning is that totalitarianism seeks not just obedience but \emph{total conversion}. The situation is hopeless by design — the novel is a \cle{negative utopia} (dystopia) that shows the logical endpoint of unchecked state power.

\begin{methode}[Analysing a Dystopian Situation]
To analyse a dystopian situation in literature, follow these steps:
\begin{enumerate}
\item Identify the \textbf{normal world destroyed} (what freedoms existed before?).
\item Identify the \textbf{mechanism of control} (surveillance, propaganda, torture, language manipulation).
\item Examine the \textbf{victim’s response} (resistance, compliance, breakdown, conversion).
\item Articulate the \textbf{author’s warning} (what does the text say about real‑world tendencies?).
\end{enumerate}
\end{methode}

\courssub{Linking Parts 2 and 3: Systematic Destruction of Symbols}

The transition from Part 2 to Part 3 is not merely chronological; it is \emph{thematic erasure}. Each symbol of Part 2 is revisited in Part 3 only to be destroyed or perverted:

\begin{popfigure}
\begin{tikzpicture}[scale=0.9, every node/.style={font=\small, align=center}]
% Flowchart of three stages
\node[draw, fill=popblueL, rounded corners, minimum width=3.5cm, minimum height=1cm, align=center] (learn) at (0,3){Stage 1: Learning\\ (Intellectual confession)};
\node[draw, fill=popgreenL, rounded corners, minimum width=3.5cm, minimum height=1cm, align=center] (understand) at (0,1){Stage 2: Understanding\\ (Grasping Party logic)};
\node[draw, fill=poporange, rounded corners, minimum width=3.5cm, minimum height=1cm, align=center] (accept) at (0,-1){Stage 3: Acceptance\\ (Love for Big Brother)};

\draw[->, thick, popdark] (learn.south) -- (understand.north);
\draw[->, thick, popdark] (understand.south) -- (accept.north);

% Side annotations
\node[text width=4cm, align=left] at (4,3) {Winston writes “2+2=5”\\ on the slate.};
\node[text width=4cm, align=left] at (4,1) {Winston sees the\\ “boot stamping\\ on a human face”.};
\node[text width=4cm, align=left] at (4,-1) {Winston betrays Julia\\ in Room 101;\\ loves Big Brother.};
\end{tikzpicture}
\legende{Flowchart of O’Brien’s three stages of reintegration.}
\end{popfigure}

\begin{itemize}
\item \textbf{The paperweight} — bought as a link to the past; smashed by the Thought Police.
\item \textbf{The room} — a sanctuary; revealed as a trap (telescreen behind picture).
\item \textbf{The diary} — Winston’s private truth; used as evidence against him.
\item \textbf{The nursery rhyme} — “Oranges and lemons”; completed by O’Brien as a mockery of shared culture.
\item \textbf{Julia} — lover and fellow rebel; becomes the instrument of Winston’s final betrayal.
\end{itemize}

This systematic destruction reinforces the novel’s central thesis: \emph{under totalitarianism, the private sphere cannot survive}. The Party does not tolerate a parallel reality; it must invade, colonise, and replace it.

\begin{savaistu}[Orwell’s Real‑World Inspiration]
Orwell drew on accounts of Soviet show trials (the Moscow Trials of the 1930s) and the NKVD’s interrogation methods. The “confessions” of Old Bolsheviks like Zinoviev and Kamenev — broken men admitting to absurd crimes — directly informed the portrayal of Winston’s reintegration. The phrase “boot stamping on a human face” echoes the brutality of totalitarian regimes Orwell witnessed in Spain and studied in Russia.
\end{savaistu}

\courssub{Key Quotation and Examination Focus}

\begin{definition}[Key Quotation]
\emph{“Power is in tearing human minds to pieces and putting them together again in new shapes of your own choosing.”} — O’Brien, Part 3, Chapter 2.
\end{definition}

This quotation is the \cle{thesis statement} of the novel’s political philosophy. In an exam, you may be asked to:
\begin{itemize}
\item Explain its meaning in context.
\item Discuss whether the novel offers any alternative to this vision.
\item Compare it with other dystopian visions (e.g., Huxley’s \emph{Brave New World}).
\end{itemize}

\begin{attention}[O’Brien is Not a Simple Villain]
A common mistake is to treat O’Brien as a sadistic individual. He is \emph{the voice of the Party’s logic}. He does not torture for pleasure; he tortures because the Party’s power \emph{requires} the destruction of independent thought. Analyse him as an \textbf{ideological instrument}, not a character with personal motives.
\end{attention}

\begin{exercice}[Exam‑Style Question]
\textbf{“Part 3 of \emph{Nineteen Eighty‑Four} offers no hope whatsoever.”} How far do you agree with this view? Support your answer with close reference to Chapters 1–4.
\end{exercice}

\begin{corrige}[Exercise 1]
A balanced answer should:
\begin{enumerate}
\item Acknowledge the overwhelming bleakness: the destruction of Winston’s mind, the boot‑stamping image, the absence of external resistance.
\item Note the \emph{brief} moments that could be read as hope: Winston’s dream of the Golden Country, his temporary refusal to betray Julia, the very fact that he \emph{remembers} his hatred (Chapter 4) — proof that the human spirit resists until the very end.
\item Conclude that Orwell’s intention is \emph{warning}, not prophecy; the horror is meant to galvanise readers to defend truth and freedom \emph{now}. The “hope” lies in the reader’s response, not in the novel’s ending.
\end{enumerate}
\end{corrige}

\begin{exoresolu}[Worked Example: Analysing the Paperweight’s Destruction]
\textbf{Question:} How does the smashing of the paperweight symbolise the transition from Part 2 to Part 3?
\textbf{Solution:}
\begin{enumerate}
\item \textbf{Object:} The paperweight contains a piece of coral — a fragment of the past, beautiful, fragile, enclosed in glass.
\item \textbf{Part 2 meaning:} It represents Winston’s desire to preserve a private history, a truth untouched by the Party.
\item \textbf{Part 3 event:} The Thought Police smash it; the coral rolls across the mat, insignificant.
\item \textbf{Symbolic shift:} The past cannot be preserved; it is either destroyed or rewritten. The private sphere (the glass) is shattered by state power. The coral’s insignificance mirrors Winston’s own fate — a small, beautiful thing crushed by a vast machine.
\end{enumerate}
\end{exoresolu}

\courssub{Summary: The Architecture of Part 3}

\begin{aretenir}[Part 3 at a Glance]
\begin{itemize}
\item \textbf{Setting:} Ministry of Love — total surveillance, no privacy.
\item \textbf{Antagonist:} O’Brien — embodiment of Party logic.
\item \textbf{Process:} Three‑stage reintegration (learning → understanding → acceptance).
\item \textbf{Climax:} Room 101 (anticipated in Chapter 4) — betrayal of Julia.
\item \textbf{Theme:} Power as an end in itself; reality control; destruction of the self.
\item \textbf{Warning:} Totalitarianism seeks not compliance but \emph{conversion}.
\end{itemize}
\end{aretenir}

\begin{exercice}[Consolidation Exercise]
Write a paragraph analysing how Orwell uses the character of Parsons in Chapter 1 to illustrate the Party’s penetration of the family unit.
\end{exercice}

\begin{corrige}[Exercise 2]
Parsons is a loyal, unthinking Party member. His arrest by his own daughter — for saying “Down with Big Brother” in his sleep — shows that:
\begin{itemize}
\item The family has been turned into a surveillance cell.
\item Children are indoctrinated to value Party loyalty over parental love.
\item Even the most orthodox citizens are not safe; the Party demands \emph{active} enthusiasm, not just passive obedience.
\item Parsons’ pride in his daughter (“She listened at the keyhole. Heard what I was saying…”) demonstrates the completeness of the psychological transformation.
\end{itemize}
\end{corrige}

\courssub{Exam Tip: The Question of Hope}

\begin{aretenir}[Exam Tip — Discussing Hope in Part 3]
Examiners often ask whether \emph{Nineteen Eighty‑Four} is entirely hopeless. A high‑scoring response will:
\begin{itemize}
\item Distinguish between \textbf{in‑text hope} (Winston’s fate) and \textbf{reader‑level hope} (the novel as warning).
\item Cite Winston’s final thought (“He loved Big Brother”) as the death of hope \emph{within} the narrative.
\item Argue that the \emph{existence of the novel itself} — written, published, read — is an act of hope; Orwell believed that exposing the mechanism could prevent its full realisation.
\item Avoid simplistic “yes/no” answers; use textual evidence (the boot image, the paperweight, the fingers scene) to support a nuanced position.
\end{itemize}
\end{aretenir}

This concludes Section 2. In Section 3 we will turn to Chaucer’s \emph{General Prologue} and examine the portraits of the skilled workers, traders and professionals — the Cook, the Manciple, and the Shipman — and their place in the social tapestry of the pilgrimage.

\coursec{Section 3: The General Prologue I — Skilled Workers, Traders and Professionals}

In Section 2 we watched Orwell dismantle a human being: in Part 3 of \emph{Nineteen Eighty-Four}, Winston's private self is interrogated, broken and rebuilt in the image of the Party. Orwell's tool is terror; his target is the inner life. We now turn to a writer who works with the opposite instrument. Geoffrey Chaucer, writing some five and a half centuries earlier, also exposes the gap between what people appear to be and what they truly are — but his weapon is not the torture chamber. It is \emph{laughter}. Where Orwell shows us a state that crushes individuals, Chaucer shows us a society of individuals who deceive themselves and one another, and he does it with such warmth that we laugh even as we judge. In this section we study four portraits from the \emph{General Prologue} to \emph{The Canterbury Tales} — the Cook, the Manciple, the Shipman and the Doctor — and we learn the formula by which Chaucer builds every portrait: appearance, profession, and then the quiet moral twist.

\courssub{Recap: The Frame Narrative}

Before examining the portraits, let us recall the framework that holds them. Chaucer begins the \emph{General Prologue} by telling us that one April, he joined a company of pilgrims at the \cle{Tabard Inn} in Southwark, just south of London Bridge. They were travelling to Canterbury to visit the shrine of Saint Thomas Becket, the martyr-archbishop murdered in his own cathedral in 1170. There were, Chaucer says, ``wel nyne and twenty'' pilgrims — about thirty in all — drawn from nearly every level of English society: knight and plowman, nun and miller, lawyer and cook.

At the inn, the Host, \cle{Harry Bailey}, proposes a game: each pilgrim shall tell two stories on the way to Canterbury and two on the way back, and whoever tells the best tale (``best sentence and moost solaas'') will win a free supper at the Tabard, paid for by the others. This \cle{storytelling contest} is the frame that justifies the whole collection.

\begin{aretenir}[The Frame Narrative]
\begin{itemize}
\item \textbf{Setting:} the Tabard Inn, Southwark; a pilgrimage to Canterbury in springtime (April).
\item \textbf{Company:} about thirty pilgrims from all social classes — a cross-section of medieval England.
\item \textbf{Device:} the Host proposes a storytelling contest (two tales each way), with a supper as prize.
\item \textbf{Importance:} the frame lets Chaucer place all of society side by side and judge each member through a portrait.
\end{itemize}
\end{aretenir}

\begin{attention}[The Narrator Is Not Chaucer]
The ``Chaucer'' who speaks in the \emph{General Prologue} is a \emph{character} — a naive, easily impressed pilgrim who admires almost everyone he meets. When the narrator praises a pilgrim, the real Chaucer, the author, may be laughing behind his back. Never confuse the narrator's voice with the author's judgement. This distinction is the key to reading Chaucer's irony.
\end{attention}

\courssub{Satire and Irony: Chaucer's Tools}

\begin{definition}[Satire]
\cle{Satire} is the use of humour, irony, exaggeration or ridicule to expose and criticise human vice or folly, often with the aim of correcting it. Chaucer's satire is usually \emph{gentle} (we call it Horatian, after the Roman poet Horace): he smiles at human weakness rather than raging against it.
\end{definition}

\begin{definition}[Irony]
\cle{Irony} is saying one thing while meaning another, so that the surface meaning and the intended meaning are opposite. In the \emph{General Prologue}, Chaucer's praise is often disguised blame: when the narrator calls a pilgrim the ``best'' at his trade, we must ask \emph{best at what, exactly} — and at whose expense.
\end{definition}

\begin{propriete}[The Moral Twist — Chaucer's Portrait Formula (Examinable)]
Chaucer builds almost every portrait in three movements:
\begin{enumerate}
\item \textbf{Appearance:} a vivid physical or material detail (clothing, face, tools, horse).
\item \textbf{Profession:} enthusiastic praise of the pilgrim's skill — often exaggerated.
\item \textbf{Moral twist:} a quiet, almost casual detail that reveals corruption, greed or hypocrisy.
\end{enumerate}
The twist is usually placed \emph{at the end} of the portrait, where it colours everything said before. This pattern is examinable: you should be able to identify it in any portrait and quote the twisting detail.
\end{propriete}

\begin{popfigure}
\begin{center}
\begin{tikzpicture}[font=\small]
\node[draw=popblue, fill=popblueL, rounded corners, minimum width=3.4cm, minimum height=1.1cm, align=center] (a) at (0,0) {\textbf{1. Appearance}\\physical detail};
\node[draw=poporange, fill=popgold!25, rounded corners, minimum width=3.4cm, minimum height=1.1cm, align=center] (b) at (4.6,0) {\textbf{2. Profession}\\praise of skill};
\node[draw=poppurple, fill=poppink!30, rounded corners, minimum width=3.4cm, minimum height=1.1cm, align=center] (c) at (9.2,0) {\textbf{3. Moral twist}\\hidden corruption};
\draw[-{latex}, very thick, popdark] (a) -- (b);
\draw[-{latex}, very thick, popdark] (b) -- (c);
\node[align=center, text=popmuted, font=\footnotesize] at (4.6,-1.2) {The reader admires first, then judges: the twist re-reads the praise as irony.};
\end{tikzpicture}
\legende{Chaucer's portrait formula: appearance, profession, moral twist.}
\end{center}
\end{popfigure}

\courssub{The Cook: Skill with an Ulcer}

Chaucer introduces the Cook, \cle{Roger of Ware} (Roger ``Hodge'' of Ware), with warm praise. He is a master of his trade: he can boil, roast, broil and fry, make thick soups (\cle{mortreux}), and bake pies. The narrator especially admires his \cle{blancmange}, a creamy white dish of chicken or fish with rice and almond milk — a luxury dish for the wealthy. Then comes the twist, slipped in with terrible casualness: on his shin, the Cook has a \cle{mormal} — an open, running ulcer or sore.

\begin{exemplebox}[The Cook's Portrait — The Formula at Work]
\begin{itemize}
\item \textbf{Appearance:} an ulcer (mormal) on his shin.
\item \textbf{Profession:} master of all cookery techniques; makes a fine blancmange; knows his London ale.
\item \textbf{Moral twist:} the running sore. The reader suddenly imagines that sore near the creamy white blancmange — and the dish is spoiled forever in our minds.
\end{itemize}
\end{exemplebox}

The effect is \cle{comic realism}. Chaucer does not sermonise about hygiene; he simply places two facts side by side — delicious white sauce, oozing sore — and lets our stomachs do the judging. There is also a hint of dishonest trade: later in the \emph{Cook's Prologue}, the Host jokes that the Cook's shop sells reheated, stale pasties and that his pies have made many a pilgrim sick. The satirical target is the \emph{unhygienic and dishonest food-seller}, a figure any market-goer in Douala or Bamenda will recognise instantly.

\courssub{The Manciple: Cunning Without Learning}

A \cle{manciple} is a steward — the officer who buys provisions for an institution such as a college, an inn of court, or a monastery. Chaucer's Manciple works for a \cle{temple} (one of the Inns of Court in London, where lawyers lived and studied), purchasing food for more than thirty learned lawyers, men expert in law and finance. The Manciple himself is \emph{unlearned} — he has no formal education. Yet Chaucer tells us, with a perfectly straight face, that this illiterate steward is so shrewd at buying and accounting that he cheats or outwits all of them, and lives comfortably on the profit.

\begin{exemplebox}[The Manciple's Portrait — The Formula at Work]
\begin{itemize}
\item \textbf{Appearance:} little physical description — the portrait is almost entirely about his cleverness.
\item \textbf{Profession:} a master buyer; ``wise in buying of victuals''; whether he pays cash or credit, he always comes out ahead.
\item \textbf{Moral twist:} he is ``alway... afore'' his masters — he skims money from men paid precisely to detect such tricks.
\end{itemize}
\end{exemplebox}

The satire here is double-edged. It mocks the Manciple's \emph{dishonesty}, yes — but it mocks the \emph{learned lawyers} even more. What is book learning worth, Chaucer asks, if thirty masters of law can be robbed blind by one unlettered steward? The portrait celebrates \cle{practical cunning} over \cle{book learning}, a theme that still resonates wherever ``street sense'' outsmarts certificates.

\courssub{The Shipman: The Morally Ambiguous Professional}

The Shipman (or Skipper) comes from the West Country — ``from \cle{Dartmouth},'' Chaucer guesses — and rides his horse awkwardly, ``as he could,'' because he is a man of the sea. He wears a woollen gown to the knee and carries a dagger on a cord round his neck. As a sailor, he is superb: he knows every harbour, every current and every tide from \cle{Gotland} (in the Baltic) to the coast of \cle{Spain} (``Cartage''); no one calculates a voyage better. His ship is called the \emph{Maudelayne} (Magdalene).

Then the twists — and there are two. First, on the voyage home, he steals wine from the merchant's cargo he is paid to transport, ``by night,'' while the merchant sleeps. Second, and far darker: in sea battles, when he takes prisoners, he sends them home ``by water'' — that is, he makes them \cle{walk the plank}; he drowns them.

\begin{exemplebox}[The Shipman's Portrait — The Formula at Work]
\begin{itemize}
\item \textbf{Appearance:} woollen gown to the knee, dagger on a lace about his neck, sun-browned face, rides awkwardly.
\item \textbf{Profession:} unmatched knowledge of tides, currents and harbours; ``ther nas noon swich from Hulle to Cartage.''
\item \textbf{Moral twist:} he steals the wine he transports and drowns his prisoners.
\end{itemize}
\end{exemplebox}

The Shipman is the \emph{morally ambiguous professional}: genuinely excellent at his work, genuinely ruthless outside it. Chaucer neither condemns him outright nor excuses him. The narrator's breezy tone — as if drowning prisoners were a minor habit — is itself the irony: the gap between the cheerful voice and the brutal fact forces the reader to become the judge.

\courssub{The Doctor: Medicine and the Love of Gold}

The Doctor of Physic (the Physician) receives one of the longest and most brilliant portraits. He is a perfect master of medieval medicine: he knows \cle{astronomy} (astrology, used to choose lucky hours for treatment, bloodletting, and administering drugs) and the theory of the four \cle{humours} — blood, phlegm, yellow bile (choler) and black bile (melancholy) — whose balance was believed to govern health. He can diagnose any illness and trace it to its humoural cause. The narrator calls him, in the famous ironic formula, a ``verray, parfit praktisour'' — a true, perfect practitioner.

Now watch the twists accumulate. His \emph{diet} is moderate — he eats little, we are told, though this may be thrift rather than virtue. He has read little of the Bible — a strange gap in a learned man of the Church's world. He works in profitable partnership with \cle{apothecaries} (druggists): each sends business to the other, and both profit (a medieval kickback scheme). He dressed in rich red and blue silk lined with taffeta — conspicuously wealthy attire. And then the final, devastating line of the portrait: he kept the gold he earned during the plague (when many doctors fled), because — Chaucer concludes — ``\cle{gold in physic is a cordial}'' (gold is good for the heart, a standard medicinal belief), and therefore ``he lovede gold in special.''

\begin{exemplebox}[The Doctor's Portrait — The Formula at Work]
\begin{itemize}
\item \textbf{Appearance:} dressed in blood-red and blue silk, lined with taffeta.
\item \textbf{Profession:} grounded in astronomy and the four humours; a ``verray, parfit praktisour.''
\item \textbf{Moral twist:} a profitable arrangement with apothecaries, little Bible reading, and above all: ``he lovede gold in special.''
\end{itemize}
\end{exemplebox}

\begin{attention}[Do Not Trust the Praise]
``Verray, parfit praktisour'' sounds like a compliment. But by the end of the portrait we know the Doctor is greedy, irreligious and profiteering. The praise is \emph{ironic}: perfect at the \emph{science} of medicine, corrupted in its \emph{practice}. In the examination, never quote Chaucer's praise without asking whether it is sincere.
\end{attention}

\begin{popfigure}
\begin{center}
\begin{tikzpicture}[font=\footnotesize]
\draw[thick, popdark] (0,0) grid (12,5);
\foreach \y in {1,2,3,4} \draw[thick, popdark] (0,\y) -- (12,\y);
\foreach \x in {3,6,9} \draw[thick, popdark] (\x,0) -- (\x,5);
\node[fill=popblueL, minimum width=3cm, minimum height=1cm] at (1.5,4.5) {\textbf{Pilgrim}};
\node[fill=popblueL, minimum width=3cm, minimum height=1cm] at (4.5,4.5) {\textbf{Skill}};
\node[fill=popblueL, minimum width=3cm, minimum height=1cm] at (7.5,4.5) {\textbf{Flaw}};
\node[fill=popblueL, minimum width=3cm, minimum height=1cm] at (10.5,4.5) {\textbf{Satirical target}};
\node[align=center] at (1.5,3.5) {Cook};
\node[align=center] at (4.5,3.5) {Fine blancmange,\\all cookery};
\node[align=center] at (7.5,3.5) {Ulcer on his shin;\\stale pasties};
\node[align=center] at (10.5,3.5) {Unhygienic,\\dishonest food trade};
\node[align=center] at (1.5,2.5) {Manciple};
\node[align=center] at (4.5,2.5) {Shrewd buying\\of victuals};
\node[align=center] at (7.5,2.5) {Cheats 30\\learned lawyers};
\node[align=center] at (10.5,2.5) {Book learning beaten\\by cunning};
\node[align=center] at (1.5,1.5) {Shipman};
\node[align=center] at (4.5,1.5) {Master of tides,\\harbours, navigation};
\node[align=center] at (7.5,1.5) {Steals wine; drowns\\prisoners};
\node[align=center] at (10.5,1.5) {The ruthless\\professional};
\node[align=center] at (1.5,0.5) {Doctor};
\node[align=center] at (4.5,0.5) {Astronomy, humours,\\perfect diagnosis};
\node[align=center] at (7.5,0.5) {Loves gold; deals\\with apothecaries};
\node[align=center] at (10.5,0.5) {Greedy, profit-driven\\medicine};
\end{tikzpicture}
\legende{Portrait grid: skill, flaw and satirical target for each pilgrim.}
\end{center}
\end{popfigure}

\courssub{The Social Situation: A Rising Middle Class}

Why do these four pilgrims matter together? Notice what they are: a cook, a steward, a merchant sailor, a physician. None is an aristocrat; none is a clergyman. They are \cle{skilled workers, traders and professionals} — the rising \cle{middle class} of fourteenth-century England. After the Black Death (1348--1350), labour became scarce and valuable; trade grew; men could rise by skill and money rather than birth. Chaucer's portraits register this new social energy — and its new temptations. These men are good at what they do, and several of them use that skill to cheat.

\begin{savaistu}[Did You Know? — A Cameroonian Mirror]
Chaucer's gallery is not far from home. Think of the skilled cook in a Yaound\'e \emph{chop house} whose hygiene you would rather not inspect; the clever purchasing officer who grows rich on small margins nobody checks; the transporter who is brilliant on the Douala--Bafoussam road but not above ``diverting'' part of the cargo; the well-paid specialist whose first question is the fee. Chaucer's satire endures because the types endure. In class discussion, ask: which modern Cameroonian professions would a twenty-first-century Chaucer portrait, and where would he place the moral twist?
\end{savaistu}

\begin{methode}[How to Master Each Portrait]
For every portrait in the \emph{General Prologue}, memorise four things:
\begin{enumerate}
\item \textbf{One physical detail} (the Cook's ulcer; the Doctor's red and blue silk).
\item \textbf{One professional detail} (the Shipman's knowledge of tides; the Manciple's shrewd buying).
\item \textbf{One moral flaw} (the Doctor's love of gold; the Shipman's drowned prisoners).
\item \textbf{One short quotation} (``he lovede gold in special''; ``verray, parfit praktisour'').
\end{enumerate}
Then, in any essay, present them in that order — appearance, profession, twist — and you will reproduce Chaucer's own method in your analysis.
\end{methode}

\begin{exoresolu}[Analysing the Moral Twist]
\textbf{Statement.} A candidate writes: ``Chaucer greatly admires the Doctor, calling him a perfect practitioner, which shows that the Doctor is the most honest pilgrim in the General Prologue.'' Identify the error in this reading and rewrite the analysis correctly.
\tcblower
\textbf{Solution.} The error is to take the narrator's praise at face value. The phrase ``verray, parfit praktisour'' is \emph{ironic}: it praises the Doctor's \emph{technical} mastery while the rest of the portrait dismantles his \emph{moral} character. Chaucer shows us a physician who reads little Scripture, who has a mutually profitable arrangement with apothecaries, who grew rich during the plague, and whose portrait closes on the line ``he lovede gold in special.'' The correct reading follows Chaucer's formula: appearance (rich silk clothing) leads to profession (perfect knowledge of astronomy and the humours), which is then undercut by the moral twist (greed). A better sentence would be: ``Chaucer's apparently admiring portrait of the Doctor is deeply ironic: the `perfect practitioner' of medicine is a perfect practitioner of profit, whose learning serves his love of gold rather than his patients.'' The lesson: in Chaucer, always ask \emph{who is speaking} and \emph{whether the praise survives the final line of the portrait}.
\end{exoresolu}

\begin{exercice}[Portrait Analysis]
\begin{enumerate}
\item For the Shipman, identify one physical detail, one professional detail and one moral flaw, and explain how the narrator's cheerful tone creates irony.
\item Explain the double satire in the portrait of the Manciple: who is mocked more, the steward or the lawyers, and why?
\item In a short paragraph, show how the detail of the Cook's ulcer transforms the earlier praise of his blancmange. What is this technique called?
\item Essay plan: ``Chaucer's skilled professionals are all admirable and all corrupt.'' Outline three body paragraphs using the portrait formula (appearance, profession, moral twist).
\end{enumerate}
\end{exercice}

\begin{corrige}[Exercise 1]
\begin{enumerate}
\item \textbf{Physical:} woollen gown to the knee, dagger on a cord, sun-browned face. \textbf{Professional:} unequalled knowledge of tides, currents and harbours. \textbf{Flaw:} he steals the wine he transports and drowns his prisoners. The breezy, admiring tone never changes when these crimes are mentioned, so the reader feels the gap between tone and fact — that gap \emph{is} the irony.
\item The lawyers are mocked more. The Manciple's dishonesty is expected satire of a cheating servant; but thirty men trained in law and finance, outwitted by an uneducated steward, make a sharper joke — book learning is shown to be helpless before practical cunning.
\item The praise of the blancmange comes first, so the reader savours the dish; then the ulcer appears and retrospectively contaminates it — we re-read the praise with disgust. This retrospective undercutting is the \emph{moral twist}, a form of structural irony.
\item Suggested plan: \textbf{P1} — the Doctor: perfect science corrupted by gold (quote ``he lovede gold in special''). \textbf{P2} — the Shipman: perfect seamanship joined to theft and murder. \textbf{P3} — the Manciple and Cook: cunning and skill serving dishonesty and filth. Conclusion: Chaucer admires skill genuinely but shows that in the rising middle class, skill without morality becomes a refined form of corruption.
\end{enumerate}
\end{corrige}

\begin{aretenir}[Key Points of the Section]
\begin{itemize}
\item The frame: the Tabard Inn, about thirty pilgrims, Harry Bailey's storytelling contest (two tales each way).
\item Chaucer's formula: \textbf{appearance} $\rightarrow$ \textbf{profession} $\rightarrow$ \textbf{moral twist}.
\item The Cook: fine blancmange, ulcer on his shin — comic realism, unhygienic trade.
\item The Manciple: unlearned steward who outwits thirty lawyers — cunning over book learning.
\item The Shipman: master sailor who steals wine and drowns prisoners — the morally ambiguous professional.
\item The Doctor: ``verray, parfit praktisour'' who ``lovede gold in special'' — satire of greedy medicine.
\item The narrator's praise is often ironic; the final line of a portrait usually carries the true judgement.
\item These pilgrims embody the rising middle class of fourteenth-century England — skilled, energetic, and morally compromised.
\end{itemize}
\end{aretenir}

\coursec{Section 4: The General Prologue II — Women, Ideal Christians and Corrupt Churchmen}

In Section 3 we examined the skilled workers and traders — the Cook, the Manciple and the Shipman — whose portraits reveal Chaucer’s keen eye for professional detail and social satire. We now turn to a group of pilgrims who embody the moral and spiritual spectrum of medieval society: the Physician, the Wife of Bath, the Parson, the Plowman, the Summoner, the Pardoner and the Host. These portraits are arranged with deliberate artistry: the ideal churchman and the ideal layman stand in stark contrast to the corrupt church officials, while the Wife of Bath and the Host occupy a complex middle ground. Understanding this structural contrast is essential for any essay on Chaucer’s satire, his view of the Church, or his portrayal of women.

\courssub{The Physician: Science, Greed and the Limits of Medieval Medicine}

Before the Wife of Bath, Chaucer places the Physician — a figure who bridges the world of learning and the world of profit.

\begin{experience}[Meeting the Physician]
Read the following lines from the General Prologue (ll. 413–444) and note every detail that builds his character:
\begin{quote}
With us ther was a doctour of phisik; \\
In al this world ne was ther noon hym lik, \\
To speke of phisik and of surgerye; \\
For he was grounded in astronomye. \\
He kepte his pacient a ful greet deel \\
In houres by his magik natureel. \\
Wel coude he fortunen the ascendent \\
Of his ymages for his pacient. \\
He knew the cause of everich maladye, \\
Were it of hoot, or coold, or moyste, or drye. \\
He knew the humour of everich man, \\
And after that he gan his medecyne \\
To yive hym, for the cause and the ground. \\
Reddy was his apothecaries \\
To senden hym drogges and his letuaries; \\
For ech of hem made oother for to wynne, \\
Hir frendshipe nas nat newe to bigynne. \\
Wel knew he the olde Esculapius, \\
And Deyscorides, and eek Rufus, \\
Ypocras, Galien, and Serapion, \\
Razis, and Avicen, and Averrois, \\
Damascien, and Constantyn, and Bernard, \\
Gatesden, and Gilbertyn. \\
His diete was mesurable and greet, \\
For ther nas no superfluitee. \\
Of his studie took he moost cure and heed; \\
Nat a word spak he moore than was neede, \\
And that was seyd in forme and reverence, \\
And short and quyck, and ful of hy sentence. \\
Sowyng in moral vertu was his speche, \\
And gladly wolde he lerne, and gladly teche. \\
Of gold he loved best, in special, \\
To make hym a ful greet chevisaunce; \\
For gold in phisik is a cordial, \\
Therfore he lovede gold in special.
\end{quote}
\end{experience}

\begin{popfigure}
\begin{tikzpicture}[scale=0.9, every node/.style={font=\small}]
% Central node
\node[draw, rounded corners, fill=popblueL, text width=3.2cm, align=center, minimum height=1.8cm] (phys) at (0,0) {\textbf{The Physician}\\Doctor of Physik};
% Knowledge
\node[draw, rounded corners, fill=popgreen!20, text width=2.5cm, align=center] (know) at (-5,2.5) {Knowledge:\\Astronomy, humours, authorities (Galen, Avicenna, etc.), surgery};
% Practice
\node[draw, rounded corners, fill=poporange!20, text width=2.5cm, align=center] (prac) at (5,2.5) {Practice:\\Horoscope-based diagnosis, apothecary kickbacks, moderate diet};
% Character
\node[draw, rounded corners, fill=poppink!20, text width=2.5cm, align=center] (char) at (-5,-2.5) {Character:\\Taciturn, sententious, morally upright in speech};
% Motivation
\node[draw, rounded corners, fill=popgold!20, text width=2.5cm, align=center] (mot) at (5,-2.5) {Motivation:\\'Of gold he loved best' — profit-driven, 'gold in phisik is a cordial'};
% Arrows
\draw[->, thick, popdark] (phys) -- (know);
\draw[->, thick, popdark] (phys) -- (prac);
\draw[->, thick, popdark] (phys) -- (char);
\draw[->, thick, popdark] (phys) -- (mot);
\end{tikzpicture}
\legende{Character web of the Physician: knowledge, practice, character and motivation.}
\end{popfigure}

\begin{aretenir}[The Physician — Key Points]
\begin{itemize}
\item \textbf{Professional authority:} 'In al this world ne was ther noon hym lik / To speke of phisik and of surgerye' — he is acknowledged as the best.
\item \textbf{Intellectual basis:} Grounded in \textbf{astronomy} (astrological medicine), the theory of the \textbf{four humours} (hot, cold, moist, dry), and a long catalogue of authorities from antiquity (Esculapius, Hippocrates/Ypocras, Galen) through Arabic scholars (Avicenna/Avicen, Averroes/Averrois, Razis) to medieval English doctors (Gatesden, Gilbertyn).
\item \textbf{Method:} Uses horoscopes ('fortunen the ascendent / Of his ymages') to time treatments; maintains a mutually profitable arrangement with apothecaries ('For ech of hem made oother for to wynne').
\item \textbf{Personal habits:} Moderate diet, no superfluity, taciturn ('Nat a word spak he moore than was neede'), speech full of 'hy sentence' (high moral wisdom).
\item \textbf{The ironic twist:} Despite his learning and apparent virtue, the final couplet exposes his ruling passion: \textbf{'Of gold he loved best... For gold in phisik is a cordial, / Therfore he lovede gold in special.'} The 'cordial' (heart medicine) becomes a pun: gold is the only medicine he truly values.
\item \textbf{Chaucer's tone:} Respect for his learning, but savage irony about his greed. He is a \textbf{competent professional corrupted by avarice}.
\end{itemize}
\end{aretenir}

\begin{attention}[The Physician vs. the Pardoner — Different Kinds of Greed]
Do not confuse the Physician's greed with the Pardoner's. The Physician is a \textbf{genuine expert} who charges excessive fees and colludes with apothecaries; his greed is professional corruption. The Pardoner is a \textbf{fraud} selling fake spiritual benefits; his greed is sacramental betrayal. In essays, distinguish: the Physician betrays his \emph{professional ethics}; the Pardoner betrays his \emph{sacred office}.
\end{attention}

\courssub{The Wife of Bath: A Complex Portrait of Female Experience}

\begin{experience}[Meeting the Wife of Bath]
Read the following lines from the General Prologue (ll. 445–476) and note every detail that builds her character:
\begin{quote}
A good wif was ther of biside Bathe, \\
But she was somdel deef, and that was scathe. \\
Of clooth-makyng she hadde swich an haunt, \\
She passed hem of Ypres and of Gaunt. \\
In al the parisshe wif ne was ther noon \\
That to the offrynge bifore hire sholde goon; \\
And if ther dide, certeyn so wrooth was she, \\
That she was out of alle charitee. \\
Hir coverchiefs ful fyne weren of ground; \\
I dorste swere they weyeden ten pound \\
That on a Sonday weren upon hir heed. \\
Hir hosen weren of fyn scarlet reed, \\
Ful streite yteyd, and shoes ful moyste and newe. \\
Boold was hir face, and fair, and reed of hewe. \\
She was a worthy womman al hir lyve: \\
Housbondes at chirche dore she hadde fyve, \\
Withouten oother compaignye in youthe, — \\
But therof nedeth nat to speke as nowthe. \\
And thries hadde she been at Jerusalem; \\
She hadde passed many a straunge strem; \\
At Rome she hadde been, and at Boloigne, \\
In Galice at Seint-Jame, and at Coloigne. \\
She coude muchel of wandrynge by the weye. \\
Gat-tothed was she, soothly for to seye. \\
Upon an amblere esily she sat, \\
Ywympled wel, and on hir heed an hat \\
As brood as is a bokeler or a targe; \\
A foot-mantel aboute hir hipes large, \\
And on hir feet a paire of spores sharpe. \\
In felaweshipe wel coude she laughe and carpe. \\
Of remedies of love she knew per chaunce, \\
For she coude of that art the olde daunce.
\end{quote}
\end{experience}

\begin{popfigure}
\begin{tikzpicture}[scale=0.9, every node/.style={font=\small}]
% Central node
\node[draw, rounded corners, fill=popblueL, text width=3.2cm, align=center, minimum height=1.8cm] (wife) at (0,0) {\textbf{The Wife of Bath}\\Alisoun of Bath};
% Appearance
\node[draw, rounded corners, fill=poppink!20, text width=2.5cm, align=center] (app) at (-5,2.5) {Appearance:\\gap-toothed, bold face, red stockings, fine kerchiefs, broad hat, spurs};
% Marital History
\node[draw, rounded corners, fill=poporange!20, text width=2.5cm, align=center] (mar) at (5,2.5) {Marital History:\\5 husbands at church door, first at age 12 (from Prologue), experienced in 'remedies of love'};
% Skills
\node[draw, rounded corners, fill=popgreen!20, text width=2.5cm, align=center] (ski) at (-5,-2.5) {Skills:\\expert cloth-maker (surpasses Ypres and Ghent), knowledgeable in love remedies};
% Travels
\node[draw, rounded corners, fill=popgold!20, text width=2.5cm, align=center] (trav) at (5,-2.5) {Travels:\\Jerusalem (3 times), Rome, Boulogne, Santiago de Compostela, Cologne};
% Personality
\node[draw, rounded corners, fill=poppurple!20, text width=2.5cm, align=center] (per) at (0,-5) {Personality:\\boisterous, deaf (from blow by 5th husband Jankyn, per her Prologue), domineering in parish, laughter and carp, authoritative};
% Arrows
\draw[->, thick, popdark] (wife) -- (app);
\draw[->, thick, popdark] (wife) -- (mar);
\draw[->, thick, popdark] (wife) -- (ski);
\draw[->, thick, popdark] (wife) -- (trav);
\draw[->, thick, popdark] (wife) -- (per);
\end{tikzpicture}
\legende{Character web of the Wife of Bath: appearance, marital history, skills, travels and personality.}
\end{popfigure}

\begin{aretenir}[The Wife of Bath — Key Points]
\begin{itemize}
\item \textbf{Name and origin:} Alisoun, from Bath (a major cloth-making centre in Somerset).
\item \textbf{Appearance:} gap-toothed (traditionally a sign of lust, travel, and boldness), bold/red face, scarlet stockings, heavy kerchiefs ('weyeden ten pound' — hyperbole for display), broad hat ('as brood as is a bokeler or a targe' — shield-sized), sharp spurs — she dresses to be seen and to ride astride.
\item \textbf{Deafness:} 'somdel deef' in GP; her Prologue reveals it was caused by a blow from her fifth husband, Jankyn — a detail that hints at domestic violence and her resilience.
\item \textbf{Marital history:} five husbands, all married 'at chirche dore' (public, legal unions); first marriage at age 12 (revealed in her Prologue); she boasts of her mastery over them.
\item \textbf{Professional skill:} a cloth-maker whose work 'passed hem of Ypres and of Gaunt' — the great Flemish centres — she is economically independent and professionally elite.
\item \textbf{Travels:} three pilgrimages to Jerusalem, plus Rome, Boulogne, Santiago de Compostela (Seint-Jame in Galice), Cologne — unusually widely travelled for a woman, suggesting wealth, piety (or its performance), and autonomy.
\item \textbf{Personality:} domineering in parish (must go first to the offering), boisterous, knowledgeable in 'remedies of love' ('the olde daunce'), laughs and jests easily ('carpe').
\item \textbf{Chaucer's tone:} admiring ('worthy womman') yet ironic; she is both a vivid individual and a type (the experienced wife).
\end{itemize}
\end{aretenir}

\begin{definition}[Frame Narrative]
A \cle{frame narrative} (or frame story) is a literary technique in which a main story (the frame) contains and introduces a series of shorter stories. In \emph{The Canterbury Tales}, the pilgrimage to Canterbury is the frame: the Host proposes a storytelling contest, and each pilgrim tells a tale. The frame provides a social context, a reason for the tales, and a structure for character interaction.
\end{definition}

\begin{attention}[Avoid One-Sided Readings of the Wife of Bath]
Many students write that Chaucer simply mocks the Wife of Bath as a lustful, domineering shrew. This is a \textbf{reductive reading}. Critics from the 20th century onward (e.g., Carolyn Dinshaw, Elaine Tuttle Hansen) argue that her portrait contains \textbf{sympathy and proto-feminist energy}: she speaks for female experience, challenges patriarchal authority (especially in her Prologue), and survives in a male-dominated world. In essays, acknowledge her flaws \emph{and} her vitality; quote 'worthy womman' and her economic independence. A balanced answer scores higher.
\end{attention}

\begin{exemplebox}[Key Quotation: The Wife of Bath's Marital Boast]
\textbf{Lines 459–460:} \\
'She was a worthy womman al hir lyve: \\
Housbondes at chirche dore she hadde fyve'

\textbf{Analysis:} 'worthy' suggests both social respect and moral worth; 'al hir lyve' implies consistency. 'At chirche dore' emphasises the public, sacramental nature of each marriage — she is not a mistress but a wife five times over. The line combines admiration with a hint of irony: is she 'worthy' because she survived five husbands, or because she mastered them?
\end{exemplebox}

\courssub{The Parson and the Plowman: The Ideal Christian Pair}

After the worldly Wife, Chaucer presents two figures who embody Christian ideals without hypocrisy. They are the only pilgrims described in wholly positive terms.

\begin{propriete}[The Parson — The Ideal Churchman]
\textbf{Portrait highlights (ll. 477–528):}
\begin{itemize}
\item 'A good man was ther of religioun' — immediate moral endorsement.
\item 'Povre' in goods, but 'rich he was of hooly thought and werk'.
\item He practices what he preaches: 'First he wroghte, and afterward he taughte.'
\item He visits his parishioners on foot, even in sickness and rain, giving his own goods to the poor.
\item His guiding principle: \textbf{'If gold ruste, what shal iren do?'} — if the priest (gold) is corrupt, what can be expected of the laity (iron)?
\item He does not seek pomp or revenue; he is a true shepherd, not a mercenary.
\item Chaucer concludes: 'I trowe ther was no bettre preest any wher.'
\end{itemize}
\end{propriete}

\begin{propriete}[The Plowman — The Ideal Layman]
\textbf{Portrait highlights (ll. 529–543):}
\begin{itemize}
\item The Parson's brother — a deliberate familial link reinforcing shared values.
\item 'A trewe swynkere and a good was he, / Lyvynge in pees and parfit charitee.'
\item He pays his tithes fully and promptly ('his tithes payde he ful faire and wel').
\item He loves his neighbour as himself, helping the poor without thought of reward.
\item He rides a mare (humble transport) and wears a tabard smocked with dirt — honest labour.
\end{itemize}
\end{propriete}

\begin{aretenir}[Key Contrast: The Parson and Plowman Prove Chaucer's Target is Corruption, Not the Church]
Chaucer places these two portraits \textbf{immediately before} the Summoner and the Pardoner. This arrangement is a structural argument: the existence of genuinely holy clergy (the Parson) and devout laypeople (the Plowman) demonstrates that Chaucer satirises \textbf{corrupt individuals}, not the institution of the Church itself. In essays, use this contrast to nuance any claim that Chaucer is anti-clerical. He respects the ideal; he exposes the failure to live up to it.
\end{aretenir}

\begin{savaistu}[Historical Context: The Parson and Wycliffite Ideas]
The Parson's principle 'If gold ruste, what shal iren do?' echoes the teachings of \textbf{John Wycliffe} (c. 1328–1384), who argued for clerical poverty and the primacy of Scripture over ecclesiastical hierarchy. Wycliffe's followers (Lollards) were persecuted as heretics. Chaucer, who knew John of Gaunt (Wycliffe's protector), may be signalling sympathy for reform — but the Parson is \emph{orthodox} in practice, not a radical. He represents the \textbf{ideal within the system}, not a rebel against it.
\end{savaistu}

\courssub{The Summoner and the Pardoner: The Corrupt Churchmen}

These two portraits form a pair of ecclesiastical corruption. They are officials of the church courts (Summoner) and the papal pardon system (Pardoner), and both abuse their offices for personal gain.

\begin{propriete}[The Summoner — The Corrupt Court Official]
\textbf{Portrait highlights (ll. 623–666):}
\begin{itemize}
\item Physical repulsiveness: 'fyred cherubynnes face', 'carbuncles', 'whelkes', 'knobby' — his outer ugliness mirrors inner corruption; 'cherubynnes' is ironic (cherubs are angelic).
\item Loves garlic, onions, leeks, and strong red wine — indulgence in sensual pleasures, contra ascetic ideal.
\item 'Questio quid juris' — he repeats a legal formula ('What is the point of law?') without understanding it, pretending learning.
\item He lets sinners off for a 'quart of wyn' — selling absolution for drink.
\item He has a network of 'fear' and 'extortion'; he knows the secrets of young girls and uses them for blackmail.
\item He wears a garland 'as it were a tavern sign' and carries a cake — he is more at home in the alehouse than the court.
\item Children are afraid of him — a detail of social terror.
\end{itemize}
\end{propriete}

\begin{propriete}[The Pardoner — The Climax of Ecclesiastical Corruption]
\textbf{Portrait highlights (ll. 667–714):}
\begin{itemize}
\item He sings 'Com hider, love, to me!' — a love song, inappropriate for a churchman, suggesting sexual ambiguity.
\item Physical description: 'yelowe' hair like wax, thin, no beard, high voice — ambiguous gender presentation, suggesting sterility or effeminacy; 'I trowe he were a geldyng or a mare.'
\item He carries a bag of fake relics: a pillowcase claimed to be the Virgin's veil ('a pilwe-beer... was Oure Lady veyl'), a piece of sail from St Peter's boat, a cross set with stones ('latoun ful of stones'), pig's bones in a glass ('pigges bones').
\item With these relics he cheats poor parsons out of more money in a day than they earn in two months.
\item His sermon theme is \textbf{avarice} ('Radix malorum est cupiditas' — 'The root of evils is greed') — yet he himself is motivated by greed: 'I preche of no thyng but for coveitise' (from his Prologue, l. 403).
\item He openly admits (in his Prologue) that he cares nothing for the correction of sin, only for profit.
\end{itemize}
\end{propriete}

\begin{savaistu}[Exam Tip: Gentle vs Savage Satire]
A frequent GCE question: \textbf{'Is Chaucer's satire gentle or savage?'}
Use the \textbf{Parson} (gentle admiration, warm humour) and the \textbf{Pardoner} (savage exposure, no redeeming trait) as your two poles. The Wife of Bath and the Host sit in the middle — affectionate irony. The Summoner is grotesque but almost farcical. Structure your essay: define the terms, analyse the Parson (gentle), the Pardoner (savage), then discuss the middle ground. Conclude that Chaucer's satire is \textbf{both}, calibrated to the moral weight of the target.
\end{savaistu}

\courssub{The Host (Harry Bailey): The Master of Ceremonies}

\begin{propriete}[The Host — Plot Device and Character]
\textbf{Portrait highlights (ll. 747–770, and throughout the links):}
\begin{itemize}
\item 'A semely man' with 'eyen stepe' and 'boold' speech — physically imposing, confident.
\item 'Faire and large' — generous, merry, a natural leader.
\item He proposes the storytelling contest: each pilgrim tells two tales on the way to Canterbury and two on the return; the best tale wins a supper at the others' cost.
\item He appoints himself judge and guide: 'And I shal be youre governour in this.'
\item He maintains order, interrupts tales (e.g., Chaucer's Tale of Sir Thopas), and mediates conflicts.
\item He is the \textbf{architect of the frame narrative} — without him, the tales would have no unifying structure.
\end{itemize}
\end{propriete}

\courssub{Structural Contrast: The Moral Spectrum of the General Prologue}

Chaucer arranges the final portraits to create a moral gradient. The following diagram visualises this spectrum.

\begin{popfigure}
\begin{tikzpicture}[scale=0.95, every node/.style={font=\small}]
% Axis
\draw[->, thick, popdark] (-7,0) -- (7,0) node[right] {Moral Corruption $\rightarrow$};
% Ideal pair
\node[draw, rounded corners, fill=popgreen!30, text width=2.2cm, align=center, minimum height=1.5cm] (parson) at (-5,1) {\textbf{The Parson}\\Ideal Churchman\\'If gold ruste...'};
\node[draw, rounded corners, fill=popgreen!30, text width=2.2cm, align=center, minimum height=1.5cm] (plow) at (-5,-1) {\textbf{The Plowman}\\Ideal Layman\\Honest labour, tithes};
% Physician (competent but greedy)
\node[draw, rounded corners, fill=popblue!30, text width=2.2cm, align=center, minimum height=1.5cm] (phys) at (-1.5,0) {\textbf{The Physician}\\Expert, greedy\\'Gold in phisik is a cordial'};
% Middle: Wife of Bath and Host
\node[draw, rounded corners, fill=popblue!20, text width=2.2cm, align=center, minimum height=1.5cm] (wife) at (1.5,1.5) {\textbf{Wife of Bath}\\Complex vitality\\Experience, authority};
\node[draw, rounded corners, fill=popblue!20, text width=2.2cm, align=center, minimum height=1.5cm] (host) at (1.5,-1.5) {\textbf{The Host}\\Master of ceremonies\\Merry, bold, organiser};
% Corrupt pair
\node[draw, rounded corners, fill=poporange!30, text width=2.2cm, align=center, minimum height=1.5cm] (sum) at (5,1) {\textbf{The Summoner}\\Corrupt court official\\Extortion, 'Questio quid juris'};
\node[draw, rounded corners, fill=popred!30, text width=2.2cm, align=center, minimum height=1.5cm] (pard) at (5,-1) {\textbf{The Pardoner}\\Climax of corruption\\Fake relics, greed};
% Labels
\node[align=center] at (-5,2.5) {\textbf{IDEAL}};
\node[align=center] at (-1.5,2.5) {\textbf{COMPROMISED}};
\node[align=center] at (1.5,3) {\textbf{AMBIGUOUS / VITAL}};
\node[align=center] at (5,2.5) {\textbf{CORRUPT}};
% Arrows showing contrast
\draw[<->, dashed, thick, popdark] (parson) -- (sum) node[midway, above, sloped, text width=3cm, align=center] {Churchmen: ideal vs corrupt};
\draw[<->, dashed, thick, popdark] (plow) -- (pard) node[midway, below, sloped, text width=3cm, align=center] {Lay/Church: virtue vs greed};
\draw[<->, dashed, thick, popdark] (phys) -- (pard) node[midway, above, sloped, text width=3cm, align=center] {Professional vs sacramental greed};
\end{tikzpicture}
\legende{Moral spectrum of the final General Prologue portraits: from ideal (Parson, Plowman) through compromised professional (Physician) and complex vitality (Wife of Bath, Host) to corruption (Summoner, Pardoner).}
\end{popfigure}

\courssub{Method: Comparing Two Portraits}

\begin{methode}[Writing a Comparison Paragraph]
When an exam question asks you to compare two portraits (e.g., 'Compare the Parson and the Pardoner'), follow this four-step structure:
\begin{enumerate}
\item \textbf{Similarity:} Identify a shared role or category (e.g., both are churchmen, both hold spiritual authority).
\item \textbf{Difference:} Contrast their moral practice (e.g., the Parson lives in poverty and visits his flock; the Pardoner exploits his flock for profit).
\item \textbf{Chaucer's Purpose:} Explain what this contrast achieves (e.g., it establishes a moral benchmark; the Parson shows what a priest \emph{should} be, making the Pardoner's corruption more shocking).
\item \textbf{Quotation:} Embed a short, precise quotation from each portrait to support your point.
\end{enumerate}
\end{methode}

\begin{exoresolu}[Worked Example: Comparing the Parson and the Pardoner]
\textbf{Question:} Compare Chaucer's presentation of the Parson and the Pardoner in the General Prologue.

\textbf{Solution:}
\begin{enumerate}
\item \textbf{Similarity:} Both are ecclesiastics who hold spiritual authority over laypeople — the Parson as a parish priest, the Pardoner as a papal pardoner licensed to sell indulgences.
\item \textbf{Difference:} The Parson embodies the ideal: he is 'povre' in goods but 'rich ... of hooly thought', visits his parishioners 'on his feet' in all weather, and lives by the principle 'If gold ruste, what shal iren do?' The Pardoner, by contrast, is a fraud who carries 'a pilwe-beer, which that he seyde was Oure Lady veyl' and 'a cros of latoun ful of stones', openly admitting 'I preche of no thyng but for coveitise.'
\item \textbf{Chaucer's Purpose:} By placing the Parson immediately before the corrupt pair (Summoner and Pardoner), Chaucer creates a \textbf{moral yardstick}. The Parson's genuine holiness proves that the Church's ideals are attainable; the Pardoner's grotesque hypocrisy is thus a betrayal of that ideal, not a reflection of it. The satire is savage against the Pardoner precisely because the Parson has shown us the alternative.
\item \textbf{Quotations:} Parson: 'First he wroghte, and afterward he taughte' (l. 498); 'If gold ruste, what shal iren do?' (l. 502). Pardoner: 'With feyned flaterye and Iapes / He made the persoun and the peple his apes' (ll. 705–6); 'I preche of no thyng but for coveitise' (from his Prologue, l. 403).
\end{enumerate}
\end{exoresolu}

\begin{exoresolu}[Worked Example: Comparing the Physician and the Pardoner]
\textbf{Question:} Compare Chaucer's presentation of the Physician and the Pardoner as figures of professional corruption.

\textbf{Solution:}
\begin{enumerate}
\item \textbf{Similarity:} Both are professionals who exploit their specialised knowledge for personal profit — the Physician through medical fees and apothecary kickbacks, the Pardoner through the sale of indulgences and fake relics.
\item \textbf{Difference:} The Physician possesses \textbf{genuine expertise} ('grounded in astronomye', knows authorities from Esculapius to Gilbertyn) and maintains a respectable exterior (moderate diet, sententious speech). The Pardoner has \textbf{no genuine spiritual power}; his 'relics' are pig's bones and pillowcases, and his only skill is 'feyned flaterye'.
\item \textbf{Chaucer's Purpose:} The Physician represents \textbf{secular professional corruption} — a good doctor gone bad. The Pardoner represents \textbf{sacramental corruption} — a holy office turned into a confidence trick. The Physician's greed is a perversion of his craft; the Pardoner's greed \emph{is} his craft.
\item \textbf{Quotations:} Physician: 'Of gold he loved best... For gold in phisik is a cordial' (ll. 442–444). Pardoner: 'I preche of no thyng but for coveitise' (Pardoner's Prologue, l. 403); 'He made the persoun and the peple his apes' (l. 706).
\end{enumerate}
\end{exoresolu}

\courssub{Consolidation Exercises}

\begin{exercice}[Practice Questions]
\begin{enumerate}
\item \textbf{Close reading:} Annotate the description of the Summoner (ll. 623–666). Identify three details that link his physical appearance to his moral corruption.
\item \textbf{Comparison:} Using the method above, write a paragraph comparing the Wife of Bath and the Prioress (from Section 3) as portraits of women in the General Prologue.
\item \textbf{Essay planning:} Plan an answer to: 'Chaucer's satire in the General Prologue is always gentle.' Discuss with reference to at least three portraits from Sections 3 and 4.
\item \textbf{Contextual link:} Explain how the Parson's principle 'If gold ruste, what shal iren do?' reflects Wycliffite ideas about clerical poverty and the responsibility of the clergy.
\item \textbf{Close reading:} Annotate the Physician's portrait (ll. 413–444). How does Chaucer use the catalogue of medical authorities to build irony?
\end{enumerate}
\end{exercice}

\begin{corrige}[Exercise 1]
\begin{itemize}
\item \textbf{Fiery face / carbuncles / whelkes} $\rightarrow$ suggests inner fire of lust and corruption; the disease mirrors the spiritual disease of extortion.
\item \textbf{Love of garlic, onions, leeks, strong wine} $\rightarrow$ sensual indulgence, lack of ascetic discipline expected of a church officer.
\item \textbf{Garland on his head 'as it were a tavern sign'} $\rightarrow$ he identifies with the alehouse, not the church; his 'office' is a licence for revelry.
\end{itemize}
\end{corrige}

\begin{corrige}[Exercise 5]
\begin{itemize}
\item The catalogue (Esculapius, Dioscorides, Rufus, Hippocrates, Galen, Serapion, Rhazes, Avicenna, Averroes, Damascenus, Constantine, Bernard, Gatesden, Gilbertyn) establishes the Physician's \textbf{genuine learning} — he knows the full tradition from antiquity through Arabic medicine to contemporary English practice.
\item This makes the final couplet more ironic: all this learning serves only \textbf{'coveitise'}. The authorities are named not to honour them but to show how far he has fallen from their ethical standards (e.g., the Hippocratic Oath).
\item The list also satirises \textbf{scholastic pedantry} — name-dropping as a substitute for true wisdom or charity.
\end{itemize}
\end{corrige}

\begin{aretenir}[Section 4 — Key Takeaways]
\begin{itemize}
\item The \textbf{Physician} is a competent professional corrupted by greed — his learning is real, his ethics are not.
\item The \textbf{Wife of Bath} is a \textbf{complex, vital figure} — not a simple satire; her economic independence, travel, and marital experience make her a unique voice.
\item The \textbf{Parson and Plowman} form an ideal pair: they prove Chaucer respects the Church's ideals and targets only hypocrites.
\item The \textbf{Summoner and Pardoner} represent institutional corruption: the Summoner sells justice, the Pardoner sells salvation; both are motivated by greed.
\item The \textbf{Host} is the narrative architect; his storytelling contest creates the frame narrative.
\item \textbf{Structural arrangement} is a moral argument: ideal $\rightarrow$ compromised professional $\rightarrow$ vital/ambiguous $\rightarrow$ corrupt.
\item For essays: use the \textbf{comparison method} (similarity $\rightarrow$ difference $\rightarrow$ purpose $\rightarrow$ quotation) and the \textbf{gentle vs savage} framework.
\end{itemize}
\end{aretenir}

\coursec{Section 5: The Merchant's Tale --- Structure, Prologue and the Marriage Debate}

In the previous section we left the General Prologue with a sharp contrast in our minds: the ideal Christians (the Parson and the Plowman) standing opposite the corrupt churchmen (the Summoner and the Pardoner). As the pilgrimage moves from \emph{portraits} to \emph{tales}, that same opposition between sincerity and corruption reappears --- but now it migrates into the most intimate of human institutions: \cle{marriage}. The Merchant, whom we met in the General Prologue as a prosperous, self-important trader, steps forward to tell his tale, and what he produces is one of the darkest, funniest and most uncomfortable stories in \emph{The Canterbury Tales}. To understand it fully, we must proceed in three stages: the teller, the structure, and the debate at the heart of the tale.

\courssub{Background: The Merchant and His Bitter Prologue}

\begin{experience}[Observing the Teller Before the Tale]
Before reading a single line of the tale, look at the teller. In the \emph{General Prologue} (ll. 270--284), the Merchant sits high on his horse, forked beard, Flemish beaver hat, boots clasped elegantly. He speaks \emph{solemnly} of his profits, his bargains, the need to keep the sea free for trade. He appears dignified, successful. Yet the narrator adds a devastating aside: \emph{``This worthy man full well his wit beset; / There wiste no wight that he was in debt''} (ll. 281--282). His prosperity is a performance; his dignity masks anxiety.
\par\smallskip
Now turn to the \cle{Prologue to the Merchant's Tale} (ll. 1217--1248). The mask drops completely. He has been married \textbf{two months} --- \emph{``two monthes, more or lesse''} --- and declares: \emph{``Weeping and wailing, care and other sorrow / I knowe enough, on even and a-morrow''} (ll. 1223--1224). He contrasts his wife with the patient Griselda of the Clerk's Tale (just told): \emph{``She is a shrew at all times''} (l. 1234). He swears he would rather be wedded to a snake. This is not a man telling a story; it is a man \emph{venting}.
\end{experience}

\begin{methode}[Connecting Prologue to Tale --- Examination Technique]
To analyse a tale's prologue, always connect the teller's personal situation to the tale's tone and bias:
\begin{enumerate}
\item Identify what the teller reveals about his or her own life in the prologue (here: two-month marriage, misery, comparison with Griselda).
\item State the emotion this produces (bitterness, humiliation, desire for revenge).
\item Show how that emotion shapes the choice of story, the narrator's asides, and the treatment of characters (the tale becomes a fabliau where marriage is a trap; the encomium of marriage is steeped in irony; the ending denies justice).
\item Conclude on the effect: the tale becomes a form of \emph{revenge} --- the Merchant punishes marriage itself by painting it as a theatre of lust, deception and ridicule.
\end{enumerate}
\end{methode}

\begin{aretenir}[Key Idea]
The Merchant's Tale is not a neutral story. It is the revenge of a newly married, deeply unhappy man. Every sarcastic remark about January, every cruel twist suffered by the characters, carries the bitterness of the teller's own two-month marriage. In the examination, always link \textbf{teller} to \textbf{tale}.
\end{aretenir}

\courssub{The Structure of the Tale}

The tale unfolds like a well-oiled machine, each part preparing the next. Follow the plot-line carefully:

\begin{popfigure}
\begin{center}
\begin{tikzpicture}[scale=1.0, every node/.style={font=\small}]
\draw[->, thick, popline] (0,0) -- (12.6,0);
\foreach \x/\lab in {0.4/{Prologue}, 2.4/{Decision\\ to marry}, 4.4/{Marriage\\ debate}, 6.4/{Wedding}, 8.4/{Garden \&\\ pear tree}, 10.4/{Blindness \&\\ trick}, 12.2/{Sight\\ restored}}{
  \fill[popblue] (\x,0) circle (2.2pt);
  \node[align=center, text width=1.7cm, above=4pt] at (\x,0) {\lab};
}
\node[below=8pt, popmuted] at (0.4,0) {Merchant's bitterness};
\node[below=8pt, popmuted] at (2.4,0) {January, aged 60};
\node[below=8pt, popmuted] at (4.4,0) {Placebo vs Justinus};
\node[below=8pt, popmuted] at (6.4,0) {January \& May};
\node[below=8pt, popmuted] at (8.4,0) {Damyan's letter};
\node[below=8pt, popmuted] at (10.4,0) {May \& Damyan};
\node[below=8pt, popmuted] at (12.2,0) {May's ``cure''};
\end{tikzpicture}
\end{center}
\legende{Plot-line of The Merchant's Tale: from the Merchant's bitter Prologue to the pear-tree trick and the restoration of January's sight.}
\end{popfigure}

\begin{enumerate}
\item \textbf{The Prologue (ll. 1217--1248):} the Merchant laments his own marriage, curses the Clerk's Griselda as an impossible ideal, and announces his tale will show the ``woe that is in marriage''.
\item \textbf{January's decision (ll. 1249--1300):} a wealthy, sixty-year-old knight of Lombardy decides, after a lifetime of lechery (\emph{``bodily delit''}), to marry --- officially for piety and an heir, in reality for a young and beautiful body.
\item \textbf{The marriage debate (ll. 1301--1480):} January consults his friends. Placebo flatters him; Justinus warns him. January follows Placebo.
\item \textbf{The wedding and encomium (ll. 1481--1700):} January marries the young and ``tender'' May, chosen like merchandise from among the town's young women. The narrator inserts a long, mock-solemn \cle{encomium} (formal praise) of marriage, heavy with irony.
\item \textbf{The garden and the pear tree (ll. 1701--1850):} January builds a private walled garden for his pleasures. Meanwhile his squire \cle{Damyan} falls sick with love for May, passes her a letter (which she reads in the \cle{privy} and destroys), and May copies the garden key in wax.
\item \textbf{The blind January (ll. 1851--1950):} January loses his sight and becomes pathologically jealous, holding May by the hand everywhere.
\item \textbf{The trick and the ``cure'' (ll. 1951--2120):} in the garden, May climbs the pear tree into Damyan's arms. The gods Pluto and Proserpina, quarrelling nearby, restore January's sight at the fatal moment --- but May, quick-witted, convinces January that his eyes deceive him and that her struggle in the tree was the ``cure'' that restored his vision. January believes her.
\end{enumerate}

\courssub{January's Delusion: Lust Disguised as Piety}

Here is the psychological core of the tale, and examiners love it. January does not merely marry; he \emph{constructs an elaborate self-deception} to make his lust look respectable.

\begin{exemplebox}[Analysing January's Self-Deception --- Close Reading]
Consider how January frames his project (ll. 1250--1300). He speaks of marriage as a \emph{``holy sacrament''}, a \emph{``gift of God''}, of his desire for an \emph{``heir''}. Yet his criteria for choosing a wife betray him completely: she must be \textbf{young} (he rejects any woman over thirty: \emph{``over thirty yeer... is but a bene-straw''}, l. 1380), \textbf{beautiful}, and \textbf{small and tender} (\emph{``tender and small''}, l. 1385). He compares old men married to young wives to the fortunate, and imagines marriage as a paradise on earth: \emph{``A wyf is Goddes yifte verraily''} (l. 1330) but he wants that gift to be a \emph{``fresshe''} (fresh) one.
\par\smallskip
The rigour of the analysis lies in the gap: \textbf{what January says} (piety, sacrament, heir, comfort in age) versus \textbf{what January wants} (a young body, available at will, in his private garden). Chaucer lets January condemn himself out of his own mouth --- a technique called \emph{self-revealing irony}. The narrator's mock-solemn encomium (ll. 1481--1600) deepens the irony: the more marriage is ``praised'' with inflated rhetoric (\emph{``O blisful ordre of wedlok precious''}), the more grotesque January's motives appear.
\end{exemplebox}

\begin{attention}[Do Not Treat This as Simple Comedy]
Do not treat \emph{The Merchant's Tale} as a simple comedy. Its bitterness and cruelty give it a dark edge that examiners expect you to notice: January's lust is predatory (he selects May like livestock at a market: \emph{``as who saith, 'Taketh oon, and lat the other go''}, l. 1420), May's adultery is cold and calculating, and Damyan's ``love'' is mere appetite (\emph{``so sore hath Venus hurt him with hire brond''}, l. 1720). Nobody in the tale is innocent. The laughter it provokes is uncomfortable laughter --- the laughter of the Merchant, not of Chaucer.
\end{attention}

\courssub{The Marriage Debate: Placebo versus Justinus}

Before marrying, January summons his friends and asks their counsel. This scene is a masterpiece of social observation, and it teaches a timeless lesson about how people misuse advice.

\textbf{Placebo} (his very name means ``I shall please'' in Latin, from the Office of the Dead: \emph{Placebo Domino in regione vivorum}) is the professional flatterer. He declares, in substance, \emph{``I have been a court-man all my life''} (l. 1350), and he has survived by never contradicting great men. His advice: marry, and do exactly as you please. He even flatters January's judgement as superior to Solomon's (who had many wives). He cites Seneca to support January's wish.

\textbf{Justinus} (``the just one'') is the honest counsellor. He warns January soberly: a wife may be a blessing or a trial (\emph{``a wyf is a mannes confusion''}, l. 1400, citing Theophrastus); a young wife to an old man is a danger (\emph{``a yong wyf and an old husbonde''}, l. 1410); think carefully before you bind yourself. He speaks from experience and good sense, citing the \emph{Golden Book of Marriage} and Jerome's \emph{Adversus Jovinianum}.

January's reaction is the point of the scene: he thanks Placebo warmly (\emph{``Now, brother,'' quod January, ``stint''}, l. 1470) and dismisses Justinus with irritation. He never wanted advice; he wanted \emph{permission}.

\begin{popfigure}
\begin{center}
\begin{tikzpicture}[font=\small]
\node[draw=popgreen, fill=popgreenL, rounded corners, text width=5.2cm, align=left, minimum height=3.4cm] (plac) at (0,0)
{\textbf{PLACEBO --- the flatterer}\\[2pt]
$\bullet$ ``I have been a court-man all my life'' (l. 1350)\\
$\bullet$ Never contradicts a lord\\
$\bullet$ Praises January's wisdom > Solomon\\
$\bullet$ Cites Seneca for pleasure\\
$\bullet$ Verdict: marry, do as you please};
\node[draw=popblue, fill=popblueL, rounded corners, text width=5.2cm, align=left, minimum height=3.4cm] (just) at (9.4,0)
{\textbf{JUSTINUS --- the honest counsellor}\\[2pt]
$\bullet$ Speaks from experience\\
$\bullet$ A wife: blessing \emph{or} trial (Theophrastus)\\
$\bullet$ Old man + young wife = danger\\
$\bullet$ Cites Jerome, Golden Book\\
$\bullet$ Verdict: reflect, be careful};
\node[draw=poporange, fill=poporange!15, rounded corners, text width=3.2cm, align=center] (jan) at (4.7,-3.4)
{\textbf{JANUARY}\\ hears only what\\ he wants to hear};
\draw[->, very thick, popgreen] (plac.south) -- (jan.north west) node[midway, left, popgreen]{chosen};
\draw[->, very thick, popblue, dashed] (just.south) -- (jan.north east) node[midway, right, popblue]{rejected};
\end{tikzpicture}
\end{center}
\legende{The marriage debate: Placebo's flattery against Justinus's honest counsel, with January's self-deceiving choice highlighted.}
\end{popfigure}

\begin{aretenir}[The Test January Fails]
The marriage debate is a test January fails. He consults advisers only to confirm his own desire. This is a timeless lesson about \cle{self-deception}: whoever asks for advice while having already decided will always find a Placebo to flatter him and will always resent the Justinus who tells the truth. The scene is also Chaucer's satire of court culture, where flattery is a career and honesty a risk. Note the names: \emph{Placebo} (I shall please) and \emph{Justinus} (the just one) are \cle{allegorical names} signalling their function.
\end{aretenir}

\courssub{Damyan and May: The Fabliau Machinery}

While January plays the devoted husband, the comic-obscene plot engine starts turning. Damyan, January's own squire, falls so violently in ``love'' with May that he takes to his bed (\emph{``lovesick''}, l. 1720). He writes her a letter; May reads it in the \cle{privy} (the toilet --- Chaucer's deliberately crude detail: \emph{``in the privee''}, l. 1780) and flushes the evidence away. She answers, arranges everything, and has the garden key copied in wax (\emph{``the waxen cliche''}, l. 1820). Note the efficiency: May, the supposed passive child-bride, is in fact the most decisive plotter in the tale. When January goes blind and clings to her hand out of jealousy (\emph{``he held hire by the hand''}, l. 1900), she simply signals Damyan to wait in the pear tree.

\begin{definition}[Fabliau]
A \cle{fabliau} is a short, comic, often bawdy verse tale popular in medieval French and English literature. Its typical ingredients are trickery, sexual deception, a gullible (often old or jealous) husband, a clever young wife and a lusty young clerk or squire. The humour is physical and cynical, and the ``clever'' characters triumph over the ``respectable'' ones. \emph{The Merchant's Tale} is a sophisticated fabliau: it borrows the machinery but adds psychological depth and bitter irony.
\end{definition}

\begin{definition}[Dramatic Irony]
\cle{Dramatic irony} occurs when the audience knows more than a character does. In the garden scene, we see May climb the pear tree into Damyan's arms while the blind January stands below, suspecting nothing. When his sight is miraculously restored at the worst possible moment, the irony doubles: he \emph{sees} the truth, yet May talks him out of believing his own eyes --- the ultimate dramatisation of wilful blindness. Her speech (ll. 2080--2100) is a masterpiece of gaslighting: \emph{``I have yow holpen on youre eyen twye''} (I have helped your two eyes).
\end{definition}

\courssub{Genre, Themes and Examination Practice}

The tale fuses two genres: the \textbf{fabliau} (bawdy trickery) and the \textbf{encomium} (a formal praise of marriage), used ironically. The narrator praises marriage in inflated, solemn terms while showing us a marriage made of lust, commerce and lies. The bitterness of the Merchant-teller saturates both.

\begin{center}
\begin{tabularx}{\linewidth}{|l|X|}
\hline
\rowcolor{popblueL} \textbf{Theme} & \textbf{How it appears in the tale} \\
\hline
Marriage as transaction & January ``shops'' for a wife (ll. 1410--1420); May trades her body for security and pleasure. \\
\hline
Age versus youth & Sixty-year-old January against young May and Damyan; youth wins through cunning. \\
\hline
Sight and blindness & Literal blindness mirrors moral blindness; restored sight changes nothing because January refuses to see. \\
\hline
Flattery versus truth & Placebo is rewarded, Justinus dismissed --- the debate prefigures the whole plot. \\
\hline
Authority and experience & The tale pits \emph{auctoritee} (cited authors) against \emph{experience} (the Merchant's, January's, May's). \\
\hline
\end{tabularx}
\end{center}

\begin{exoresolu}[Worked Examination Paragraph]
\textbf{Question:} ``The marriage debate in The Merchant's Tale shows that January's tragedy is self-inflicted.'' Discuss with close reference to the text.
\tcblower
\textbf{Model paragraph:} January's downfall is prepared long before the pear tree, in the very scene where he seeks advice. Chaucer stages the consultation as a clear moral test: on one side stands Placebo, whose name --- ``I shall please'' --- announces his function, and who boasts that he has been a court-man all his life by never opposing a lord (l. 1350); on the other stands Justinus, who warns soberly that a wife may prove a blessing or a trial (citing Theophrastus, l. 1400) and that an old man wedding a young girl invites disaster (l. 1410). January's response reveals that he never wanted counsel at all, only confirmation: he embraces the flatterer (l. 1470) and brushes aside the honest man. The debate thus functions as a miniature of the whole tale --- January will later refuse to believe even his own restored eyes, just as here he refuses to hear Justinus. His blindness, literal and moral, is chosen, not imposed; and through it Chaucer, channelling the bitter Merchant, suggests that marriages built on lust and flattery contain their own punishment.
\end{exoresolu}

\begin{exercice}[Practice Questions --- GCE Style]
\begin{enumerate}
\item Explain how the Merchant's Prologue determines the tone of his tale. Support your answer with two precise references to the text.
\item Write a detailed essay plan on: ``Sight and blindness, literal and moral, in The Merchant's Tale.''
\item Compare Placebo and Justinus as satirical portraits of court life and counsellors. How does January's choice between them characterise him?
\item \textbf{Passage-based:} Read the encomium of marriage (ll. 1481--1520). How does Chaucer use irony in this passage to undermine January's motives?
\end{enumerate}
\end{exercice}

\begin{corrige}[Exercise 1]
The Prologue reveals that the Merchant has been married only two months and already considers his wife a torment, contrasting her with the patient Griselda of the Clerk's Tale (ll. 1225--1235). This personal bitterness (i) dictates his choice of a story in which marriage is a trap of lust and deception rather than a sacrament, and (ii) colours the narration with sarcasm --- for instance the mock-solemn encomium praising marriage (ll. 1481--1600) while January selects May like merchandise (l. 1420), and the cynical ending in which May's lie is believed and January is duped into thinking her adultery cured him. The tale is thus best read as the teller's revenge on marriage itself.
\end{corrige}

\begin{savaistu}[Did You Know?]
The quarrel between Pluto and Proserpina (ll. 1980--2050), who restore January's sight mid-trick, is Chaucer's ironic borrowing of classical myth: the god and goddess of the underworld argue about whether women are faithless, and each ``proves'' their case through January and May. Pluto restores January's sight to punish May; Proserpina gives May the wit to talk her way out of it. It is a reminder that in medieval estates satire, even pagan gods could be drafted into a Christian-age comedy of marriage --- and that Chaucer's pilgrims, traders and churchmen alike, would have relished the joke. The pear tree itself echoes the Garden of Eden: a forbidden fruit, a serpent (Damyan), a fall --- but here the ``Eve'' outwits both God and husband.
\end{savaistu}

\coursec{Section 6: A Raisin in the Sun Act 2 Scene 1 — American Dream, Money and Cultural Identity}

Having explored the medieval world of Chaucer's pilgrims and the intricate marriage debates of \emph{The Merchant's Tale}, we turn now to mid‑twentieth‑century Chicago. Lorraine Hansberry's \emph{A Raisin in the Sun} (1959) shifts our attention from the Canterbury road to a cramped apartment on the South Side, where a Black family awaits a \$10,000 insurance cheque. The play's title comes from Langston Hughes's poem \emph{Harlem}: \emph{``What happens to a dream deferred? / Does it dry up / like a raisin in the sun?''} In Act 2 Scene 1, the deferral becomes acute: the money arrives, dreams collide, and the family's identity — as Africans, as Americans, as a unit — is tested. This scene is the dramatic hinge on which the rest of the play turns.

\courssub{Recap of Act 1: The Insurance Cheque and Conflicting Dreams}

Before we enter Act 2 Scene 1, we must fix the situation established in Act 1. The Younger family — Mama (Lena), her son Walter Lee, his wife Ruth, their son Travis, and Walter's sister Beneatha — live in a worn‑out three‑room flat. The catalyst is the \$10,000 life‑insurance payment following the death of Big Walter, Mama's husband and the children's father. Each character invests the money with a different dream:

\begin{popfigure}
\begin{tikzpicture}[scale=0.9, every node/.style={font=\small}]
% Central cheque
\node[draw=popblue, thick, fill=popblueL, rounded corners, minimum width=3.5cm, minimum height=1.5cm, align=center, font=\bfseries\small] (cheque) at (0,0){\$10,000\\Insurance Cheque};

% Mama's dream - House
\node[draw=popgreen, thick, fill=popgreenL, rounded corners, minimum width=3cm, minimum height=2cm, text width=2.8cm, align=center] (mama) at (-5,3) {\textbf{Mama (Lena)}\\\emph{A house with a garden}\\``A place for Travis to play''\\Down payment: \$3,500};
\draw[->, thick, popgreen] (cheque.north west) -- (mama.south east) node[midway, above, sloped, font=\tiny] {Security, roots, dignity};

% Walter's dream - Liquor store
\node[draw=poporange, thick, fill=poporange!20, rounded corners, minimum width=3cm, minimum height=2cm, text width=2.8cm, align=center] (walter) at (5,3) {\textbf{Walter Lee}\\\emph{Liquor store investment}\\``Business partnership with Willy \& Bobo''\\Full amount: \$10,000};
\draw[->, thick, poporange] (cheque.north east) -- (walter.south west) node[midway, above, sloped, font=\tiny] {Manhood, economic power, respect};

% Beneatha's dream - Medical school
\node[draw=poppurple, thick, fill=poppurple!20, rounded corners, minimum width=3cm, minimum height=2cm, text width=2.8cm, align=center] (beneatha) at (0,-3.5) {\textbf{Beneatha}\\\emph{Medical school tuition}\\``I'm going to be a doctor''\\Tuition: ~\$3,000+/year};
\draw[->, thick, poppurple] (cheque.south) -- (beneatha.north) node[midway, right, font=\tiny] {Intellectual fulfilment, independence};

% Conflict arrows
\draw[<->, dashed, thick, popdark] (mama.east) to[bend left=20] node[midway, above, font=\tiny, text width=2.5cm, align=center] {Clash: practicality vs. risk} (walter.west);
\draw[<->, dashed, thick, popdark] (mama.south east) to[bend right=15] node[midway, right, font=\tiny, text width=2.5cm, align=center] {Generational tension} (beneatha.north west);
\draw[<->, dashed, thick, popdark] (walter.south west) to[bend left=15] node[midway, left, font=\tiny, text width=2.5cm, align=center] {Sibling rivalry} (beneatha.north east);
\end{tikzpicture}
\legende{The three competing dreams anchored on the \$10,000 insurance cheque. Each dream reflects a different vision of the American Dream.}
\end{popfigure}

\begin{definition}[The American Dream in \emph{A Raisin in the Sun}]
The \cle{American Dream} is the national ethos that every citizen, through hard work and determination, can achieve prosperity, upward mobility and a better life for their children. In Hansberry's play, this dream is \emph{tested against racial reality}: the Youngers are Black in 1950s America, where segregation, redlining and employment discrimination systematically block the very pathways the dream promises. The insurance cheque is a rare lump sum — a ``stake'' — that could buy entry, but the family disagrees on \emph{how} to spend it.
\end{definition}

\courssub{Scene 1 Opening: Cultural Identity and Afrocentrism versus Assimilation}

Act 2 Scene 1 opens not with the cheque but with \cle{cultural performance}. Beneatha has cut her hair short and wears a Nigerian robe (\emph{``a very beautiful piece of clothing''}) sent by Joseph Asagai, a Nigerian student she met on campus. She dances to a record of Nigerian folk music, moving ``with a serious, almost religious concentration''. This is a deliberate \emph{rejection of assimilation} — the pressure to conform to white, middle‑class norms.

Enter \cle{George Murchison}, Beneatha's wealthy, college‑educated boyfriend. He represents the \cle{assimilationist} path: straight hair, Western suits, dismissal of African heritage as ``garbage''. When he sees Beneatha's natural hair, he says: \emph{``You look eccentric... What have you done to your head?''} Beneatha retorts: \emph{``It's natural. It's me.''}

\begin{popfigure}
\begin{tikzpicture}[scale=0.95, every node/.style={font=\small}]
% Continuum line
\draw[thick, popline] (-6,0) -- (6,0);
\foreach \x in {-5,-3,-1,1,3,5} \draw (\x,-0.1) -- (\x,0.1);

% Labels
\node[below=0.3cm] at (-5,0) {Assimilation};
\node[below=0.3cm] at (5,0) {Afrocentrism / Cultural Nationalism};

% George Murchison
\node[draw=poporange, thick, fill=poporange!15, rounded corners, minimum width=2.5cm, minimum height=1.2cm, text width=2.2cm, align=center] (george) at (-4,1.5) {\textbf{George Murchison}\\\emph{``We're Americans!''}\\Rejects African past; seeks acceptance in white society};
\draw[->, poporange] (george.south) -- (-4,0.2);

% Asagai
\node[draw=popgreen, thick, fill=popgreenL, rounded corners, minimum width=2.5cm, minimum height=1.2cm, text width=2.2cm, align=center] (asagai) at (4,1.5) {\textbf{Joseph Asagai}\\\emph{``You are not fully alive''}\\Embraces African identity; sees struggle as global};
\draw[->, popgreen] (asagai.south) -- (4,0.2);

% Beneatha - positioned between but leaning toward Asagai
\node[draw=poppurple, thick, fill=poppurple!15, rounded corners, minimum width=2.8cm, minimum height=1.5cm, text width=2.5cm, align=center] (beneatha) at (0.5,-1.8) {\textbf{Beneatha Younger}\\\emph{Searching}\\``I'm looking for my identity!''\\Rejects George's shallowness; drawn to Asagai's vision};
\draw[->, poppurple] (beneatha.north) -- (0.5,-0.2);

% Anchors
\node[font=\tiny, popmuted] at (-5,-0.8) {Straight hair, Western clothes, country club};
\node[font=\tiny, popmuted] at (5,-0.8) {Natural hair, African robes, Yoruba language, pan-Africanism};
\end{tikzpicture}
\legende{Beneatha's identity continuum: assimilation (George Murchison) $\leftrightarrow$ Afrocentrism (Asagai). Her choice is political, not merely romantic.}
\end{popfigure}

\begin{attention}[Beneatha's Storyline Is Not a Romance Subplot]
Do not reduce Beneatha's choice between Asagai and George to a love triangle. Her decision is a \cle{choice between two visions of Black identity}:
\begin{itemize}
\item \textbf{George} embodies \emph{integration at the cost of erasure} — ``assimilationist'' in Hansberry's terms. He tells Beneatha: \emph{``You read books... to get a degree. That's all.''} He wants a wife who decorates his status.
\item \textbf{Asagai} embodies \emph{cultural recovery and political consciousness}. He challenges her: \emph{``You are not fully alive... because you have no past.''} He offers a connection to a pre‑colonial history and a future of African liberation.
\end{itemize}
Examiners reward candidates who see Beneatha as the play's \emph{intellectual and ideological conscience}.
\end{attention}

\courssub{Walter's Despair and Mama's Response: ``We Ain't Never Been That Dead Inside''}

While Beneatha dances, Walter is absent — drinking, wandering, despairing. He returns drunk, bitter, feeling his manhood denied because Mama controls the money. He lashes out: \emph{``You run our lives like you want to... You butchered up a dream of mine!''}

Mama's response is the moral centre of the scene. She reveals she has already spent \$3,500 as a down payment on a house in \cle{Clybourne Park} — a white neighbourhood. She entrusts the remaining \$6,500 to Walter: \$3,000 for Beneatha's medical school, \$3,500 for his own discretion. Her speech is a masterclass in \cle{matriarchal authority} and \cle{historical memory}:

\begin{aretenir}[Mama's Defining Speech (Act 2, Scene 1)]
\begin{quote}
\emph{``I come from five generations of people who was slaves and sharecroppers... but ain't nobody in my family never let nobody pay 'em no money that was a way of telling us we wasn't fit to walk the earth. We ain't never been that poor... We ain't never been that dead inside.''}
\end{quote}
\textbf{Key points for analysis:}
\begin{itemize}
\item \textbf{Historical continuum:} She links the present to slavery and sharecropping — the family's survival is an act of resistance.
\item \textbf{Dignity over money:} The refusal to be ``bought out'' (later, by Lindner) is prefigured here.
\item \textbf{Trust as empowerment:} Giving Walter the money is a calculated risk to restore his agency.
\item \textbf{The phrase ``dead inside'':} Spiritual death is worse than material poverty.
\end{itemize}
\end{aretenir}

\courssub{Mama's Decision: The Plot's Crucial Turning Point}

Mama's purchase of the house and her entrusting of the balance to Walter constitute the \cle{peripeteia} (reversal) of the play's rising action. Structurally:

\begin{methode}[Tracking the Turning Point in Act 2 Scene 1]
\begin{enumerate}
\item \textbf{Before the scene:} Walter holds the dream of the liquor store; Mama holds the cheque; power is with Mama.
\item \textbf{The action:} Mama spends part of the money on a house (her dream) and gives Walter the rest (his chance).
\item \textbf{Immediate effect:} Walter is transformed — sober, hopeful, visionary.
\item \textbf{Dramatic irony:} The audience knows the money is vulnerable (Willy Harris is untrustworthy; the house is in a hostile white area). Walter's hope is \emph{precarious}.
\item \textbf{Consequence:} This decision sets up the catastrophe of Act 2 Scene 2/3 (the loss of the money) and the climax of Act 3 (the Lindner confrontation).
\end{enumerate}
\end{methode}

\courssub{Walter's Renewed Hope: Dramatic Irony and the Speech to Travis}

With the money in his hands, Walter undergoes a striking metamorphosis. He sits with his son Travis at bedtime and paints a luminous future:

\begin{exemplebox}[Walter's Vision Speech — Close Analysis]
\textbf{Extract (paraphrased):} \emph{``Just tell me where you want to go to school... You just name it... and I'll hand you the world!''}

\textbf{Technique: Dramatic Irony}
\begin{itemize}
\item \textbf{Walter's state:} He believes the \$3,500 is his to invest; he feels ``like a man'' for the first time.
\item \textbf{Audience's knowledge:} We have seen Willy Harris's unreliability (Act 1); we sense the money will be stolen. The house in Clybourne Park will provoke racist backlash.
\item \textbf{Effect:} The tenderness of the father‑son moment is shadowed by impending betrayal. The ``world'' he promises is built on sand.
\item \textbf{Exam tip:} When asked about dramatic irony, quote the contrast between Walter's \emph{``I'm going to make a transaction... that's going to change our lives''} and the audience's foreboding.
\end{itemize}
\end{exemplebox}

\courssub{Theme Analysis: The American Dream Deferred — Money as Opportunity and Poison}

Hansberry uses the \$10,000 to expose the \cle{dual nature of money} in a racially structured society:

\begin{propriete}[Money as Both Opportunity and Poison]
\begin{itemize}
\item \textbf{Opportunity:} The cheque is the only capital the family has ever possessed. It can buy a house (asset accumulation), fund education (human capital), or start a business (entrepreneurship). These are the classic levers of the American Dream.
\item \textbf{Poison:} Money \emph{divides} the family (Walter vs. Mama, Walter vs. Beneatha). It exposes Walter's vulnerability to exploitation (Willy Harris). It attracts the racist buyout offer (Lindner). It makes the family a target.
\item \textbf{Hughes's poem realised:} The dream ``dries up like a raisin in the sun'' when Walter loses the money; it ``festers like a sore'' in his bitterness; it ``stinks like rotten meat'' in the moral decay of considering Lindner's offer; it ``crusts and sugars over'' in Mama's weary patience; it ``sags like a heavy load'' on Ruth's pregnancy dilemma; it may ``explode'' in the final act of defiance.
\end{itemize}
\end{propriete}

\begin{savaistu}[Cameroonian Parallel: The ``Chop Money'' Dilemma]
In many Cameroonian families, a sudden windfall — a gratuity, a lottery win, a diaspora remittance — triggers exactly the same tensions Hansberry dramatises: the elder wants to build a house in the village (Mama), the young man wants to start a business or buy a taxi (Walter), the daughter needs university fees (Beneatha). The \cle{communal vs. individual} allocation of scarce capital is a universal dramatic situation, not only an African‑American one.
\end{savaistu}

\courssub{Dramatic Technique: Music, Costume and Domestic Setting}

Hansberry's stagecraft in this scene is precise and symbolic:

\begin{center}
\begin{tabularx}{\linewidth}{|l|X|X|}
\hline
\rowcolor{popblueL}
\textbf{Element} & \textbf{Description in Scene} & \textbf{Dramatic Function} \\
\hline
\textbf{Music} & Nigerian folk record (Beneatha); later, blues/jazz on radio (Walter's mood) & \emph{Cultural counterpoint}: African roots vs. African‑American urban reality. Music marks Beneatha's ideological shift. \\
\hline
\textbf{Costume} & Beneatha's Nigerian robe + natural hair; George's suit; Walter's rumpled clothes & \emph{Visual coding of identity}: Beneatha's robe is a \cle{prop of resistance}; George's suit signals assimilation; Walter's dishevelment shows internal collapse. \\
\hline
\textbf{Domestic setting} & Cramped kitchen/living room; morning light; Travis's bed on the sofa & \emph{Pressure cooker}: The lack of space forces conflict into the open. The ``rat trap'' apartment makes the house purchase visceral. \\
\hline
\textbf{Lighting} & ``Gray morning light'' at opening; warmer when Walter receives money & \emph{Mood tracking}: Bleakness $\to$ false dawn $\to$ (later) harsh reality. \\
\hline
\end{tabularx}
\end{center}

\begin{methode}[How to Analyse a Drama Scene for GCE Advanced Level]
When the exam asks you to comment on a scene, structure your answer around four pillars:
\begin{enumerate}
\item \textbf{Plot advancement:} What new information? What decisions? What changes? (Here: house bought, money given to Walter.)
\item \textbf{Character revelation:} What do we learn? How do they change? (Walter's vulnerability $\to$ hope; Mama's wisdom; Beneatha's politicisation.)
\item \textbf{Stagecraft:} Music, costume, props, lighting, set, entrances/exits. (Nigerian record, robe, cheque envelope, Travis's bed.)
\item \textbf{Theme:} Which major themes are developed? (American Dream, identity, racism, family, money, dignity.)
\end{enumerate}
\textbf{Always link technique to meaning.} Do not list devices; explain \emph{why} Hansberry chooses them at this moment.
\end{methode}

\courssub{Worked Example: Exam‑Style Question}

\begin{exoresolu}[Question]
\textbf{``In Act 2 Scene 1 of \emph{A Raisin in the Sun}, Hansberry uses the character of Beneatha to explore the tension between assimilation and cultural identity.'' Discuss this statement with close reference to the scene.}
\tcblower
\textbf{Model Response Plan:}

\textbf{Introduction:}
\begin{itemize}
\item Identify the scene's place: pivotal moment after cheque arrival, before Clybourne Park crisis.
\item Thesis: Beneatha is the \emph{ideological site} where the play debates Black identity; her costume, language and suitors externalise this debate.
\end{itemize}

\textbf{Paragraph 1: Visual and Performative Resistance (Costume + Dance)}
\begin{itemize}
\item Nigerian robe, natural hair, folk music — \emph{``she is dancing... with a serious, almost religious concentration''}.
\item Contrast with Act 1: straightened hair, Western clothes. The change is \emph{deliberate and political}.
\item Stage direction: ``a very beautiful piece of clothing'' — beauty as reclamation.
\end{itemize}

\textbf{Paragraph 2: The Two Suitors as Embodied Ideologies}
\begin{itemize}
\item \textbf{George Murchison:} ``assimilationist junk'', dismisses African heritage, values status (``I go to college to get a degree''). Represents \emph{integration without power}.
\item \textbf{Joseph Asagai:} Gifts the robe, teaches Yoruba, challenges her: ``You are not fully alive.'' Represents \emph{rootedness and pan-African consciousness}.
\item Beneatha's rejection of George (``You're a fool'') and embrace of Asagai's critique show her \emph{agency}.
\end{itemize}

\textbf{Paragraph 3: Language and Intellectual Awakening}
\begin{itemize}
\item Beneatha's vocabulary shifts: ``assimilationist'', ``mutilation'' (of hair), ``identity''.
\item She articulates the \emph{psychic cost} of assimilation: ``I am not interested in being someone's little episode in America.''
\item This is not teenage rebellion; it is \emph{theoretical engagement} with Fanon‑type ideas (Hansberry read Fanon).
\end{itemize}

\textbf{Paragraph 4: Connection to Wider Themes}
\begin{itemize}
\item Her search mirrors the family's search: \emph{What kind of life can Black people build in America?}
\item Mama's house in Clybourne Park = spatial integration; Beneatha's Afrocentrism = cultural integrity. Both are necessary.
\item Walter's later growth (Act 3) echoes Beneatha's: he reclaims dignity by rejecting Lindner, just as she rejects George.
\end{itemize}

\textbf{Conclusion:}
Beneatha is not a peripheral character; she is the play's \emph{conscience}. Hansberry uses her to insist that the American Dream cannot be bought at the price of self‑erasure. The tension in this scene foreshadows the family's ultimate triumph: they move to Clybourne Park \emph{with} their identity intact.
\end{exoresolu}

\courssub{Key Quotations for Revision}

\begin{aretenir}[Essential Quotations — Act 2 Scene 1]
\begin{itemize}
\item \textbf{Mama:} \emph{``There is always something left to love. And if you ain't learned that, you ain't learned nothing.''} (To Beneatha, after Walter's outburst) — \cle{Unconditional love as survival strategy}.
\item \textbf{Mama:} \emph{``We ain't never been that dead inside.''} — \cle{Dignity as non‑negotiable}.
\item \textbf{Beneatha:} \emph{``I'm looking for my identity!''} — \cle{The central quest of the play's youngest generation}.
\item \textbf{Asagai (to Beneatha):} \emph{``You are not fully alive... because you have no past.''} — \cle{History as prerequisite for agency}.
\item \textbf{Walter (to Travis):} \emph{``I'm going to make a transaction... that's going to change our lives.''} — \cle{Tragic irony: the transaction destroys him}.
\item \textbf{George:} \emph{``Let's face it, baby, your heritage is nothing but a bunch of raggedy‑assed spirituals and some grass huts!''} — \cle{Internalised racism masquerading as sophistication}.
\end{itemize}
\end{aretenir}

\courssub{Summary: Why This Scene Matters}

Act 2 Scene 1 is the \cle{fulcrum} of \emph{A Raisin in the Sun}. It accomplishes four things simultaneously:
\begin{enumerate}
\item \textbf{Advances the plot:} The house is bought; the money is distributed; the Clybourne Park conflict is seeded.
\item \textbf{Deepens character:} Walter's fragility and potential; Mama's strategic wisdom; Beneatha's intellectual courage.
\item \textbf{Stages the central ideological debate:} Assimilation vs. cultural recovery, acted out through Beneatha's body (hair, robe) and her suitors.
\item \textbf{Establishes the dramatic irony} that will make the coming loss (Act 2 Scene 2/3) and the final stand (Act 3) resonate with tragic and then heroic weight.
\end{enumerate}

As you prepare for the GCE, remember: examiners look for \emph{integration} of textual evidence, dramatic technique and thematic insight. Do not treat the ``American Dream'', ``cultural identity'' and ``money'' as separate topics — they are \emph{one dramatic complex} in this scene. The cheque is the catalyst; the family is the crucible; the dream is the product — or the casualty.

\begin{exercice}[Practice Questions for Act 2 Scene 1]
\begin{enumerate}
\item How does Hansberry use \emph{stage directions} (music, costume, lighting) to externalise Beneatha's internal conflict in the opening of Act 2 Scene 1?
\item ``Mama is the play's moral anchor.'' Discuss with reference to her handling of the insurance money in this scene.
\item Analyse the \emph{dramatic irony} in Walter's speech to Travis. How does it prepare the audience for the events of Act 2 Scene 2?
\item Compare and contrast George Murchison and Joseph Asagai as \emph{dramatic foils} for Beneatha's development.
\item ``Money is both the key and the trap.'' Explore this paradox in Act 2 Scene 1.
\end{enumerate}
\end{exercice}

\begin{corrige}[Exercise 1 — Stage Directions and Beneatha's Conflict]
\textbf{Key points to include:}
\begin{itemize}
\item \textbf{Nigerian folk music:} Non‑Western soundscape invades the Chicago apartment — auditory claim of African space.
\item \textbf{Robe and natural hair:} Visual rejection of Eurocentric beauty standards; the robe is a \emph{gift from Asagai}, making the political personal.
\item \textbf{``Serious, almost religious concentration'':} Elevates the dance from entertainment to ritual; identity recovery as sacred act.
\item \textbf{Contrast with George's entrance:} He brings the ``assimilationist'' world (suit, dismissive tone) into the same space — collision staged in real time.
\item \textbf{Lighting:} ``Gray morning light'' suggests the ambiguity of her position — neither fully African nor fully accepted American.
\end{itemize}
\textbf{Examiner's note:} Always quote stage directions exactly (or paraphrase precisely) and link to \emph{character psychology} and \emph{theme}.
\end{corrige}

\coursec{Section 7: A Raisin in the Sun Act 2 Scenes 2–3 — Loss, Racism and the Clybourne Park Decision}

In Section 6 we saw Act 2 Scene 1 open in celebration: Walter and Ruth dancing, Beneatha in Nigerian dress, the family briefly united around the promise of the insurance money. Hansberry deliberately builds that joy so that what follows will hurt more. Scenes 2 and 3 of Act 2 contain the play's double catastrophe and double climax: first the family loses everything Walter invested, then Walter must decide what kind of man he will be in front of his son. These scenes answer the question the whole play has been asking: when a dream is deferred, does it dry up like a raisin in the sun, or does it explode into dignity?

\courssub{Scene 2: Mrs Johnson's Visit — Foreshadowing and White Flight}

Scene 2 opens with a moment of fragile happiness: the family is packing, Beneatha teases Walter, and Mama has even bought a little plant to take to the new house. Into this warmth steps \cle{Mrs Johnson}, the Youngers' neighbour, a comic figure on the surface — nosy, envious, self-important — but carrying a dark message. She shows the family a newspaper report that another black family that moved into a white neighbourhood has been \emph{bombed}. She laughs it off with her catchphrase that she is a ``proud'' coloured person, but the warning is real: the Youngers' move to Clybourne Park may be answered with violence.

\begin{definition}[White Flight]
\cle{White flight} is the historical movement of white families out of a neighbourhood when black families move in, often encouraged by estate agents, residents' associations and discriminatory lending practices. In 1950s Chicago — the play's setting — black families who crossed the colour line faced hostile ``improvement associations'', boycotts, vandalism and sometimes bombings. Clybourne Park is Hansberry's fictionalised version of this reality: her own family had fought a famous court case (\emph{Hansberry v. Lee}, 1940) after moving into a white Chicago neighbourhood.
\end{definition}

Dramatically, Mrs Johnson's visit performs two functions. First, it is \cle{foreshadowing}: the audience is told, before Mr Lindner returns, that the threat to the family is not only financial but physical. Second, it exposes \emph{internalised racism}: Mrs Johnson mocks Beneatha's Africanism and Walter's ambitions, showing how oppression teaches the oppressed to police each other. Mama's quiet response — she does not argue, she simply continues packing — shows a dignity Mrs Johnson cannot touch.

\begin{savaistu}[Hansberry's personal history]
Lorraine Hansberry's father, Carl Hansberry, bought a house in a white neighbourhood in Chicago in 1937. The family faced a mob, legal challenges, and a Supreme Court case (\emph{Hansberry v. Lee}, 1940) that struck down a restrictive covenant. This autobiographical core gives the play its authentic tension: the Youngers' struggle is not invented — it is lived history.
\end{savaistu}

\courssub{Climax 1: The Loss — Willy's Betrayal}

The scene's mood collapses when Bobo arrives, hesitant and sweating, to deliver the news Walter cannot say himself: \cle{Willy Harris} never went to Springfield to get the liquor-store licence. He has run off with all the money — Walter's investment \emph{and} the portion Mama had entrusted to Walter for Beneatha's medical school fees. Walter's earlier boast that he would make a transaction ``that will change our lives forever'' returns as bitter irony.

\begin{definition}[Climax in Drama]
The \cle{climax} is the scene of greatest tension in a play, the point at which the protagonist must make an irreversible choice or suffer an irreversible event, after which the action turns towards its resolution. \emph{A Raisin in the Sun} is unusual in having \textbf{two} climaxes: the \emph{loss} (Act 2 Scene 2), where Walter is destroyed by trusting the wrong man, and the \emph{refusal} (Act 2 Scene 3), where he is redeemed by defying the wrong man.
\end{definition}

The family's reactions map the theme of the play. Beneatha's despair is ideological: her dream of healing and of Africa seems to die with the money, and she turns on Walter with contempt. Mama's grief is the deepest and most surprising — she does not rage about the money but remembers \cle{Big Walter}, her dead husband, who literally worked himself to death: she cries that she saw him ``killing himself'' day by day, and now the fruit of his labour has been thrown away in a single afternoon. The money was never just money; it was a man's life converted into cash. Walter, shattered, can only repeat Willy's name and his own foolishness.

\begin{attention}[Common mistake]
Do not write that Mama gave Walter the money ``for the liquor store''. She gave him \$6,500 with clear instructions: \$3,000 for Beneatha's schooling, the rest for him to manage. Walter disobeyed by handing \emph{everything} to Willy. The moral weight of the catastrophe depends on this disobedience — he gambled his sister's future, not only his own.
\end{attention}

\courssub{Asagai's Challenge: Idealism Tested by Catastrophe}

In the wreckage, \cle{Asagai} comes to say goodbye to Beneatha and finds her packing in despair. His long speech is the philosophical heart of the scene. He observes that it is strange that ``those who see the changes'' in the world — the dreamers, the reformers — ``are always the ones who are most disappointed''. He challenges Beneatha: was her dream of being a doctor dependent on her father's death-money? Then it was never fully hers. He asks her to choose between \emph{idealism} and \emph{cynicism}, and renews his invitation: come home to Nigeria, where there is real work of liberation to do.

This exchange matters for two reasons. First, it prevents the play from becoming a simple story about money: Asagai insists that the future is something you \emph{build}, not something you \emph{inherit}. Second, it gives Beneatha a genuine choice between two visions of black identity — George Murchison's assimilation into white American wealth, and Asagai's return to African roots — which she is still weighing as the play ends.

\begin{aretenir}[Asagai's key distinction]
Asagai separates \emph{idealism} (working for a future you may not see) from \emph{cynicism} (giving up because the present is hard). For the GCE essay, link this to the play's title: the raisin in the sun may dry up, but the \emph{seed} can still be planted.
\end{aretenir}

\courssub{Scene 3: Mr Lindner's Second Visit — Racism in a Polite Mask}

Scene 3 opens with Walter inventing a desperate plan: he has telephoned \cle{Mr Lindner}, the representative of the Clybourne Park Improvement Association, to accept the buy-out offer from Act 1 — the association will pay the Youngers \emph{not} to move in. Walter rehearses the performance he will give: he will be the grinning, shuffling, humble Negro that white people expect, ``Mistah Charlie's'' obedient boy. Mama and Beneatha watch this rehearsal with horror. Beneatha declares he is ``no brother of mine''; Mama answers with one of the play's most important lines — that the time to love someone most is when he is at his lowest and believes he deserves no love.

\begin{exemplebox}[Lindner's rhetoric: politeness as a weapon]
Lindner never insults the family. He speaks of ``community harmony'', says the residents are ``not bad people'', just people who want to protect their ``dream'' of the neighbourhood, and frames the buy-out as generous. This is \cle{structural racism} in its most effective costume: no slurs, no threats — only money, smiles and the assumption that black people simply do not belong. Hansberry's point is that polite racism is harder to fight than open hatred, because it makes the victim look unreasonable for refusing.
\end{exemplebox}

\begin{popfigure}
\begin{tikzpicture}[font=\small]
\draw[thick, popline, ->] (0,0) -- (10.5,0) node[below left=2pt and -40pt, popmuted]{};
\draw[thick, popline, ->] (0,0) -- (0,5.2) node[rotate=90, above left=18pt and -52pt, popmuted]{};
\node[popmuted, below] at (5.2,-0.15) {Dramatic time (Act 2)};
\node[popmuted, rotate=90] at (-0.55,2.6) {Tension};
\coordinate (A) at (0.4,1.2);
\coordinate (B) at (2.4,2.6);
\coordinate (C) at (4.6,4.6);
\coordinate (D) at (6.6,0.9);
\coordinate (E) at (8.6,4.8);
\coordinate (F) at (10.2,3.1);
\draw[very thick, popblue] (A) -- (B) -- (C) -- (D) -- (E) -- (F);
\fill[poporange] (B) circle (2.5pt);
\fill[popdark] (C) circle (3pt);
\fill[popgreen] (D) circle (2.5pt);
\fill[popdark] (E) circle (3pt);
\fill[poppurple] (F) circle (2.5pt);
\node[above, poporange, align=center, text width=2.2cm] at (B) {Mrs Johnson: bombing warning};
\node[above, popdark, align=center, text width=2.4cm] at (C) {CLIMAX 1: Willy steals the money};
\node[below, popgreen, align=center, text width=2.2cm] at (D) {Despair; Walter phones Lindner};
\node[above, popdark, align=center, text width=2.4cm] at (E) {CLIMAX 2: Walter refuses Lindner};
\node[right, poppurple, align=center, text width=2.2cm] at (F) {Resolution: the family moves out};
\end{tikzpicture}
\legende{Plot structure of Act 2 Scenes 2--3: rising tension, the double climax and the qualified resolution.}
\end{popfigure}

\begin{popfigure}
\begin{tikzpicture}[font=\small]
\node[draw=popblue, fill=popblueL, rounded corners, align=center, text width=4.6cm, minimum height=1.1cm] (mask) at (0,0) {\textbf{What Lindner says}\\ ``community harmony''};
\node[draw=popdark, fill=white, rounded corners, align=center, text width=4.6cm, minimum height=1.1cm] (real) at (7.2,0) {\textbf{What it means}\\ ``black families are not welcome''};
\draw[very thick, poporange, ->] (mask) -- (real) node[midway, above, popmuted]{masks};
\node[draw=popblue, fill=popblueL, rounded corners, align=center, text width=4.6cm, minimum height=1.1cm] (mask2) at (0,-2) {\textbf{What Lindner says}\\ ``a generous business offer''};
\node[draw=popdark, fill=white, rounded corners, align=center, text width=4.6cm, minimum height=1.1cm] (real2) at (7.2,-2) {\textbf{What it means}\\ ``we will pay to keep you out''};
\draw[very thick, poporange, ->] (mask2) -- (real2) node[midway, above, popmuted]{masks};
\node[draw=popblue, fill=popblueL, rounded corners, align=center, text width=4.6cm, minimum height=1.1cm] (mask3) at (0,-4) {\textbf{What Lindner says}\\ ``people can get awful worked up''};
\node[draw=popdark, fill=white, rounded corners, align=center, text width=4.6cm, minimum height=1.1cm] (real3) at (7.2,-4) {\textbf{What it means}\\ a veiled threat of violence};
\draw[very thick, poporange, ->] (mask3) -- (real3) node[midway, above, popmuted]{masks};
\end{tikzpicture}
\legende{Lindner's rhetoric: courteous language concealing structural racism and intimidation.}
\end{popfigure}

\courssub{Climax 2: The Decision — ``Brick by Brick''}

When Lindner arrives with his papers, Walter begins his humiliating performance — and then sees his young son \cle{Travis} watching him. Mama has insisted that Travis stay: Walter must make his son proud, or teach him shame. Looking at the boy, Walter cannot do it. He straightens up and tells Lindner, quietly and clearly, that the family will move into their house, because his father ``earned it for us brick by brick''. He adds that they are plain, hardworking people, and that his sister will be a doctor. Lindner leaves, defeated; the family erupts into motion, packing for the move.

\begin{aretenir}[Key quotation]
Mama's verdict on the scene is the play's emotional summary: Walter ``finally come into his manhood today, didn't he? Kind of like a rainbow after the rain.'' Manhood, for Hansberry, is not money or success — it is the courage to claim dignity in front of one's child, after failure.
\end{aretenir}

\begin{methode}[Analysing a climax in drama]
\begin{enumerate}
\item \textbf{Identify the choice} facing the character (Walter: take Lindner's money and stay, or refuse and move).
\item \textbf{State what is at stake} (money the family desperately needs versus self-respect, Travis's example, Big Walter's legacy).
\item \textbf{Describe the decision and its trigger} (the sight of Travis transforms Walter's performance into a refusal).
\item \textbf{Conclude on character and theme} (Walter matures from dreamer to responsible man; the play affirms dignity over material security).
\end{enumerate}
\end{methode}

\courssub{Resolution: A Qualified Hope}

The play ends with the family leaving the cramped apartment, poorer in cash but richer in unity and self-respect. Mama comes back for her little plant — the symbol of her stubborn, nurturing dream — and the stage empties. Hansberry offers \emph{hope, not certainty}: the Youngers are moving into a neighbourhood that has already sent Lindner and may send worse.

\begin{attention}[The ending is not a simple happy ending]
A frequent exam error is to call the conclusion ``happy''. The family has lost the investment money, Beneatha's schooling is uncertain, and Clybourne Park is hostile — Mrs Johnson's bombing story hangs over the final scene. The ending is a \cle{qualified hope}: the raisin in the sun refuses to dry up, but the sun is still hot. Always mention both the victory (dignity, unity) and the danger (racism, poverty).
\end{attention}

\begin{exoresolu}[Essay question: the double climax]
\textbf{Question:} ``\emph{A Raisin in the Sun} has two climaxes: one destroys Walter, the other saves him.'' Discuss with close reference to Act 2 Scenes 2 and 3.
\tcblower
\textbf{Plan and model paragraph.} Introduction: define climax; state the thesis that the two climaxes mirror each other — both involve Walter, money and another man's offer, but his responses are opposite. Body 1 (Climax 1): Willy's theft; Walter's passivity and despair; the loss of Beneatha's fees; Mama's memory of Big Walter; Walter destroyed because he trusted the wrong man. Body 2 (Climax 2): Lindner's polite offer; Walter's planned ``grinning'' performance; Travis's presence as the turning point; the refusal — ``my father... earned it for us brick by brick''; Walter redeemed because he defied the wrong man. Conclusion: the symmetry shows that Hansberry measures manhood not by wealth but by moral choice; the ending is hopeful but qualified by the racism awaiting the family.

\textbf{Model paragraph (Climax 2):} In Scene 3, Hansberry replays the first climax with its terms reversed. Once again Walter stands between his family and a man offering a shortcut to money; but where he had passively trusted Willy, he now actively confronts Lindner. The stage presence of Travis is decisive: unable to perform the ``grinning'' humility he had rehearsed before his son's eyes, Walter instead claims the house as his father's legacy, earned ``brick by brick''. The alliteration and the image of manual labour root his dignity in honest work rather than speculation. Mama's comment that he has ``come into his manhood... like a rainbow after the rain'' confirms that this second climax completes the moral education the first one began.
\end{exoresolu}

\begin{exercice}[Comparing the two climaxes]
In 250--300 words, compare Walter's behaviour in the two climaxes of Act 2. Show how the first (trusting Willy) destroys him while the second (defying Lindner) redeems him, and explain what this reveals about Hansberry's idea of manhood.
\end{exercice}

\begin{corrige}[Exercise 1]
A strong answer notes the deliberate symmetry: both scenes involve money, a white or absent man, and a decision about the family's future. In Scene 2 Walter is passive and credulous — he has handed everything, including Beneatha's fees, to Willy, and collapses into despair when the theft is revealed; Mama's grief for Big Walter shows he has betrayed his father's legacy. In Scene 3 Walter is active and lucid: facing Lindner's polite racism, and watched by Travis, he refuses the buy-out and claims the house his father earned ``brick by brick''. The contrast reveals Hansberry's definition of manhood: not financial success but moral courage, responsibility to family and the transmission of dignity to the next generation. The best answers add that the redemption is qualified — the family still faces hostility in Clybourne Park — so the play ends in hope, not triumph.
\end{corrige}

\newpage

\coursec{Integration activity}

\courssub{Aims of the Activity}
\begin{aretenir}[Aims of the Activity]
\begin{itemize}
    \item To consolidate plot analysis, character study and thematic comparison across the three set texts: \emph{Nineteen Eighty-Four} (Parts 2 and 3), \emph{The Canterbury Tales} (General Prologue portraits and \emph{The Merchant's Tale}) and \emph{A Raisin in the Sun} (Act 2, Scenes 1--3).
    \item To develop collaborative critical thinking by organising the class into expert groups that must defend a thesis using only textual evidence.
    \item To practise the synthesis skill required for the GCE Advanced Level comparative essay: evaluating how storyline functions as a vehicle for social criticism.
    \item To connect literary analysis to the Cameroonian classroom context, encouraging students to see literature as a tool for interrogating their own society.
\end{itemize}
\end{aretenir}

\courssub{Situation: The Plot Detectives' Symposium}
\begin{experience}[The Plot Detectives' Symposium]
The Literature Club of \textbf{Government Bilingual High School, Yaound\'e} is preparing for the upcoming GCE Advanced Level examination. The club patron, Mr.~Fonge, has designed a \emph{Plot Detectives' Symposium} to help Upper Sixth students move beyond isolated text study and engage in inter‑textual dialogue. The symposium will take place in the school library over two consecutive literature periods (80 minutes total). 

The class of 45 students is divided into three expert groups of 15 each. Each group receives a \textbf{Case File} containing the relevant plot outlines, character lists and thematic prompts for their assigned text. The groups have 30 minutes to prepare a five‑minute oral presentation supported by a \textbf{Plot‑Structure Poster} (A1 paper, markers provided). The poster must display: (1) a timeline of key plot events, (2) a tension graph (rising action, climax, falling action), and (3) three annotated quotations (or close paraphrases) that substantiate the group's argument.

\textbf{Group A (Orwell)} must prove, using only plot events from \emph{Nineteen Eighty-Four} Parts 2 and 3, that \emph{``the Party wins because it destroys private life''}. They may refer to Winston's affair with Julia, the betrayal in Room 101, and the final scene in the Chestnut Tree Café.

\textbf{Group B (Chaucer)} must rank the studied pilgrims — Cook, Manciple, Shipman, Physician, Wife of Bath, Parson, Plowman, Summoner, Pardoner, Host (Harry Bailey) — from most to least admirable. The ranking must be defended with textual evidence from the General Prologue portraits and, where relevant, from \emph{The Merchant's Tale} (January's delusion, the marriage debate between Placebo and Justinus).

\textbf{Group C (Hansberry)} must argue whether the Younger family's ending in Act 2 constitutes a victory or a defeat. They must cite the two climaxes of Act 2: (i) the loss of the insurance money (Scene 2) and (ii) the confrontation with Karl Lindner in Clybourne Park (Scene 3).

After the three presentations, the whole class votes on the synthesis question: \emph{``Which writer uses storyline most powerfully to criticise society?''} The teacher then models a comparative essay plan that students can adapt for the GCE examination.
\end{experience}

\courssub{Tasks}
\begin{methode}[Symposium Procedure]
\begin{enumerate}
    \item \textbf{Preparation (30 minutes)}: In your expert group, read the Case File, discuss the thesis, select three key plot moments, and design the Plot‑Structure Poster. Assign roles: \emph{Timekeeper}, \emph{Scribe}, \emph{Presenter}, \emph{Evidence Checker}.
    \item \textbf{Presentation (5 minutes per group)}: Deliver your case. The Presenter speaks; the Evidence Checker displays the poster. The class takes notes using the \emph{Comparative Observation Grid} (provided by the teacher).
    \item \textbf{Class Vote and Discussion (10 minutes)}: Each student casts a secret ballot for the synthesis question. The teacher tallies votes and leads a brief discussion on why different groups may have chosen different writers.
    \item \textbf{Teacher Modelling (15 minutes)}: The teacher displays a model comparative essay plan on the board, highlighting: (a) introduction with a clear thesis, (b) three body paragraphs each focusing on one text's plot structure as social critique, (c) a synthesis paragraph linking the three, (d) conclusion evaluating relative effectiveness.
    \item \textbf{Individual Write‑Up (Homework)}: Using the essay plan, each student writes a 500‑word comparative essay for submission in the next lesson. The essay must reference at least two plot events per text and use appropriate literary terminology (e.g., \emph{peripeteia}, \emph{anagnorisis}, \emph{exposition}, \emph{denouement}).
\end{enumerate}
\end{methode}

\courssub{Model Answer}
\begin{exemplebox}[Teacher's Model Comparative Essay Plan]
\begin{center}
\begin{tabularx}{\linewidth}{|l|X|}
\hline
\rowcolor{popblueL}
\textbf{Section} & \textbf{Content and Key Points} \\
\hline
\textbf{Title} & \emph{Plot as Scalpel: How Orwell, Chaucer and Hansberry Use Storyline to Dissect Their Societies} \\
\hline
\textbf{Introduction} & 
\begin{itemize}
    \item Hook: A society's deepest flaws are often exposed not by sermons but by the stories it tells about itself.
    \item Context: Three canonical texts — \emph{Nineteen Eighty-Four} (1949), \emph{The Canterbury Tales} (late 14th century), \emph{A Raisin in the Sun} (1959) — each uses plot to critique totalitarianism, ecclesiastical corruption, and racial capitalism respectively.
    \item Thesis: While all three writers employ storyline to criticise society, Orwell's tightly controlled narrative arc in Parts 2 and 3 most powerfully demonstrates the eradication of private life as the mechanism of totalitarian victory; Chaucer's episodic portrait gallery offers a panoramic moral audit; Hansberry's domestic tragedy turns a single family's decision into a national indictment.
\end{itemize} \\
\hline
\textbf{Body Paragraph 1: Orwell — The Architecture of Betrayal} & 
\begin{itemize}
    \item \textbf{Plot structure}: Part 2 builds a fragile private world (Winston–Julia affair, the room above Charrington's shop) — a rising action of rebellion. Part 3 systematically dismantles it: arrest, Ministry of Love, Room 101, final capitulation in the Chestnut Tree Café.
    \item \textbf{Key events as evidence}: 
    \begin{enumerate}
        \item The purchase of the paperweight (symbol of a past the Party cannot touch) — later smashed during arrest.
        \item The betrayal in Room 101 where Winston screams ``Do it to Julia!'' — the destruction of love as the final private bond.
        \item The closing scene: Winston's alcoholic resignation, ``He loved Big Brother.'' The plot's denouement confirms the Party's total victory.
    \end{enumerate}
    \item \textbf{Critical commentary}: The linear, inevitable trajectory mirrors the Party's logic; there is no subplot to offer relief. The storyline itself becomes an instrument of oppression.
\end{itemize} \\
\hline
\textbf{Body Paragraph 2: Chaucer — The Mosaic of Moral Portraiture} & 
\begin{itemize}
    \item \textbf{Plot structure}: The General Prologue is not a single narrative but a series of static portraits; \emph{The Merchant's Tale} provides a narrative exemplum. The ``plot'' is the pilgrimage frame, but the social critique lies in the juxtaposition of characters.
    \item \textbf{Ranking evidence (Group B's task)}: 
    \begin{enumerate}
        \item \textbf{Parson} and \textbf{Plowman} — ideal Christian charity (``He was a shepherd and no mercenary'').
        \item \textbf{Wife of Bath} — complex, vital, challenges patriarchy through experience.
        \item \textbf{Shipman}, \textbf{Cook}, \textbf{Manciple} — skilled but morally ambiguous (the Cook's ulcer, the Manciple's cunning).
        \item \textbf{Physician} — love of gold over cure.
        \item \textbf{Summoner} and \textbf{Pardoner} — grotesque corruption (``His wallet lay before him in his lap / Bretful of pardons come from Rome all hot'').
        \item \textbf{Host (Harry Bailey)} — genial manipulator, controls the storytelling contest.
    \end{enumerate}
    \item \textbf{Critical commentary}: The ``plot'' of the pilgrimage serves as a social microscope; the critique is cumulative, not teleological. The Merchant's Tale (January's delusion, the Placebo/Justinus debate) extends the critique to marriage as a social contract.
\end{itemize} \\
\hline
\textbf{Body Paragraph 3: Hansberry — Domestic Tragedy as Public Protest} & 
\begin{itemize}
    \item \textbf{Plot structure}: Act 2 contains two climaxes. First, the loss of the \$10,000 insurance money (Scene 2) — a peripeteia that shatters Walter's dream of economic autonomy. Second, the confrontation with Karl Lindner (Scene 3) — an anagnorisis where the family reclaims dignity by refusing the buyout.
    \item \textbf{Key events as evidence}: 
    \begin{enumerate}
        \item Walter's confession to Mama: ``That money is made out of my father's flesh.'' The loss exposes the fragility of the Black American Dream.
        \item Beneatha's disillusionment with assimilation (Asagai's challenge).
        \item The final scene: Walter tells Lindner, ``We have decided to move into our house because my father — my father — he earned it.'' The storyline transforms a private family decision into a public act of resistance.
    \end{enumerate}
    \item \textbf{Critical commentary}: The plot's tight domestic scope magnifies the societal forces (racism, capitalism, generational conflict). The victory is moral, not material — a nuanced critique.
\end{itemize} \\
\hline
\textbf{Synthesis Paragraph} & 
\begin{itemize}
    \item Compare the three plot strategies: Orwell's \emph{inexorable linear destruction} vs. Chaucer's \emph{static comparative mosaic} vs. Hansberry's \emph{double climax of loss and reclamation}.
    \item Argue that Orwell's storyline is most powerful \emph{as a warning} because its structure leaves no escape; the reader experiences the same suffocation as Winston. Chaucer's method is more diagnostic than prophetic. Hansberry's offers hope but within a limited sphere.
    \item Link to Cameroonian context: Just as Orwell's Party destroys private life to secure power, students can reflect on how contemporary societies (including Cameroon) may erode personal freedoms through surveillance, corruption, or ethnic division. Literature equips us to recognise these patterns.
\end{itemize} \\
\hline
\textbf{Conclusion} & 
\begin{itemize}
    \item Restate thesis in fresh words.
    \item Final thought: The storyline is not merely a sequence of events; it is the writer's scalpel. Orwell cuts deepest because his plot admits no redemption — a deliberate artistic choice that forces the reader to confront the logic of totalitarianism.
    \item Call to action for GCE candidates: In the examination, always anchor thematic claims in specific plot moments; use terms like \emph{exposition}, \emph{rising action}, \emph{climax}, \emph{denouement} to show structural awareness.
\end{itemize} \\
\hline
\end{tabularx}
\end{center}
\end{exemplebox}

\courssub{Sample Opening Paragraph}
\begin{exemplebox}[Sample Opening Paragraph (for student reference)]
\begin{quote}
In \emph{Nineteen Eighty-Four}, George Orwell constructs a plot that functions as a machine for the annihilation of the private self. The narrative arc of Parts 2 and 3 moves from the tentative creation of a secret world — Winston's diary, his affair with Julia, the rented room — to its systematic demolition in the Ministry of Love. Each stage of the plot mirrors the Party's doctrine: control the past, control the present, control the future. By contrast, Geoffrey Chaucer's \emph{General Prologue} offers no single narrative arc but a gallery of portraits that collectively audit the moral health of 14th-century English society. Lorraine Hansberry's \emph{A Raisin in the Sun} compresses social critique into the domestic crucible of the Younger family, where two climaxes in Act 2 — the loss of the insurance money and the rejection of Karl Lindner's bribe — transform a private tragedy into a public declaration of dignity. While all three writers use storyline to criticise society, Orwell's relentless linear plot most powerfully demonstrates how the destruction of private life secures totalitarian victory.
\end{quote}
\end{exemplebox}

\courssub{Examiner's Tips for the GCE}
\begin{attention}[Examiner's Tips]
\begin{itemize}
    \item Always quote or closely paraphrase; vague references earn no marks.
    \item Use the metalanguage of plot: \emph{exposition, inciting incident, rising action, climax, falling action, denouement}.
    \item In comparative essays, devote one paragraph to each text, then a synthesis paragraph — do not mix texts within a single paragraph.
    \item Contextualise briefly (e.g., post‑war Britain, medieval estates satire, Civil Rights era USA) but keep focus on the plot's critical function.
    \item Proofread for spelling of character names: \emph{Winston Smith, Julia, O'Brien, Big Brother; Chaucer, Harry Bailey, January, Placebo, Justinus; Walter Lee, Mama, Beneatha, Karl Lindner}.
\end{itemize}
\end{attention}

\newpage

\coursec{Methods and worked examples}

\courssub{Skill 1: Analysing Plot Development and Narrative Turning Points}

Before any analysis can begin, the candidate must be able to \emph{see} a narrative as a constructed whole: a sequence of episodes chosen and ordered by the author to produce an effect. The GCE questions on \emph{Nineteen Eighty-Four}, \emph{The Canterbury Tales} and \emph{A Raisin in the Sun} very often ask you to "discuss the significance" of one chapter or scene. This is a question about \cle{plot architecture}, not about storytelling.

\begin{definition}[Key Narrative Terms]
\begin{itemize}
\item \textbf{Exposition:} the opening material that establishes setting, characters and situation (e.g. Part 1 of \emph{Nineteen Eighty-Four}; Act 1 Scene 1 of \emph{A Raisin in the Sun}, with its tired living-room and the expected cheque).
\item \textbf{Rising action:} the series of complications that intensify the central conflict (Winston and Julia's affair; the Youngers' quarrels over the money).
\item \textbf{Climax / turning point (peripeteia):} the episode in which the fortune of the protagonist is decisively reversed (the arrest in Part 2 Chapter 10; the loss of the money in Act 2 Scene 2).
\item \textbf{Falling action and resolution:} the consequences of the climax, leading to the final state of affairs (Part 3 of \emph{Nineteen Eighty-Four}; the move to Clybourne Park).
\item \textbf{Foreshadowing:} a detail that anticipates a later event (the rhyme "here comes a chopper to chop off your head"; the hidden telescreen behind the picture).
\end{itemize}
\end{definition}

\begin{methode}[How to Analyse a Plot Episode]
When the GCE asks you to "discuss the significance" of a chapter or scene (e.g. \emph{Nineteen Eighty-Four} Part 2 Chapter 9, or \emph{A Raisin in the Sun} Act 2 Scene 2), follow these steps:
\begin{enumerate}
\item \textbf{Locate} the episode precisely in the whole work (part, chapter, act, scene) and state what comes immediately before and after it.
\item \textbf{Summarise} the events in two or three sentences only — never retell the whole story.
\item \textbf{Identify the narrative function}: is it exposition, rising action, a turning point (climax), or falling action? Justify your label with one concrete detail.
\item \textbf{Link to character}: show what the episode reveals or changes in the protagonist (e.g. Winston's fatalism, Walter Lee's manhood).
\item \textbf{Link to theme}: connect the episode to at least one major theme (love as rebellion, the American Dream, racism).
\item \textbf{Conclude} by showing how the episode prepares what follows (foreshadowing, irony).
\end{enumerate}
\textbf{Reflex:} always ask yourself \emph{"What changes because of this scene?"} If nothing changes, it is probably exposition; if everything changes, it is a climax.
\end{methode}

\begin{exoresolu}[Worked Example: The Loss of the Money in \emph{A Raisin in the Sun}]
\textbf{Statement:} "Act 2 Scene 2 of \emph{A Raisin in the Sun}, in which Walter Lee loses the insurance money, is the first climax of the play." Discuss this view, showing how Hansberry constructs the scene as a turning point. (GCE-style, 25 marks)
\tcblower
\textbf{Step 1 — Location.} The scene falls near the end of Act 2, after Mama has entrusted Walter with the remaining \$6,500 (with instructions to bank \$3,000 for Beneatha's schooling and keep \$3,500 for himself) and immediately before the expected second visit of Mr Lindner and the move to Clybourne Park in Act 2 Scene 3 and Act 3.

\textbf{Step 2 — Brief summary.} Bobo arrives, agitated, to announce that Willy Harris has disappeared with all the money Walter invested in the liquor-store project — including Beneatha's \$3,000. Walter collapses; Mama realises that her husband's insurance money, the fruit of his whole working life, is gone.

\textbf{Step 3 — Narrative function.} This is unmistakably a \cle{climax} (the first of two, the second being Walter's confrontation with Lindner in Act 3). Everything changes: the family's financial hope, Walter's authority as "head of the family," and Beneatha's dream of medicine are destroyed in a single entrance. Hansberry marks the moment theatrically: Bobo's stammering hesitation delays the blow, and the stage directions show Walter's physical collapse, an external sign of inner ruin.

\textbf{Step 4 — Character.} The scene exposes Walter's tragic flaw: his equation of \cle{manhood with money}. He entrusted the cash to Willy without any written guarantee, because he wanted to feel like a businessman overnight. Yet Hansberry prevents us from simply condemning him: his desperation springs from years of humiliation as a chauffeur, "a servant" in his own eyes.

\textbf{Step 5 — Theme.} The loss dramatises the fragility of the \cle{American Dream} for Black families: the insurance money — the only capital the Youngers will ever possess — vanishes into a fraudulent scheme, echoing how African-Americans were systematically excluded from legitimate business and credit. It also tests \cle{family solidarity}: Mama's grief is not for the money but for her husband's wasted life.

\textbf{Step 6 — Forward link.} The loss is the necessary precondition for the true climax: only because the family is ruined does Lindner's buy-out offer become tempting, and only then can Walter's final refusal prove his moral growth. The scene is therefore both an ending (of illusions) and a trigger (of redemption).

\textbf{Conclusion.} The statement is justified: Hansberry constructs the scene as a classical peripeteia — a reversal of fortune — which strips the characters of material hope so that the play's real question, that of dignity, can be posed.
\end{exoresolu}

\begin{exemplebox}[Applying the Same Method to \emph{Nineteen Eighty-Four} Part 2 Chapter 9]
Chapter 9 is the long central chapter in which Winston reads extracts from Goldstein's \emph{Theory and Practice of Oligarchical Collectivism} in the room above Charrington's shop, with Julia asleep beside him. Its narrative function is double. Structurally, it is the \emph{calm before the catastrophe}: the lovers are at their safest and most united, and the very next chapter brings the arrest. Thematically, it supplies the \cle{intellectual explanation} of Oceania — perpetual war, the mutability of the past, the Inner Party's monopoly of power — so that the reader understands the machine that is about to crush the lovers. Note the irony: Winston reads a book he believes to be the truth of the resistance, yet Part 3 reveals that the book itself was written by a Party committee including O'Brien. Even the rebels' "truth" is manufactured by the enemy.
\end{exemplebox}

\courssub{Skill 2: Analysing Character and Character Development}

A character in a set text is not a real person but a \emph{construction}: the sum of what the character says, does, and what others say about him or her, filtered through the author's techniques (narrative voice, irony, stage directions). The examiner rewards candidates who analyse the \emph{construction}, not those who describe the person.

\begin{methode}[How to Write a Character Analysis Paragraph]
\begin{enumerate}
\item \textbf{Claim:} open with a precise, arguable statement about the character (not "Winston is the hero" but "Winston's rebellion is doomed from the start because it is nostalgic rather than political").
\item \textbf{Contextualise:} name the part/chapter/scene where your evidence occurs.
\item \textbf{Evidence:} quote a short phrase (three to eight words) or paraphrase a precise action.
\item \textbf{Analyse the technique:} comment on \emph{how} the author creates the effect — diction, imagery, irony, stage directions, symbolism.
\item \textbf{Evaluate:} link back to the claim and to a theme; show development if the character changes across the work.
\end{enumerate}
\textbf{Reflex:} one paragraph = one idea + one quotation + one technique + one thematic link.
\end{methode}

\begin{exoresolu}[Worked Example: O'Brien in Part 3 of \emph{Nineteen Eighty-Four}]
\textbf{Statement:} "O'Brien is not merely a torturer; he is the true believer who completes the Party's destruction of Winston's self." Analyse Orwell's presentation of O'Brien in Part 3, Chapters 1--4.
\tcblower
\textbf{Claim and context.} In the Ministry of Love (Part 3), Orwell transforms O'Brien from the hoped-for ally of Part 2 into Winston's interrogator, teacher and — most disturbingly — intimate. The presentation is built on a deliberate \cle{irony}: the man Winston trusted becomes the instrument of his annihilation.

\textbf{Evidence and technique.} First, Orwell gives O'Brien the manner of a \emph{doctor} or \emph{priest} rather than a thug: he speaks patiently, almost tenderly, while ordering the dial on the pain machine to be turned up. This contrast between gentle tone and brutal action embodies the Party's perversion of care — Winston is treated as a "patient" whose "defective" thinking must be healed. Secondly, O'Brien articulates the Party's philosophy with chilling clarity: power is not a means but an end, and reality exists only in the collective mind of the Party. Orwell thus makes him the \cle{mouthpiece of totalitarian logic}, forcing the reader to confront the ideology directly rather than through Winston's confused resistance. Thirdly, the relationship is intimate: O'Brien tells Winston that he has always known him, and Winston feels that O'Brien's mind contains his own. This perverse closeness — torturer as confessor — shows that the Party does not merely kill rebels; it \emph{understands} and \emph{absorbs} them.

\textbf{Development and evaluation.} Across Chapters 1 to 4, O'Brien moves Winston through the three stages of "reintegration" that he himself names: learning (the beatings), understanding (the philosophical dialogues), and acceptance (the final betrayal of Julia in Room 101). Each stage strips away one layer of selfhood: bodily dignity, intellectual certainty, and finally love — the last private loyalty. By the end, Winston's capitulation is not external but internal: he genuinely loves Big Brother.

\textbf{Conclusion.} O'Brien is therefore far more than a villain: he is the novel's dark centre, the proof that totalitarianism's ultimate weapon is not fear but the seduction of the rebel's own intelligence. Orwell's presentation warns that the destruction of the self is completed not by strangers but by those who know us best.
\end{exoresolu}

\begin{exemplebox}[Short Model: Winston in Part 2 Chapters 7--8]
In Part 2 Chapter 7, Winston's memory of his mother defines his moral standard: her love was disinterested loyalty in a world where loyalty no longer exists, and he concludes that the proles, who have kept such feelings, are the only hope — "If there is hope, it lies in the proles." In Chapter 8, however, his visit to O'Brien shows the limits of his rebellion: he and Julia swear to commit murder, sabotage and even to throw sulphuric acid in a child's face for the Brotherhood, but refuse to promise to separate. Orwell's technique is the \cle{loaded exception}: the one thing they will not sacrifice is precisely the thing Room 101 will take. Character analysis here becomes plot analysis — Winston's deepest value is his point of maximum vulnerability.
\end{exemplebox}

\courssub{Skill 3: Close Reading of a Chaucerian Portrait}

The \emph{General Prologue} presents the pilgrims through the medieval theory of \cle{estates}: society was imagined as ordered groups (those who pray, those who fight, those who work), each with expected duties. Chaucer's art is to let concrete detail imply moral judgement, so the reader must learn to \emph{decode} detail.

\begin{methode}[How to Analyse a Portrait in the \emph{General Prologue}]
\begin{enumerate}
\item \textbf{Identify the estate:} state which social category the pilgrim belongs to (skilled worker, trader, professional, clergy) and what medieval society expected of that estate.
\item \textbf{List the details Chaucer selects:} clothing, physical features, tools, speech, habits, money. Chaucer characterises through \cle{concrete detail}, not abstract judgement.
\item \textbf{Detect the tone:} is the portrait admiring (Parson, Plowman), neutral-comic (Cook), or satirical (Summoner, Pardoner)? Look for irony — praise that is really blame.
\item \textbf{Decode the satire:} ask what vice the detail implies (the Cook's running sore suggests dubious hygiene; the Physician's love of gold suggests greed).
\item \textbf{Generalise:} conclude on what the portrait reveals about Chaucer's view of his society — corruption of the Church, dignity of honest labour, etc.
\end{enumerate}
\textbf{Reflex:} in Chaucer, \emph{every physical detail is a moral sign}. Never quote a detail without interpreting it.
\end{methode}

\begin{exemplebox}[The Skilled Workers and Traders: Cook, Manciple, Shipman]
\begin{itemize}
\item \textbf{The Cook} is a master of his craft — he can roast, boil, broil and fry, and knows his London ale — but Chaucer plants one devastating detail: the open sore (\emph{mormal}) on his shin. The comic portrait curdles: the reader cannot enjoy the description of his dishes once hygiene is in doubt. Tone: comic with a satirical sting.
\item \textbf{The Manciple} (steward of a lawyers' college) is illiterate, yet Chaucer ironically notes that he outwits the thirty learned lawyers he serves: he buys provisions shrewdly and always comes out ahead in his accounts. The portrait satirises the gap between \emph{book learning} and \emph{practical cunning} — and hints at dishonesty in his "tally."
\item \textbf{The Shipman (Skipper)} from Dartmouth is a superb seaman who knows every coast from Gotland to Cape Finisterre — but he steals wine from the merchants' cargo he carries, and in sea fights he sends his prisoners home "by water" (they are drowned). Chaucer's praise of his skill is genuine; the moral judgement is left to the reader. This is \cle{ironic neutrality}: the narrator admires the craftsman and ignores the criminal.
\end{itemize}
\textbf{Generalisation:} with the professional class, Chaucer's satire is mild and amused; he reserves his sharpest weapons for the corrupt clergy.
\end{exemplebox}

\begin{exoresolu}[Worked Example: The Pardoner and the Summoner]
\textbf{Statement:} "In the portraits of the Summoner and the Pardoner, Chaucer's satire of the corrupt Church is at its sharpest." Discuss with close reference to the \emph{General Prologue}.
\tcblower
\textbf{Estate and expectation.} The Summoner is an officer of the ecclesiastical court who cites sinners — especially for sexual offences — before the archdeacon; the Pardoner sells papal indulgences (remissions of penance). Both offices were notorious in the fourteenth century for extortion, and Chaucer exploits this expectation fully.

\textbf{Selected details and interpretation.} Chaucer begins with the body. The Summoner's face is "fire-red," covered with pustules that no ointment can cure — a detail medieval readers would read as the outward sign of inner corruption, possibly of lechery or disease. He frightens children, drinks wine until he is beside himself, and then speaks Latin only when drunk — a comic proof that his learning is a thin veneer. Crucially, Chaucer tells us he will let a sinner off "for a quart of wine": the instrument of Church discipline is himself for sale.

The Pardoner is presented through his \cle{professional tricks}: he carries a pillowcase which he claims is Our Lady's veil, and a jar of "pigges bones" which he sells as saints' relics to poor parsons, making more money in a day than the parson earns in two months. Chaucer's narrator adds, with devastating irony, that in church he sings the offertory beautifully — because a good song means a full collection plate. His wallet lies before him in his lap, "brimful of pardons come from Rome all hot": the comic exaggeration ("all hot") mocks the traffic in grace as if pardons were fresh loaves.

\textbf{Tone and satire.} The tone is not open denunciation but \cle{ironic admiration}: the narrator calls the Pardoner a "noble ecclesiaste" and the Summoner a "gentil harlot" — a perfect oxymoron that fuses false gentility with vice. Chaucer lets the rogues condemn themselves through their own practices.

\textbf{Generalisation.} Together the two portraits complete Chaucer's anatomy of the corrupt Church: the Summoner corrupts its \emph{justice}, the Pardoner its \emph{grace}. Set against the Parson and the Plowman — the ideal Christians of the same Prologue — they show that Chaucer's target is not religion but its mercenary betrayal. The satire is "at its sharpest" precisely because it is delivered with a smile.
\end{exoresolu}

\begin{exemplebox}[The Contrasting Portraits: Physician, Wife of Bath, Parson, Plowman, Host]
\begin{itemize}
\item \textbf{The Doctor (Physician)} is a perfect practitioner of medicine and astronomy (astrology), "grounded in astronomye" — but the portrait closes on the famous sting: "gold in physic is a cordial" (gold was believed a medicine), therefore he loved gold in special. He also made money during the plague and reads the Bible little. Skill is praised; greed and worldliness are implied.
\item \textbf{The Wife of Bath} is a skilled cloth-maker, gap-toothed, boldly dressed in scarlet stockings, and five times widowed — "housbondes at chirche dore she hadde fyve." Chaucer presents her with amused admiration: she embodies female energy, experience and the claim to mastery in marriage, which her own Prologue and Tale will develop. She is the Prologue's great study of \cle{women and marriage}.
\item \textbf{The Parson} is the ideal churchman: poor in goods but "rich in holy thought and work," patient in adversity, refusing to curse his parishioners for unpaid tithes, and practising before he preaches — he first "wroghte, and afterward he taughte." He visits his flock on foot, in rain and thunder. Chaucer's tone here is wholly serious: this is the standard by which the Summoner and Pardoner are condemned.
\item \textbf{The Plowman}, the Parson's brother, is the ideal layman: he works hard, lives in peace and charity, loves God and his neighbour, and pays his tithes faithfully. Honest labour is given full dignity.
\item \textbf{The Host, Harry Bailey}, is not satirised at all: a large, merry, masterful innkeeper, he proposes the tale-telling contest and appoints himself judge and guide. Structurally, he is the \cle{engine of the whole work} — without him there is no \emph{Canterbury Tales}.
\end{itemize}
\textbf{Generalisation:} the pairing of ideal (Parson, Plowman) and corrupt (Summoner, Pardoner) church figures is deliberate: Chaucer's satire works by \emph{contrast}, not by blanket condemnation.
\end{exemplebox}

\courssub{Skill 4: Analysing Theme and Author's Message}

A \cle{theme} is not a topic ("love," "racism") but a \emph{statement the work makes about the topic} ("love is a political act"; "the American Dream is both necessary and poisoned"). The examiner expects you to formulate the theme as an argument and to trace it through the work.

\begin{methode}[How to Build a Theme Essay]
\begin{enumerate}
\item \textbf{Define the theme} in one sentence as the author treats it (e.g. "In \emph{Nineteen Eighty-Four}, love is not a private feeling but a political act of rebellion").
\item \textbf{Trace the theme} through at least three moments of the work, in chronological order (beginning, middle, end).
\item \textbf{For each moment,} give one precise piece of evidence and one comment on technique.
\item \textbf{Show evolution or contrast:} how does the theme change, deepen or get destroyed? (Winston and Julia's love is born, flourishes above Charrington's shop, and is murdered in the Ministry of Love.)
\item \textbf{Conclude with the author's purpose:} what warning, celebration or question does the treatment of the theme leave with the reader?
\end{enumerate}
\textbf{Reflex:} a theme essay is a \emph{journey}, not a list. Use connectives of progression: "at first…", "this deepens when…", "finally…".
\end{methode}

\begin{exoresolu}[Worked Example: Love as Rebellion in \emph{Nineteen Eighty-Four} Part 2]
\textbf{Statement:} Trace the theme of love as an act of rebellion in Part 2 of \emph{Nineteen Eighty-Four} (Chapters 1--9) and show how Orwell prepares its destruction.
\tcblower
\textbf{Definition.} In Oceania, where the Party seeks to abolish private loyalty and even to make the sexual act a joyless "duty to the Party," Winston and Julia's affair is from the outset political: their embrace is, in the novel's own terms, "a blow struck against the Party."

\textbf{Stage 1 — Birth of the rebellion (Chapters 1--2).} Julia's note, "I love you," shatters Winston's isolation. Orwell stresses the risk: their first meeting in the crowd at Victory Square is choreographed like a military operation. The theme is born in fear — love and danger are inseparable.

\textbf{Stage 2 — The private world (Chapters 3--5).} The rented room above Mr Charrington's shop becomes a counter-state: the paperweight, the real coffee, the sugar, the make-up Julia puts on are all small restorations of beauty and the past. Orwell's technique is symbolic: the coral inside the glass paperweight figures their fragile, enclosed world — beautiful, anachronistic, and easily smashed. Winston's reflections on the proles (Chapter 7) widen the theme: genuine human feeling survives only among the proles, whom the Party does not bother to indoctrinate. His memory of his mother — whose love was an end in itself, a matter of loyalty — defines what the Party has destroyed.

\textbf{Stage 3 — The political commitment (Chapters 8--9).} The visit to O'Brien and the reading of Goldstein's book transform private love into organised conspiracy. Yet Orwell sows the seeds of destruction here: the lovers swear to do anything for the Brotherhood \emph{except} to stop loving each other — the very exception Room 101 will exploit. Chapter 9's long extracts from Goldstein's book also slow the love plot deliberately, suggesting that private happiness cannot survive while the system is unchallenged.

\textbf{Preparation of the destruction.} Foreshadowing accumulates: the picture of St Clement's Church hides the telescreen; the rhyme "here comes a chopper to chop off your head" recurs; Winston repeatedly says "we are the dead." The arrest smashes the paperweight — the symbol of their world — into fragments.

\textbf{Conclusion.} Orwell traces love from instinctive revolt to conscious conspiracy, only to show that the Party's power reaches precisely into the last refuge of the self. The theme thus carries the novel's bleakest warning: a regime that controls truth can also murder love.
\end{exoresolu}

\begin{exemplebox}[Theme Bank for the Three Set Texts]
\begin{itemize}
\item \emph{Nineteen Eighty-Four}: the destruction of the self (Part 3); the mutability of truth and the past ("Who controls the past controls the future"); language as an instrument of power (Newspeak); love and loyalty as rebellion (Part 2).
\item \emph{The Canterbury Tales}: the corruption of the Church versus true Christianity; the dignity of honest labour; women and mastery in marriage (Wife of Bath, Merchant's Tale); flattery versus truth (Placebo and Justinus).
\item \emph{A Raisin in the Sun}: the American Dream and its cost; money versus dignity; cultural identity and African heritage; racism and housing segregation (Clybourne Park); the strength of family and of women (Mama, Ruth, Beneatha).
\end{itemize}
\end{exemplebox}

\courssub{Skill 5: Analysing a Dramatic Scene (Dialogue, Stage Directions, Symbolism)}

Drama is the only one of the three genres on your syllabus that is written to be \emph{performed}. A candidate who analyses \emph{A Raisin in the Sun} exactly like a novel loses marks: the examiner expects attention to \cle{stagecraft}.

\begin{methode}[How to Analyse a Scene from a Play]
\begin{enumerate}
\item \textbf{Remember a play is performance:} always consider stage directions, lighting, music, props and movement, not only words.
\item \textbf{Identify the conflict} of the scene: who wants what, and who or what opposes them?
\item \textbf{Track the power shifts:} note where the upper hand changes (e.g. Mama giving Walter the money transfers authority — and responsibility — to him).
\item \textbf{Analyse key symbols:} the insurance cheque, Mama's plant, Beneatha's Nigerian dress and hair, Lindner's briefcase.
\item \textbf{Connect to context:} 1950s Chicago, segregation, the Great Migration heritage, African independence movements (Asagai's references).
\item \textbf{Conclude} on the scene's contribution to the play's central question: what should a family do with its dreams in a racist society?
\end{enumerate}
\textbf{Reflex:} for every quotation of dialogue, ask "how would this be \emph{played} on stage?"
\end{methode}

\begin{exoresolu}[Worked Example: Beneatha, Asagai and Cultural Identity (Act 2 Scene 1)]
\textbf{Statement:} Analyse how Hansberry uses Act 2 Scene 1 — the Nigerian dress and the "assimilationist" debate — to explore the theme of cultural identity.
\tcblower
\textbf{Conflict of the scene.} The scene stages a friendly but pointed clash between two models of Black identity: Asagai, the Nigerian intellectual who urges pride in African roots, and George Murchison, the wealthy Black American who dismisses African heritage and represents assimilation into white middle-class culture. Beneatha stands between them — literally and symbolically.

\textbf{Stagecraft and symbolism.} Hansberry's stage directions do the thematic work. Beneatha enters wearing the Nigerian robes Asagai has given her, her hair cut "natural" — a radical choice for 1959, when straightened hair was the norm. The visual shock on stage embodies her search for an identity that is not defined by white standards. The music matters too: the Nigerian folk song she plays, and the dance Walter drunkenly joins, transform the cramped living-room for a moment into an imagined Africa — the stage itself seems to dream. When George arrives and mocks her hair, the collision of the two worlds is played out on her own body.

\textbf{Dialogue and ideas.} Asagai's teasing name for Beneatha — "Alaiyo," "One for Whom Bread Is Not Enough" — names her spiritual hunger, which money cannot satisfy. George, by contrast, calls her interest in Africa sentimental nonsense and tells her to stop being "eccentric." Hansberry gives each man a coherent position, but the scene's sympathy is clear: George's sophistication is shown as borrowed, while Asagai connects identity to history and to the anti-colonial struggle then unfolding in Africa.

\textbf{Context.} Written at the dawn of African independence and the American civil-rights movement, the scene asks a question that was urgent for Hansberry's audience: can Black Americans be free while despising their own origins?

\textbf{Conclusion.} Through costume, music and debate rather than sermon, Hansberry makes cultural identity a living stage conflict. The scene establishes that the Youngers' struggle is not only economic (the cheque) but existential: the question of \emph{who they are} must be answered before the question of \emph{where they will live}.
\end{exoresolu}

\begin{exemplebox}[White Flight and Racism: Act 2 Scene 2 (opening) and Scene 3]
In Act 2 Scene 2, Mrs Johnson, the neighbour, brings a newspaper report of a bomb set off in a Black family's new house in a white area — a piece of \cle{foreshadowing} that names the danger awaiting the Youngers. Her visit also dramatises "white flight": whites abandon a neighbourhood as soon as Black families move in, so that integration never stabilises. In Act 2 Scene 3, Mr Lindner of the Clybourne Park "Improvement Association" embodies \cle{polite racism}: he never insults the family; he speaks of "the way we do things" and offers money to buy the house back. Hansberry's point is sharp: segregation works not only through violence but through reasonable-sounding men with briefcases. The family's stunned silence, then Walter's controlled fury, show that courtesy does not soften the insult — it sharpens it.
\end{exemplebox}

\courssub{Skill 6: Analysing Structure and Debate in \emph{The Merchant's Tale}}

\emph{The Merchant's Tale} belongs to the genre of the \cle{fabliau}: a short comic verse tale of trickery, usually involving an old husband, a young wife and a clever lover. But Chaucer complicates the genre by framing the tale with the Merchant's bitter Prologue and by inserting a formal \cle{debate} on marriage before the action begins.

\begin{methode}[How to Analyse a Tale's Prologue, Structure and Debate]
\begin{enumerate}
\item \textbf{Read the Prologue as a key:} the pilgrim's own words before the tale reveal his bias. The Merchant, newly and bitterly married, announces a tale born of personal disillusion — so his tale is not objective.
\item \textbf{Map the structure:} identify the fabliau elements (old husband, young wife, clever lover, trick) and the debate passages that interrupt the narrative.
\item \textbf{Analyse the debate as drama:} in the marriage debate, Placebo (the flatterer, whose name means "I shall please") tells January what he wants to hear; Justinus (the "just one") warns against marrying a young wife. Note whose advice January follows and why.
\item \textbf{Track the irony:} January's "delusion" — his belief that marriage to young May will be paradise on earth — is systematically undercut by the narrator's sarcasm and by events (the garden, the pear tree, the blindness).
\item \textbf{Conclude on meaning:} is the tale a cynical attack on marriage, on old men's folly, or on self-deception generally? Support your verdict with the Prologue.
\end{enumerate}
\textbf{Reflex:} in Chaucer, always separate \emph{what the tale shows} from \emph{what the teller intends}.
\end{methode}

\begin{exoresolu}[Worked Example: January's Delusion and the Marriage Debate]
\textbf{Statement:} "In \emph{The Merchant's Tale}, the debate between Placebo and Justinus exposes January's folly before the tale's action even begins." Discuss.
\tcblower
\textbf{The Prologue as key.} The Merchant opens by lamenting his own two-month-old marriage: his wife is, he says, a torment compared to whom even the shrewish wives of other tales would seem mild. This confession frames the tale as an act of personal bitterness — the reader is warned to expect a jaundiced view of marriage.

\textbf{Structure of the debate.} January, a sixty-year-old knight of Lombardy, decides to marry a girl under twenty and calls his friends to counsel him — but, significantly, he asks for counsel only \emph{after} his decision is made. The debate is therefore a formality, and Chaucer exploits this irony. Placebo, the court flatterer, argues that January is wise and needs no advice — his very name ("I shall please") announces his function: he tells the old man that his own judgement is best. Justinus, by contrast, speaks frankly: he warns that a young wife may bring an old husband misery and cuckoldry, and cites the difficulty of pleasing a young woman. The debate thus dramatises the eternal choice between \cle{flattery} and \cle{truth}.

\textbf{January's delusion.} January's reasoning reveals that he does not want advice but confirmation. He imagines marriage as a private paradise — he speaks of wedlock as an earthly Eden — while simultaneously treating his future wife as property, a possession to be chosen like livestock for her youth and beauty. Chaucer's irony is surgical: the very qualities January selects (May's youth, her freshness) are what will make her desire the young squire Damian. His later physical blindness is the literal fulfilment of the moral blindness he displays in the debate: he refuses to \emph{see} what Justinus tells him.

\textbf{Function within the tale.} The debate is thus not a digression but the tale's thesis stated in advance: the action (the garden, the pear-tree trick, the quarrel of the gods Pluto and Proserpina) merely proves what Justinus predicted. Chaucer invites us to watch a man construct his own trap.

\textbf{Conclusion.} The statement is justified. By placing the debate before the marriage, Chaucer makes the tale a demonstration of self-deception: January hears the truth, chooses flattery, and is duly punished — though the Merchant's bitter narration leaves us unsure whether the target is foolish old husbands or wives themselves.
\end{exoresolu}

\courssub{Skill 7: Planning and Writing the GCE Literature Essay}

All the skills above converge in the examination essay. The GCE Advanced Level rewards a clear thesis, precise textual evidence, commentary on technique, and a weighed conclusion.

\begin{methode}[The Five-Paragraph GCE Essay Plan]
\begin{enumerate}
\item \textbf{Decode the question (3 minutes):} underline the key terms ("discuss," "trace," "to what extent") and define the exact task. Never write on the work in general.
\item \textbf{Brainstorm and select (5 minutes):} list all relevant episodes/quotations, then choose the three or four strongest and order them logically (usually chronologically through the text).
\item \textbf{Introduction:} name author, work and genre; define the key term of the question; announce your line of argument (thesis).
\item \textbf{Body paragraphs:} one idea each, following the pattern claim -- evidence -- technique -- thematic link.
\item \textbf{Conclusion:} answer the question directly, weigh the evidence ("to what extent…"), and end with the author's wider purpose.
\end{enumerate}
\textbf{Reflex:} examiners reward \emph{argument + evidence + technique}. Retelling the plot without analysis is the most common cause of failure.
\end{methode}

\begin{exoresolu}[Worked Example: Full Essay Plan with Model Paragraph]
\textbf{Statement:} "The Youngers' decision to move to Clybourne Park is a victory, but a fragile one." To what extent do you agree? (Refer to \emph{A Raisin in the Sun}, especially Act 2 Scene 3 and Act 3.)
\tcblower
\textbf{Introduction (model).} In \emph{A Raisin in the Sun} (1959), Lorraine Hansberry dramatises a Black family's struggle to turn an insurance cheque into a dignified life. The purchase of a house in the white neighbourhood of Clybourne Park, and the family's final refusal of Mr Lindner's buy-out offer, constitute the play's resolution. This essay argues that the decision is a genuine moral victory — the family chooses dignity over money — but a deliberately fragile one, since Hansberry leaves the economic and racial dangers unresolved.

\textbf{Plan of the body.}
\begin{itemize}
\item \textbf{Paragraph 1 — The obstacle:} Lindner's visit (Act 2 Scene 3). The "welcoming committee" offers to buy the house back; his polite vocabulary ("you people," "the way we do things") shows racism wearing a courteous mask. Evidence: his euphemisms and the family's stunned silence.
\item \textbf{Paragraph 2 — The collapse that makes surrender tempting:} the loss of the money (Act 2 Scene 2) means the family now \emph{needs} Lindner's cash. Walter's plan to accept the offer, performing humiliation "like a slave," is the play's darkest moment.
\item \textbf{Paragraph 3 — The victory:} before Lindner and his own son Travis, Walter refuses: "we are very proud people" and they will move into the house. Technique: Hansberry makes him speak as head of the family for the first time; Mama confirms he has "come into his manhood."
\item \textbf{Paragraph 4 — The fragility:} the money is still gone; Beneatha's studies are uncertain; Clybourne Park's hostility awaits (Hansberry's audience knew of real bombings of Black homes in white districts, already evoked by Mrs Johnson's newspaper). The final image — Mama returning for her struggling plant — symbolises hope that is alive but delicate.
\end{itemize}

\textbf{Model body paragraph (Paragraph 3).} Walter's refusal of Lindner's offer is the play's second and decisive climax. Hansberry stages it as a public test: Travis is deliberately kept in the room, so that Walter must perform manhood before his son. His speech to Lindner reverses the language of the earlier humiliation he had rehearsed: instead of the grinning clown, he stands "like a man" and declares that the family will occupy the house his father earned. The technique is one of \cle{echo and reversal}: the same money that destroyed him now becomes irrelevant, proving that his identity no longer depends on it. Thematically, the moment answers the play's central question — the dream that matters is not the liquor store but self-respect transmitted across generations.

\textbf{Conclusion (model).} The decision is therefore a victory of values over circumstances; but Hansberry, writing before the civil-rights victories of the 1960s, honestly refuses a fairy-tale ending. The Youngers move out with nothing secured except their unity and pride — which, the play suggests, is precisely the foundation on which everything else must be built.
\end{exoresolu}

\begin{attention}[Common Mistakes to Avoid in the GCE Literature Essay]
\begin{itemize}
\item \textbf{Retelling the plot.} Summary earns almost no marks; every sentence of narration must serve an argument.
\item \textbf{Invented quotations.} Quote only short phrases you are certain of ("I love you"; "we are the dead"); otherwise paraphrase precisely with the chapter or scene reference.
\item \textbf{Ignoring technique.} In Literature, \emph{how} the author writes (irony, symbolism, stage directions, foreshadowing) is as important as \emph{what} happens.
\item \textbf{Confusing the Chaucer pilgrims.} Do not mix up the Summoner (court officer) with the Pardoner (seller of indulgences), or the Physician with the Parson.
\item \textbf{Treating the play as a novel.} For \emph{A Raisin in the Sun}, always mention at least one element of stagecraft (props, lighting, costume, movement).
\item \textbf{One-sided answers.} Questions with "to what extent" or "discuss" require you to weigh both sides before concluding.
\end{itemize}
\end{attention}

\begin{exercice}[Practice Questions]
\begin{enumerate}
\item "Part 2 of \emph{Nineteen Eighty-Four} is the story of a doomed love." Discuss with close reference to Chapters 7--9.
\item "In the \emph{General Prologue}, Chaucer praises skill and punishes hypocrisy." Test this view against the portraits of the Shipman, the Physician, the Parson and the Pardoner.
\item Analyse the dramatic function of Mr Lindner's visit in Act 2 Scene 3 of \emph{A Raisin in the Sun}.
\item "January's tragedy is that he hears the truth and chooses flattery." Discuss the role of Placebo and Justinus in \emph{The Merchant's Tale}.
\item "In Part 3 of \emph{Nineteen Eighty-Four}, Orwell shows that the self can be completely destroyed." To what extent do you agree?
\end{enumerate}
\end{exercice}

\begin{corrige}[Exercise 1]
Strong answers define love as rebellion (the Party abolishes private loyalty), trace the affair from Julia's note through the room above Charrington's shop to the visit to O'Brien, and show the preparation of doom: the loaded exception in the lovers' oath, the recurring rhyme of the chopper, "we are the dead," the hidden telescreen. The best candidates add Chapter 9's irony — Goldstein's book, the rebels' supposed truth, is itself a Party product — and conclude that Orwell makes the love story carry his warning about totalitarian power over the inner life.
\end{corrige}

\begin{corrige}[Exercise 2]
The statement is largely confirmed but needs nuance. Skill is genuinely praised even in flawed figures: the Shipman is the best seaman "from Hull to Cartagena," the Physician is "a verray, parfit praktisour," the Cook a master of sauces. Hypocrisy is punished by irony: the Physician's gold, the Pardoner's relics, the Summoner's quart of wine. But the Parson and Plowman show that Chaucer also praises \emph{virtue} directly, without irony — so the formula should be completed: Chaucer praises skill, reveres goodness, and reserves satire for those who betray their office.
\end{corrige}

\begin{corrige}[Exercise 3]
Lindner's visit is the externalisation of the play's central conflict: racism arrives not as a mob but as a courteous offer of money. Dramatically, it plants the temptation that the loss of the money (Scene 2) will make irresistible, thus linking the two climaxes. Hansberry's technique: euphemistic diction ("you people"), the prop of the briefcase, and the family's frozen silence. The scene also unites the family for the first time in a shared refusal — a rehearsal for Walter's final stand in Act 3.
\end{corrige}

\begin{corrige}[Exercise 4]
Candidates should establish that January requests counsel only after deciding, making the debate a formality; contrast Placebo (name: "I shall please"; advice: trust yourself) with Justinus (name: the just one; warning: a young wife may cuckold an old husband); and show that the tale's action — May, Damian, the garden, the pear tree — fulfils Justinus's prediction. The conclusion should weigh the Merchant's bitter Prologue: the tale demonstrates self-deception, but its teller's bias leaves the final target ambiguous.
\end{corrige}

\begin{corrige}[Exercise 5]
Strong answers follow the three stages of reintegration (learning, understanding, acceptance), analyse O'Brien's dual role as torturer and teacher, and show what is destroyed at each stage: bodily dignity, intellectual certainty (two and two make five), and finally love (the betrayal of Julia in Room 101). The closing state — Winston loving Big Brother — confirms total destruction. Nuanced answers may note the small residue of memory earlier in Part 3, but the ending leaves no space for inner resistance: the victory of the Party is complete, which is Orwell's final warning.
\end{corrige}

\newpage

\coursec{Exercises}

\courssub{I check my knowledge}
\begin{exercice}[MCQ: Nineteen Eighty-Four Part 2]
Which of the following best describes the significance of Winston's memory of his mother and sister in Chapter 7?
\begin{enumerate}
\item It reveals the Party's complete control over family life.
\item It shows that Winston's rebellion is rooted in a longing for a lost humanity.
\item It proves that the Party can erase even the deepest personal memories.
\item It foreshadows Winston's eventual betrayal of Julia.
\end{enumerate}
\end{exercice}

\begin{exercice}[True/False: The General Prologue Portraits]
State whether each statement is true or false. Justify your answer with a brief reference to the text.
\begin{enumerate}
\item The Cook is praised solely for his culinary skills without any moral judgement.
\item The Shipman (Skipper) is portrayed as an honest and law-abiding merchant.
\item The Doctor's knowledge of astronomy is presented as a genuine scientific pursuit.
\item The Parson and the Plowman are the only pilgrims who fully embody Christian ideals.
\end{enumerate}
\end{exercice}

\begin{exercice}[Definitions: Key Terms in A Raisin in the Sun Act 2]
Define the following terms as they relate to the play's themes and context:
\begin{enumerate}
\item \cle{White flight}
\item \cle{Assimilationism} (as represented by George Murchison)
\item \cle{The American Dream} (as understood by Walter Lee)
\item \cle{Climax} (identify the climax of Act 2)
\end{enumerate}
\end{exercice}

\begin{exercice}[Direct Application: Identifying Plot Points]
Match each event to the correct chapter or scene:
\begin{enumerate}
\item Winston reads Goldstein's book in the room above Mr Charrington's shop.
\item The Merchant tells his prologue, revealing his own unhappy marriage.
\item Walter receives the insurance money and gives it to Willy Harris.
\item The Pardoner admits his relics are fake but claims he can still preach against greed.
\item Mr Lindner visits the Younger family to offer a buyout.
\item Winston is tortured in the Ministry of Love and learns the Party seeks power for its own sake.
\end{enumerate}
\end{exercice}

\courssub{I practise}
\begin{exercice}[Passage Analysis: Winston and Julia's Rebellion]
Read the following paraphrase of a key moment in Part 2, Chapter 7: \emph{Winston realises that the proles have stayed human because they have not been hardened by the Party's ideology; he and Julia, by contrast, are "the dead" who only think they are alive.} Write a short analysis (150–200 words) explaining how this moment develops the theme of hope versus futility. Focus on Orwell's use of contrast and the symbolic weight of the paperweight.
\end{exercice}

\begin{exercice}[Character Study: The Doctor and the Wife of Bath]
Compare Chaucer's presentation of the Doctor (Physician) and the Wife of Bath in the General Prologue. In about 200 words, discuss how each portrait blends professional skill with moral ambiguity. Use at least two specific details from each portrait.
\end{exercice}

\begin{exercice}[Theme Exploration: Money and the American Dream]
In Act 2 Scene 1, the Younger family debates how to use the insurance money. Write a structured response (two paragraphs) showing how Hansberry uses this debate to expose:
\begin{itemize}
\item The different definitions of the \cle{American Dream} held by Mama, Walter, and Beneatha.
\item The tension between cultural identity and material success.
\end{itemize}
\end{exercice}

\begin{exercice}[Structural Analysis: The Merchant's Tale Prologue]
Explain how the Merchant's Prologue functions as a frame for the tale that follows. Your answer should cover:
\begin{enumerate}
\item The Merchant's own marital experience and its influence on his narrative stance.
\item The debate between Placebo and Justinus as a structural device.
\item The irony of the Merchant telling a tale about a deluded old husband.
\end{enumerate}
\end{exercice}

\courssub{I prepare for the GCE}
\begin{exercice}[Essay: Hope and Destruction in Part 2 of Nineteen Eighty-Four (25 marks)]
\textbf{Question:} "In Part 2 of \emph{Nineteen Eighty-Four}, Orwell builds hope only to destroy it." Discuss this view with close reference to the plot of Chapters 7–9. Your essay should analyse the development of Winston and Julia's rebellion, the role of O'Brien, and the significance of Goldstein's book.
\end{exercice}

\begin{exercice}[Essay: Ideal versus Corrupt in the General Prologue (25 marks)]
\textbf{Question:} "Chaucer's ideal characters (the Parson and the Plowman) exist primarily to expose the corrupt ones." How far do you agree? Discuss with reference to at least four portraits from the General Prologue, including the Summoner, the Pardoner, the Doctor, and the Wife of Bath.
\end{exercice}

\begin{exercice}[Comparative Essay: Plot as Social Criticism (30 marks)]
\textbf{Question:} "Writers use plot to criticise the societies they describe." Compare how \emph{any two} of your set texts (\emph{Nineteen Eighty-Four}, \emph{The Canterbury Tales: The General Prologue and The Merchant's Tale}, \emph{A Raisin in the Sun}) achieve this. In your answer, consider narrative structure, key turning points, and the resolution of each plot.
\end{exercice}

\newpage

\coursec{Solutions to the exercises}

\begin{corrige}[Exercise 1 — MCQ: Nineteen Eighty-Four Part 2]
\textbf{Correct answer: (b)} — It shows that Winston's rebellion is rooted in a longing for a lost humanity.

\textbf{Justification.} In Part 2 Chapter 7, Winston remembers a childhood scene of hunger and selfishness during which he stole his little sister's chocolate ration, after which his mother and sister disappeared (presumably ``vaporised''). The memory returns to him as a dream in which his mother's death seems to belong to a time when feelings such as love, loyalty and private grief still had meaning. Orwell makes Winston reflect that his mother's death was tragic in a way that is no longer possible, because tragedy depends on \cle{private loyalties} that the Party has abolished. This is the psychological foundation of his rebellion: he does not revolt for political doctrine but because he senses that something human --- love, memory, family devotion --- has been stolen from him and his generation.

\textbf{Why the other options fail:}
\begin{itemize}
\item (a) is partially true of the Party's aims (Junior Spies, the destruction of the family), but the memory's \emph{significance} in the chapter is emotional and moral, not merely institutional.
\item (c) is contradicted by the text: the memory survives precisely because the Party has \emph{not} erased it; the chapter stresses the persistence of memory.
\item (d) is tempting but wrong: the memory does not foreshadow the betrayal of Julia; the betrayal is prepared by other elements (the fear of rats, the declaration ``we are the dead'').
\end{itemize}
\end{corrige}

\begin{corrige}[Exercise 2 — True/False: The General Prologue Portraits]
\begin{enumerate}
\item \textbf{False.} The Cook is an excellent craftsman --- Chaucer praises his skill with dishes and his famous ``blankmanger'' --- but the portrait contains a sharp moral and physical judgement: he has an ulcerous sore (``mormal'') on his shin, and the Host later jokes about his dubious hygiene. Skill and moral discomfort coexist.
\item \textbf{False.} The Shipman is a superb sailor but a lawless man: he steals wine from the cargo he transports (from Bordeaux, ``while the chapman sleeps''), has no scruples of conscience, and sends his prisoners home ``by water'' --- that is, he makes them walk the plank. He is competent, not honest.
\item \textbf{False.} The Doctor's astronomy is \emph{astrology} used for profit: he chooses ``fortunate'' hours for treatment according to the stars (``magik natureel''). Chaucer adds the devastating line that his study of the Bible was small, and that he loved gold above all --- ``for gold in phisik is a cordial.'' His science serves avarice.
\item \textbf{True (with nuance).} The Parson is a poor but holy priest who practises what he preaches (``first he wroghte, and afterward he taughte''), and the Plowman, his brother, is a true labourer who lives in charity, loves God and his neighbour, and pays his tithes faithfully. Among the pilgrims described in the syllabus, they alone fully embody the Christian ideal; Chaucer's praise of them is unironic, unlike nearly every other portrait.
\end{enumerate}
\end{corrige}

\begin{corrige}[Exercise 3 — Definitions: Key Terms in A Raisin in the Sun Act 2]
\begin{enumerate}
\item \cle{White flight}: the movement of white families out of a neighbourhood when Black families move in, driven by racial prejudice and fear of falling property values. In the play it is the background to the Clybourne Park conflict: the white residents' association sends Mr Lindner to buy the Youngers out rather than live beside them.
\item \cle{Assimilationism}: the attitude of abandoning one's African or Black cultural heritage to adopt white, middle-class American values. George Murchison, Beneatha's wealthy suitor, embodies it: he mocks her natural hair and her interest in African dress and history, telling her to stop being ``eccentric'' and simply conform. Hansberry opposes him to Joseph Asagai, who encourages Beneatha's African identity.
\item \cle{The American Dream} (Walter Lee's version): for Walter, the dream is reduced to \cle{material success} --- money, a business of his own (the liquor store), a car, pearls for his wife, the power to give his son ``the world.'' Hansberry shows this as a distorted, money-centred version of the dream, contrasted with Mama's dream of a house and dignity and Beneatha's dream of self-realisation through medicine.
\item \cle{Climax}: the decisive turning point of highest tension after which the outcome becomes inevitable. In Act 2, the climax occurs in Scene 2 when Walter learns that Willy Harris has \textbf{absconded with all the money} --- both the liquor-store investment and Beneatha's school money. This catastrophe (``Climax 1: Loss'') destroys Walter's dream and forces the family's final moral test in Act 3.
\end{enumerate}
\end{corrige}

\begin{corrige}[Exercise 4 — Direct Application: Identifying Plot Points]
\begin{enumerate}
\item Winston reads Goldstein's book in the room above Mr Charrington's shop. $\rightarrow$ \textbf{Nineteen Eighty-Four, Part 2, Chapter 9.}
\item The Merchant tells his prologue, revealing his own unhappy marriage. $\rightarrow$ \textbf{The Canterbury Tales --- The Merchant's Prologue} (he laments two months of marriage to a ``shrewe,'' contrasting his sorrow with the Clerk's patient Griselda).
\item Walter receives the insurance money and gives it to Willy Harris. $\rightarrow$ \textbf{A Raisin in the Sun, Act 2, Scene 2} (Mama entrusts him with the remaining \$6{,}500; he hands it over for the liquor-store scheme).
\item The Pardoner admits his relics are fake but claims he can still preach against greed. $\rightarrow$ \textbf{The General Prologue} (portrait of the Pardoner: pigs' bones sold as saints' relics, yet ``in chirche... a noble ecclesiaste'').
\item Mr Lindner visits the Younger family to offer a buyout. $\rightarrow$ \textbf{A Raisin in the Sun, Act 2, Scene 3.}
\item Winston is tortured in the Ministry of Love and learns the Party seeks power for its own sake. $\rightarrow$ \textbf{Nineteen Eighty-Four, Part 3, Chapters 2--3} (O'Brien's interrogation: ``The Party seeks power entirely for its own sake'').
\end{enumerate}
\end{corrige}

\begin{corrige}[Exercise 5 — Passage Analysis: Winston and Julia's Rebellion]
\textbf{Model answer (approximately 180 words).}

This moment crystallises Part 2's central tension between hope and futility through a carefully built contrast. Winston's famous reflection that ``if there is hope, it lies in the proles'' rests on the idea that the proles, ignored by the Party, have preserved instinctive human loyalties --- love, family, memory --- which Party members have lost. Against them Orwell sets Winston and Julia, who call themselves ``the dead'': intellectually awakened but already doomed, alive only in anticipation of their inevitable capture. The contrast is double-edged. The proles are human but unconscious, unable to rebel because they cannot understand their condition; Winston and Julia understand everything but are powerless. Hope and futility are thus distributed between two groups, and neither possesses both. The glass paperweight concentrates this symbolism: the tiny coral embedded in crystal is a miniature of the private, beautiful, useless world the lovers try to inhabit --- a fragment of the past preserved inside transparent fragility. Its very transparency anticipates the moment the Thought Police will smash it: hope in Part 2 is real but made of glass, and Orwell lets the reader cherish it precisely so that its destruction will hurt.

\textbf{Examiner's expectations:} identification of the contrast (proles/Party members), the paradox of ``the dead,'' the paperweight as symbol of fragile private life, and a conclusion linking hope to its foreshadowed destruction.
\end{corrige}

\begin{corrige}[Exercise 6 — Character Study: The Doctor and the Wife of Bath]
\textbf{Model answer (approximately 210 words).}

Both portraits show Chaucer's characteristic method: professional excellence entangled with moral ambiguity.

The Doctor is a perfect practitioner --- ``verray, a parfit praktisour'' --- grounded in astronomy (astrology) and ``magik natureel,'' able to diagnose by the humours and to choose propitious hours for treatment. Yet the portrait darkens steadily: he works in profitable collusion with his apothecaries, each sending business to the other; his diet is moderate but his appetite for money is not; he has read little of the Bible; and the climactic couplet reveals that he ``lovede gold in special,'' since gold is a cordial in medicine. Skill serves avarice.

The Wife of Bath is equally skilled in her trade: her cloth-making surpasses the weavers of Ypres and Ghent, and she is a veteran pilgrim who knows much of ``wandrynge by the weye.'' Her ambiguity is social and sexual rather than financial: she insists on taking precedence at the offering, wears flamboyant scarlet stockings and enormous hats, is ``gat-tothed'' (a traditional sign of sensuality), and has had five husbands, ``withouten oother compaignye in youthe.'' Chaucer's tone is amused rather than condemnatory.

Together the portraits show that in the General Prologue, competence is never the whole of a person: the Doctor cures bodies but is spiritually sick; the Wife is morally irregular but vividly, generously alive.

\textbf{Examiner's expectations:} two precise details per portrait, the formula ``skill + moral flaw,'' and a comparative conclusion on Chaucer's irony.
\end{corrige}

\begin{corrige}[Exercise 7 — Theme Exploration: Money and the American Dream]
\textbf{Model structured response.}

\emph{Paragraph 1 --- Three dreams, one cheque.} The \$10{,}000 insurance payment forces each Younger to define the American Dream aloud. For Mama, the dream is \cle{security and dignity}: a house with a yard and a garden, the dream she and Big Walter deferred all their lives --- money is a means to shelter the family, never an end. For Walter, the dream is \cle{entrepreneurial success}: the liquor store represents the power to stop being a chauffeur, to ``make transactions'' like the rich white men he serves, to buy pearls for Ruth and choices for Travis. For Beneatha, the dream is \cle{self-fulfilment and identity}: medical school, independence, and --- through Asagai --- a recovered African heritage. Hansberry's stagecraft makes the cheque a mirror: each character sees a different America in it.

\emph{Paragraph 2 --- Identity against materialism.} The debate exposes the cost of defining the dream in cash. Walter's contempt for George Murchison is double-edged: George has money but has purchased it with assimilation, mocking African dress and heritage, while Walter wants money without realising he risks the same self-loss. Beneatha, cutting her hair and wearing Asagai's Nigerian robes, asserts that identity cannot be bought. Mama's decision to buy the Clybourne Park house --- and her later entrusting of the rest to Walter --- dramatises the play's question: can money serve the family's dignity and roots, or will it master them? Act 2 Scene 2's catastrophe answers provisionally: money pursued as the dream itself destroys; money subordinated to love and identity may yet build.

\textbf{Examiner's expectations:} clear differentiation of the three dreams, reference to concrete stage moments (the house, the liquor-store plan, the robes and hair), and explicit linkage of money to the identity theme.
\end{corrige}

\begin{corrige}[Exercise 8 — Structural Analysis: The Merchant's Tale Prologue]
\begin{enumerate}
\item \textbf{The Merchant's marital experience.} In his Prologue the Merchant confesses he has been married only two months, yet his wife is a ``shrewe'' whose cruelty, he claims, exceeds anything the Clerk's patient Griselda endured. This personal bitterness determines his \cle{narrative stance}: he tells the tale of January not as a neutral storyteller but as a wounded husband eager to prove that marriage is a snare and wives are deceivers. The Prologue thus functions as a frame of bias: the reader is warned to read the tale as an act of revenge, not objective truth.

\item \textbf{Placebo and Justinus as structural device.} Before the marriage, January consults his two counsellors. Placebo (whose name means ``I shall please'') flatters him and approves everything; Justinus honestly warns him against marrying a young wife in old age. The debate dramatises the choice between flattery and truth --- and January, predictably, follows Placebo. Structurally, the debate delays the narrative, exposes January's self-delusion before the action begins, and establishes the tale's central irony: the delusion is chosen, not imposed. It also mirrors the Prologue: Justinus speaks what the Merchant, in his bitterness, now believes.

\item \textbf{The irony of the teller.} The Merchant, freshly embittered by his own marriage, tells of an old knight deluded into marrying young May, who cuckolds him with Damyan. The irony is layered: the Merchant mocks January's blindness while displaying his own --- he generalises from two months of marriage to all women; he condemns January's folly yet has himself just rushed into marriage; and his tale's cynicism about wives rebounds on him, since his own account of his marriage may be as one-sided as January's self-justifications. The frame turns the tale into an unintentional self-portrait.
\end{enumerate}
\end{corrige}

\begin{corrige}[Exercise 9 — Essay: Hope and Destruction in Part 2 of Nineteen Eighty-Four (25 marks)]
\textbf{Indicative mark scheme (25 marks):}
\begin{itemize}
\item Knowledge of plot, Chapters 7--9 (5 marks): accurate sequence of events and textual reference.
\item Analysis of the construction of hope (7 marks): the affair, the room, the proles, the paperweight, O'Brien, Goldstein's book.
\item Analysis of the destruction of hope (7 marks): ``we are the dead,'' the irony of the Hate Week arrest context, the falsity of the sanctuary.
\item Evaluation of the statement and personal, argued response (4 marks).
\item Expression, structure, and quotation/paraphrase accuracy (2 marks).
\end{itemize}

\textbf{Model plan and key content.}

\emph{Introduction.} Agree largely with the statement: Part 2 is engineered as a rising curve of hope --- love, privacy, conspiracy, knowledge --- each element of which contains the seed of its own destruction. Orwell's method is dramatic irony: the reader is allowed to hope in order to feel the loss.

\emph{Paragraph 1 --- Love as constructed hope (Ch. 7 context).} Winston's memories of his mother and the prole woman hanging washing establish what hope is \emph{for}: a lost humanity. The affair with Julia, the room above Charrington's shop, the paperweight: Orwell builds a private counter-world. Yet each sanctuary is compromised from the start --- the room is rented from a Thought Police agent, the lovers themselves say ``we are the dead.''

\emph{Paragraph 2 --- O'Brien: hope personified as trap.} Winston's long-held belief that O'Brien is a secret dissident is the novel's cruellest hope. The initiation scene, the wine, the Brotherhood's ritual questions (``You are prepared to... throw sulphuric acid in a child's face?'') should warn the reader: the rebellion demands inhumanity, mirroring the Party it opposes. O'Brien gives hope an address --- and the address is false.

\emph{Paragraph 3 --- Goldstein's book: the hope of understanding (Ch. 9).} The book explains the mechanics of oligarchical collectivism, war as economic control, doublethink. For Winston, to understand is to hope: ``the book fascinated him, or more exactly it reassured him.'' But the book never answers \emph{why} --- and it is itself a Party fabrication, written by a collective including O'Brien. Even knowledge is manufactured.

\emph{Paragraph 4 --- Destruction.} The arrest at the chapter's end: the smashed paperweight, the voice behind the picture, ``the house is surrounded.'' Every hope of Part 2 is revealed as curated by the Party.

\emph{Conclusion.} Orwell builds hope with real tenderness --- the coral, the thrush, Julia's arm --- precisely because the political point requires it: a totalitarianism that merely crushed would be less terrifying than one that first teaches its victims to hope.

\textbf{Examiner's expectations:} candidates must go beyond plot summary; top-band answers track the \emph{pattern} (hope built then annulled) across all three elements and comment on Orwell's structural irony.
\end{corrige}

\begin{corrige}[Exercise 10 — Essay: Ideal versus Corrupt in the General Prologue (25 marks)]
\textbf{Indicative mark scheme (25 marks):}
\begin{itemize}
\item Knowledge of at least four required portraits (6 marks).
\item Analysis of the ideal portraits and their function (6 marks).
\item Analysis of the corrupt portraits and Chaucer's satiric technique (6 marks).
\item Argued judgement on ``how far'' the claim holds (5 marks).
\item Expression and structure (2 marks).
\end{itemize}

\textbf{Model plan and key content.}

\emph{Introduction.} The claim is half true: the Parson and Plowman do provide the moral yardstick by which the Summoner, Pardoner and Doctor are measured --- but Chaucer's art is richer than a simple exposure device, since the corrupt portraits are also celebrated for their vitality.

\emph{Paragraph 1 --- The ideals.} The Parson: poor in goods, rich in holy thought; patient in adversity; visits his flock on foot, in rain, far apart; ``first he wroghte, and afterward he taughte.'' The Plowman: honest labourer, lives in peace and charity, loves God and neighbour, pays tithes. Crucially, Chaucer's irony --- everywhere else his tool --- is absent here: the unironic praise marks them as the standard.

\emph{Paragraph 2 --- The corrupt measured against the standard.} The Summoner: lecherous, drunk on wine, bribable with a quart of ale, abuses his ecclesiastical court for blackmail; the Pardoner: sells fake relics (pigs' bones, a pillowcase as Mary's veil), preaches against avarice while practising it, and boasts of his own hypocrisy; the Doctor: colludes with apothecaries, loves gold, reads little Bible. Each fails precisely where the Parson succeeds: the Parson is ``noght... mercenarie'' while these three sell the Church's office.

\emph{Paragraph 3 --- The Wife of Bath complicates the scheme.} She is morally irregular (five husbands, precedence at the offering) yet Chaucer's tone is warm, comic, admiring of her energy and skill. If the ideals existed \emph{primarily} to expose corruption, the Wife would be condemned --- she is not. This weakens the claim.

\emph{Paragraph 4 --- Judgement.} The ideals function as a \emph{foil} and a \emph{standard}, yes; but the corrupt portraits have independent life --- they are studies of human energy, professional pride, and self-deception. Chaucer's satire is corrective, not merely denunciatory: the Parson shows what the Church could be, which is why the satire has a target at all.

\emph{Conclusion.} Agree ``to a large extent'': structurally the ideals anchor the satire; but Chaucer's comic generosity means the corrupt pilgrims are exposed \emph{and} enjoyed.

\textbf{Examiner's expectations:} the phrase ``how far'' demands a qualified thesis; candidates who treat the question as pure agreement rarely exceed the middle band.
\end{corrige}

\begin{corrige}[Exercise 11 — Comparative Essay: Plot as Social Criticism (30 marks)]
\textbf{Indicative mark scheme (30 marks):}
\begin{itemize}
\item Knowledge of both chosen plots, with accurate turning points (8 marks).
\item Analysis of narrative structure as a critical instrument (8 marks).
\item Analysis of resolutions and their social meaning (6 marks).
\item Genuine comparison (links, contrasts, connectives) (5 marks).
\item Expression, organisation, register (3 marks).
\end{itemize}

\textbf{Model approach --- exemplified with \emph{Nineteen Eighty-Four} and \emph{A Raisin in the Sun} (other pairings acceptable).}

\emph{Introduction.} Both Orwell and Hansberry make plot itself an argument: the \emph{shape} of events --- what is attempted, what is destroyed, what is salvaged --- constitutes the criticism of Oceania's totalitarianism and of 1950s American racism respectively.

\emph{Paragraph 1 --- Structure as criticism.} Orwell's tripartite structure (misery --- hope --- annihilation) is a syllogism: Part 1 establishes the crime of the system, Part 2 tests the possibility of escape, Part 3 proves escape impossible. The plot's very architecture denies the reader consolation, which \emph{is} the critique of totalitarianism. Hansberry's three-act domestic structure works differently: a single inheritance cheque organises the plot so that every social force --- segregation (Lindner), economic exclusion (Walter's chauffeur job), assimilation (George) --- must pass through one family's living room. The compressed setting turns private decisions into public indictments.

\emph{Paragraph 2 --- Turning points.} In \emph{Nineteen Eighty-Four}, the arrest and Room 101 reverse every hope: the turning point proves that the Party controls not only actions but inner reality (``two and two make five''). In \emph{A Raisin in the Sun}, Willy Harris's theft (Act 2 climax) appears to be Walter's destruction --- yet Hansberry redirects the plot: the true turning point becomes Walter's refusal of Lindner's money in Act 3. Where Orwell's turning point closes possibility, Hansberry's opens it --- but only through dignity, not cash.

\emph{Paragraph 3 --- Resolutions.} Orwell ends with Winston loving Big Brother: the resolution is the total victory of the system over the self, a warning aimed at the reader's own society. Hansberry ends with the family moving into Clybourne Park, facing an uncertain, hostile future: the resolution is a partial, costly victory that criticises America precisely by showing how much courage an ordinary family needs merely to occupy a house. One plot says: this is what happens if power goes unchecked. The other says: this is what resistance costs.

\emph{Paragraph 4 --- Comparison.} Both writers refuse comfortable endings, but for opposite critical purposes: Orwell's closed plot indicts a system that permits no future; Hansberry's open plot indicts a society that makes an ordinary dream require heroism. Plot, in both, is not decoration but verdict.

\emph{Conclusion.} Structure, turning point and resolution are the three places where plot becomes social criticism; Orwell weaponises inevitability, Hansberry weaponises difficulty.

\textbf{Examiner's expectations:} the word ``compare'' must be honoured continuously (connectives such as ``whereas,'' ``by contrast,'' ``similarly''); two parallel mini-essays with a comparative conclusion cap the mark in the middle band. Candidates choosing \emph{The Canterbury Tales} should exploit the pilgrimage frame and the Merchant's Prologue-to-tale structure as their ``narrative structure'' evidence.
\end{corrige}

\newpage

\coursec{Summary sheet}

\begin{popfigure}
\begin{tikzpicture}[
  mindmap,
  concept color=popblue,
  level 1 concept/.append style={level distance=4.5cm,sibling angle=360/7},
  level 2 concept/.append style={level distance=3cm},
  every node/.style={font=\small, text width=2.8cm, align=center},
  grow cyclic
]
\node[concept, align=center]{Chapter LIT02\\Emphasis on Storyline,\\Plot \& Situation}
  child[concept color=poporange] {node[concept, align=center]{1984 Part 2\\Love as Rebellion}}
  child[concept color=popgreen] {node[concept, align=center]{1984 Part 3\\Destruction of the Self}}
  child[concept color=poppurple] {node[concept, align=center]{General Prologue I\\Skilled Workers \& Traders}}
  child[concept color=poppink] {node[concept, align=center]{General Prologue II\\Women, Ideal Christians,\\Corrupt Churchmen}}
  child[concept color=popgold] {node[concept, align=center]{The Merchant's Tale\\Structure, Prologue,\\Marriage Debate}}
  child[concept color=popblue] {node[concept, align=center]{Raisin Act 2 Sc1\\American Dream,\\Money, Cultural Identity}}
  child[concept color=popmuted] {node[concept, align=center]{Raisin Act 2 Sc2--3\\Loss, Racism,\\Clybourne Park}};
\end{tikzpicture}
\legende{Mind map of Chapter LIT02: seven thematic branches covering the three set texts.}
\end{popfigure}

\begin{bilan}
\textbf{Key Definitions}
\begin{itemize}
  \item \cle{Plot} -- the causal and logical arrangement of events in a narrative; distinct from \cle{story} (chronological sequence).
  \item \cle{Situation} -- the initial set of circumstances (time, place, social context) that frames the action.
  \item \cle{Storyline} -- the thread of events as experienced by the reader, often shaped by narrative perspective.
  \item \cle{Totalitarianism} -- a political system where the state seeks total control over public and private life (Orwell).
  \item \cle{American Dream} -- the belief that prosperity and success are attainable through hard work, central to Hansberry's drama.
  \item \cle{Estates Satire} -- Chaucer's technique of portraying representatives of medieval social estates to expose corruption.
  \item \cle{Marriage Debate} -- the medieval literary topos (Placebo vs Justinus) questioning the merits of marriage.
\end{itemize}

\textbf{Core Concepts \& Themes to Master}
\begin{itemize}
  \item \emph{Nineteen Eighty-Four Part 2 (Ch 7--9)}: Winston and Julia's affair as political rebellion; the role of memory, the paperweight, and the prole woman's song; O'Brien's entrapment.
  \item \emph{Part 3 (Ch 1--4)}: Ministry of Love, systematic torture, \cle{doublethink}, \cle{Room 101}, the final betrayal (Julia) and spiritual death.
  \item \emph{General Prologue I}: Cook, Manciple, Shipman -- skilled workers/traders; realism vs idealisation; irony in physical descriptions.
  \item \emph{General Prologue II}: Physician (greed), Wife of Bath (female agency, experience vs authority), Parson \& Plowman (ideal Christian poverty), Summoner \& Pardoner (corruption, sale of indulgences), Host (frame narrator).
  \item \emph{The Merchant's Tale}: January's delusion (senex amans), the garden as locus of deceit, the Pluto/Proserpina interlude, the debate on marriage (Placebo's flattery vs Justinus' caution).
  \item \emph{A Raisin in the Sun Act 2 Sc1}: Walter's investment dream, Beneatha's African identity (Asagai), Mama's plant as symbol of deferred dream.
  \item \emph{Act 2 Sc2--3}: Walter's loss of the insurance money (climax 1), Lindner's offer (white flight), the family's ultimate decision to move -- affirmation of dignity over money.
\end{itemize}

\textbf{Analytical Methods}
\begin{enumerate}
  \item \textbf{Plot Mapping}: Draw a linear diagram for each text (exposition $\rightarrow$ rising action $\rightarrow$ climax $\rightarrow$ falling action $\rightarrow$ resolution). Note parallel plots (e.g., Winston/Julia vs Party).
  \item \textbf{Character-Function Grid}: List characters, their social estate (Chaucer), or ideological role (Orwell, Hansberry), and their function in advancing theme.
  \item \textbf{Contextual Anchor}: For every major episode, write one sentence linking it to historical context (post-war Britain, 1950s US segregation, 14th-century Church).
  \item \textbf{Comparative Paragraphs}: Use a two-column table to contrast similar themes across texts (e.g., rebellion in 1984 vs Raisin; corruption in Pardoner vs Party).
  \item \textbf{Close Reading Checklist}: Identify \cle{imagery}, \cle{irony}, \cle{symbolism}, \cle{narrative voice}, and \cle{dramatic irony} in selected passages.
\end{enumerate}

\textbf{Common Mistakes (Traps)}
\begin{itemize}
  \item Confusing \emph{plot} with \emph{story} -- always explain causality, not just sequence.
  \item Quoting long passages from memory; the exam rewards \textbf{short, precise} quotations (max 8 words) or accurate paraphrase.
  \item Treating Chaucer's pilgrims as mere stereotypes; they are \emph{individualised} through detail (e.g., the Cook's ulcer, the Manciple's cunning).
  \item Overlooking the \emph{frame narrative} in Canterbury Tales (Host's role, storytelling contest).
  \item Ignoring stage directions in \emph{Raisin} (e.g., Mama's plant, the lighting shift at Lindner's entrance).
  \item Writing a plot summary instead of an analytical essay; every paragraph must answer the question.
\end{itemize}

\textbf{GCE Advanced Level Tips}
\begin{itemize}
  \item \textbf{Essay Structure}: Introduction (thesis + roadmap) $\rightarrow$ 3--4 body paragraphs (PEEL: Point, Evidence, Explanation, Link) $\rightarrow$ Conclusion (synthesis, not repetition).
  \item \textbf{Time Management}: 45 min per essay; spend 5 min planning, 35 min writing, 5 min proofreading.
  \item \textbf{Quotation Bank}: Memorise 2--3 key lines per set text (e.g., ``Freedom is the freedom to say that two plus two make four'', ``What happens to a dream deferred?'', ``The lyf so short, the craft so long to lerne'').
  \item \textbf{Context Marks}: Explicitly mention \textbf{Orwell's Spanish Civil War experience}, \textbf{Hansberry's family legal battle (Hansberry v. Lee)}, \textbf{Chaucer's diplomatic career and knowledge of estates satire}.
  \item \textbf{Comparative Questions}: When asked to compare, use a \emph{point-by-point} structure rather than text-by-text.
  \item \textbf{Language Analysis}: Comment on \textbf{Orwell's Newspeak}, \textbf{Chaucer's rhyming couplets and iambic pentameter}, \textbf{Hansberry's use of African American Vernacular English}.
\end{itemize}
\end{bilan}

\findecours

\end{document}
